\section{Base lattices II: stages at members and heights}\label{sec:pf:baselat2}

\providecommand{\bltwoL}[1]{\mathcal{L}_{\mathrm{#1}}}
\providecommand{\bltwoN}[1]{N_{\mathrm{#1}}}
\providecommand{\bltwop}{\mathfrak{p}}
\providecommand{\bltwoSt}{\mathrm{St}}
\providecommand{\bltwoWt}{W_{\mathrm{t}}}
\providecommand{\bltwoOO}{\mathcal{O}}
\providecommand{\bltwoS}[1]{\mathbf{S}_{\mathrm{#1}}}

\subsection{Introduction}\label{sec:pf:baselat2:intro}

This section proves the member-existence theorem for the \emph{origin parameter}
(\cref{pf:baselat2:main} below). It is one of the two inputs of the member-bundle theorem
\cref{thm:member-bundles}; the other one, and the passage from the origin parameter to the parameter
$p_A$ of \cref{thm:member-bundles}, are treated in the section of
\lref{GlobalTheory.T.large_unramified_exponent_bundles_exist_alt}.

\paragraph{What is proved.} Recall from \cref{sec:families:nonempty} the notions of twisting data,
of \emph{large unramified exponents} ($\sfname{EvLgG}$), of a parameter $p$ and of the two statements
$\sfname{MemberE}(p)$ and $\sfname{ExBL}(p)$. For $h_L\in\sfname{RecTExists}$ (\cref{def:rectexists}) and
$c\in\R$ the development defines three nested \emph{norm-choice extensions}
\[
  \bltwoN{law}\ \subset\ \bltwoN{list}(c)\ \subset\ \bltwoN{or}(c),
\]
and a fourth one $\bltwoN{geo}(c)$ refining $\bltwoN{or}(c)$ by a goodness condition
(\cref{pf:baselat2:ext}). The \emph{origin parameter} $\bltwop_{\mathrm{or}}(c)$ is the parameter defined by
$\bltwoN{or}(c)$, the trivial spectral extension and the trivial geometric predicate. The main theorem
(\cref{pf:baselat2:main}) states: under four packages of literature statements and of results proved
elsewhere, the existence of unlocalised saturated states, and the multiplicity condition for
$\bltwop_{\mathrm{or}}(c)$, both $\sfname{MemberE}(\bltwop_{\mathrm{or}}(c))$ and $\sfname{ExBL}(\bltwop_{\mathrm{or}}(c))$ hold.

\paragraph{The idea.} The theorem is deduced from the member-existence theorem for a \emph{pair} of
extensions, \lref{GlobalTheory.Mem.Top.Bl.exists_baseLattice_of_pair_alt}, applied to the pair
$(\bltwoN{or}(c),\bltwoN{geo}(c))$. That theorem needs four inputs which are produced here:
\begin{enumerate}[label=(\roman*)]
\item \emph{level-three data at starting witnesses}: eventually, for large twisting data, every starting
  column witness has a level carrying level-three data for $\bltwop_{\mathrm{or}}(c)$;
\item the \emph{level-three movement conditions} for $\bltwop_{\mathrm{or}}(c)$ and
  $\bltwop_{\mathrm{geo}}(c)$: level-three data survive the replacement of the spectral stage, of the
  division lattices and of the heights of a level;
\item the \emph{second pass}: every built core for $\bltwop_{\mathrm{or}}(c)$ satisfying the export
  condition can be moved, by enlarging the base and re-twisting with the same prime set, to a starting
  witness carrying level-three data for $\bltwop_{\mathrm{geo}}(c)$;
\item the position-movement property for $\bltwop_{\mathrm{or}}(c)$, which is proved in the section of
  \lref{GlobalTheory.T.large_unramified_exponent_bundles_exist_alt}.
\end{enumerate}
Input (i) comes from a \emph{chain of stages}. To a starting column witness the development attaches
pre-data $P$ and \emph{states} $x$ over $P$, and to a state a sequence of increasingly rich
\emph{levels}: a same-Galois level, an intermediate level, three further intermediate levels (settings
and forms, series, moduli points) and the strict-Henselian-acyclic level, which carries the
height-counting character data. The \emph{all-stages theorem} (\cref{pf:baselat2:allstages}) asserts
that each step of the chain can be carried out; composing the steps, starting from an unlocalised
saturated state, gives (i) (\cref{pf:baselat2:has3}). The origin extension is designed so that the
norm-choice datum of a built core remembers the state from which it came; this is what makes the
second pass (iii) possible (\cref{sec:pf:baselat2:second}): from a built core one recovers a state with a
strict-Henselian-acyclic level, constructs new twisting characters from a \emph{seed} by global
class field theory (\cref{pf:baselat2:twlift}), carries the level to the enlarged base and re-applies
the chain.

\paragraph{Organisation.} \Cref{sec:pf:baselat2:conv} fixes the vocabulary (eventual statements,
states, levels, extensions, stage statements). \Cref{sec:pf:baselat2:main} states and proves the main
theorem from the results of the later subsections. \Cref{sec:pf:baselat2:chain} composes the chain of
stages. \Cref{sec:pf:baselat2:allstages} proves the all-stages theorem from the stage constructions,
each of which is reduced to its arithmetic ingredients. \Cref{sec:pf:baselat2:second} proves spectral
finiteness and the second pass, and \cref{sec:pf:baselat2:lv3} the level-three movement conditions.
\Cref{sec:pf:baselat2:ingr} proves the principal ingredients of the stage constructions and of the second
pass (true types over ray-class primes, existence of the automorphic side, forcing by large data and the eta
equality, twisted forms, series data on twisted copies, geometric laws, character data, the comparison of levels
and the movement to the base). The remaining subsections contain complete proofs of arithmetic results which
are used in the stage constructions and elsewhere: a coefficient bound for theta series and the integrality of
Hecke operators on sup lattices (\cref{sec:pf:baselat2:theta}), a bound for the order of the relative different
(\cref{sec:pf:baselat2:diff}), additivity of congruence lengths, a base-change property of lattices and the
equality of local and rational congruence lengths (\cref{sec:pf:baselat2:len}), and auxiliary lemmas on local
fields, Galois groups and characters (\cref{sec:pf:baselat2:aux}).

Several intermediate steps of the stage constructions are cited through their formal statements and
are not converted here; their content is paraphrased where they are used. The formal development is
\cite{LeanSource2026}.

% ---------------------------------------------------------------------------
\subsection{Conventions}\label{sec:pf:baselat2:conv}

\paragraph{Scale contexts, bases, frames, twisting data.} We use the vocabulary of
\cref{sec:families:nonempty} and of the sections of
\lref{GlobalTheory.Mem.Top.Bl.exists_baseLattice_of_pair_alt} and
\lref{GlobalTheory.T.large_unramified_exponent_bundles_exist_alt}: $K$ is a \emph{scale context}
(a scale $(x,\eps,J)$, a constant record and auxiliary data), $B$ a \emph{base} (a pre-base, i.e.\ ranks,
a CM host $E/F$ with auxiliary primes $g_0,g_1,d$, Hermitian spaces and words, together with a
coefficient field $L=B.\mathrm{LL}$ with an embedding $L\to\C$ and subfields $K_s\subseteq L_0\subseteq L$),
$f$ a \emph{frame} of the pre-base of $B$, and $T$ a \emph{twisting datum} over $B$: a finite set
$T.\mathrm{pl}$ of rational primes and, for each place $v$ of $L$ with residue characteristic
$\ell_v$, a finite set $Q_v$ of primes of $\OO_E$, an exponent $k_v\in\N$ and a Hecke character
$\psi_v$ of $E$. The \emph{untwisted} datum $\mathrm{untw}(B,\mathrm{pl})$ has $Q_v=\emptyset$, $k_v=0$,
$\psi_v=1$. $B'$ \emph{enlarges} $B$ if they have the same pre-base, the same subfields $K_s,L_0$ in $\C$
and $L\subseteq L'$ in $\C$. The base predicate
\[
  \kappa_0(B,f)\ :\Longleftrightarrow\ \bigl(\exists\delta,\ K_s \text{ is exact for } (\delta,f)\bigr)\ \wedge\
  \text{the pre-base of $B$ is pinned}
\]
is the one used in $\sfname{MemberE}$.

\paragraph{Eventual statements.} For a predicate $p$ on scale contexts, $\mathcal{E}(p)$ (``$p$ holds
eventually in the scale'') means: there is $\eps_0>0$ such that for all $\eps\in(0,\eps_0]$ with $\eps<1/52$
and all constants $C_8\ge1$, $C_5,c_6,C_b$ there is a transported constant $K_0$ such that for every standard
constant record $C$ whose constant $C_{\mathrm{pde}}$ equals $C_8$ and which is $K_0$-large there is $c_1$ such that for all
$c_D\ge c_1$ there is $c_2$ such that for all $c_h\ge c_2$ there is $x_0$ with $p(K)$ for every
$(C_5,c_6,C_b)$-good context $K$ with these constants and $K.x\ge x_0$; $\mathcal E_0(p)$ is the same
statement without the constants $C_5,c_6,C_b$ and the goodness premise. For a predicate $p(B,f,T)$,
$\mathrm{lg}_2(K,p)$ (``$p$ holds for large twisting data'') means that there are finite sets
$S_1(B,f,\cdot),S_2(B,f,\cdot,\cdot,\cdot)$ of algebraic numbers, a finite set $P_1(B,f,\cdot,\cdot)$ of primes
and finite sets $\mathcal A(B,f,\cdot,\cdot)(\ell)$ of primes, depending on the pre-base and the frame (and on
the subfields $K_s$, $L_0$ and the character set of $T$), such that $p(B,f,T)$ holds for all $B,f,T$ over $K$
whose fixed coefficient field $L_0$ contains $S_1$ and whose coefficient field contains $S_2$, with
$P_1\subseteq T.\mathrm{pl}$ and $T$ avoiding $\mathcal A$ (no prime of $Q_v$ has absolute norm in
$\mathcal A(\ell_v)$). Then $\sfname{EvLgG}(p)=\mathcal{E}(K\mapsto\mathrm{lg}_2(K,p_K))$, and for a base
predicate $\kappa$
\[
  \mathrm{Step}(\kappa,p,q)=\sfname{EvLgG}\bigl(\kappa(B,f)\wedge\mathrm{Forced}(K,B,T)\wedge p\ \Rightarrow\ q\bigr),
\]
where $\mathrm{Forced}(K,B,T)$ says that $T.\mathrm{pl}$ contains $g_0,g_1,d$ and every prime $p\le X''(n)$
of the scale. For predicates $p,q$ on column witnesses (\emph{witness predicates}) put
\begin{equation}\label{eq:pf:baselat2:stepw}
  \mathcal{W}(p,q)\ :=\ \mathrm{Step}\bigl(\kappa_0,\ \top,\ \forall w,\ p(K,B,f,T,w)\Rightarrow q(K,B,f,T,w)\bigr):
\end{equation}
\emph{eventually in the scale, for large twisting data with $\kappa_0(B,f)$ and $T$ forced, every
column witness $w$ over $(B,f,T)$ with $p(w)$ satisfies $q(w)$.} We use three closure properties
proved in the development: if $\mathcal{W}(p,q)$ and $\mathcal{W}(p,r)$ then $\mathcal{W}(p,q\wedge r)$
(\lref{GlobalTheory.Mem.Top.Bl.StepW.and}); if $\mathcal{W}(p,q)$ and $\mathcal{W}(q,r)$ then
$\mathcal{W}(p,r)$ (\lref{GlobalTheory.Mem.Top.Bl.StepW.trans}); and $\mathcal{W}(p,q)$ holds as
soon as $p(w)\Rightarrow q(w)$ for all $K,B,f,T,w$ with $K$ over the pre-base of $B$ and $T$ forced
(\lref{GlobalTheory.Mem.Top.Bl.StepW.of_all}). Consequently $\mathcal{W}(p,q)$ may be proved by
proving finitely many statements $\mathcal{W}(p,q_i)$ and deducing $q$ from $q_1\wedge\dots\wedge q_n$ for
each witness; all assembly proofs below have this form, and we say that \emph{$q$ follows pointwise from
the $q_i$}.

\paragraph{Witnesses, pre-data and states.} Column witnesses $w$ over $(B,f,T)$ and their \emph{levels}
$Z\in w.\mathrm{Lt}$ are as in the section of
\lref{GlobalTheory.Mem.Top.Bl.exists_baseLattice_of_pair_alt}; a level $Z$ determines stages
$(W,V,S,N,G)$ and a global datum $w.D(Z)$. A witness is \emph{pinned} if it carries a column pin; it has
\emph{good data} if some level carries a fine real datum satisfying the conditions of a pre-datum (below);
and it is a \emph{starting witness}, $\mathrm{start}(w)$, if it is pinned, $T$ is untwisted or $w$ has good
data, and $(B,T)$ is \emph{coherent}: $(Q_v,k_v,\psi_v)$ depends only on the restriction of $v$ to $L_0$, and
$L$ contains a primitive $\ell_v^{k_v}$-th root of unity for every $v$.

A \emph{pre-datum} $P$ over $(K,w)$ consists of a level $Z_0$, a set $\bltwoWt=P.\bltwoWt$ of places of $L$, a
real datum $z$ over $(w,Z_0,T,\bltwoWt)$ and a row package, subject to group, coherence, geometric and
standardness conditions, to $\bltwoWt$ being the set of places determined by $z$, to $\ell_v\in T.\mathrm{pl}$
for $v\in\bltwoWt$ and to $\psi_v=1$ for $v\notin\bltwoWt$. A \emph{state} $x$ over $P$ consists of an
\emph{over-datum} $x.o$ (a level $x.o.Z$ with the same $G$-stage as $Z_0$, a map $\eta$ and stages $V,S,N$), a
Bloch--Kato law at $w.D(x.o.Z)$, a moved datum, a side datum and a pinned read package; $\bltwoSt(P)$ is the
set of states over $P$, and we write $D_x=w.D(x.o.Z)$. For a predicate $\varphi$ on states put
\[
  \forall^{\mathrm{st}}(\varphi):=\mathcal{W}\bigl(\mathrm{start},\ w\mapsto \forall P\ \forall x\in\bltwoSt(P),\ \varphi(x)\bigr).
\]
A state $x$ is \emph{unlocalised} if its set of localisation primes is empty. A pre-datum $P$ is
\emph{saturated}, $\mathrm{Sat}(P)$, if, when $T$ is untwisted, $\bltwoWt$ is closed under the relation ``lying
over one and the same place'' of $w.D(Z_0)$. For $c\in\R$, $\mathrm{List}_c(x)$ is the \emph{list pin} of
the level characters $\lambda_x$ of $x$ at level $(1,c,b)$, $b$ the scale parameter: a grid datum for
$\lambda_x$ whose values lie in $L_0$ and whose $\lambda$-tie is full.

\paragraph{Levels.} Let $x$ be a state. The development defines types of \emph{levels} of $x$, each
extending the previous one:
\begin{enumerate}[label=(L\arabic*)]
\item a \emph{same-Galois level} $\bltwoL{sg}(x)$: height-counting base data
  (a host tie, an embedding $j_E:E\to\Qbar$ compatible with the frame, a character-frame embedding $A$,
  the statements that the Galois side of $A$ is \emph{true} and that its automorphic side is realised,
  the representation parameters of the rows, a frame for $\Omega$, Hodge--Tate data, a row depth, an
  elementary datum, a valuation datum, labels, generators $\zeta_k$ of the label data at $g_k$, and the
  injectivity of $\OO\to L$ with integral range), together with a map $P_q$ from the primes of the
  maximal real subfield of $E$ to primes of $E$ above them whose range is the selection of $w$;
\item an \emph{intermediate level} $\bltwoL{int}(x)$: a same-Galois level together with a \emph{module
  input} (a family $M$ of modules adapted to the rows, with input characters $\mu_0\xi$), a
  \emph{rationality datum}, and \emph{row cuspidality} (\cref{pf:baselat2:cusp});
\item \emph{first, second and third intermediate levels} $\bltwoL{1}(x),\bltwoL{2}(x),\bltwoL{3}(x)$:
  an intermediate level together with settings-and-forms data, series data, respectively moduli-points
  data;
\item the \emph{strict-Henselian-acyclic level} $\bltwoL{sha}(x)$: a third intermediate level $l$
  together with \emph{height-counting character data} $c$ for $(l.\mathrm{mt},l.b,T.k,T.\psi)$; $c$ contains a
  norm choice $c.\mathrm{ch}$.
\end{enumerate}
Two norm choices are \emph{kept}, $c\approx c'$, if their data have the same Galois data, agreeing
twist data and forms, and the same twisted indices, signs, sets $Q$, exponents and characters.

\begin{definition}[Norm-choice extensions]\label{pf:baselat2:ext}
Fix $h_L$ and $c\in\R$. The norm-choice extensions used here assign to a configuration
$(K,B,f,T,w,Z,\bltwoWt,p_0,t,\mathrm{nch},n,x_d)$ with a norm-choice core the following types:
\begin{enumerate}[label=(\alph*)]
\item $\bltwoN{law}$: the default extension (real expansion data of the row) times the Bloch--Kato law
  for $w.D(Z)$ (rational map values imply rationality);
\item $\bltwoN{list}(c)=\bltwoN{law}\times\mathrm{Pin}_c$, where $\mathrm{Pin}_c$ is the list pin of the
  level characters of the rows $t$ at level $(1,c,b)$;
\item $\bltwoN{or}(c)=\bltwoN{list}(c)\times\mathrm{Orig}_c$, where the \emph{origin datum}
  $\mathrm{Orig}_c(K,w,\bltwoWt,\mathrm{nch})$ says: there are a pre-datum $P$ over $(K,w)$, a state $x$
  over $P$ and a strict-Henselian-acyclic level $l$ of $x$ with $\mathrm{List}_c(x)\wedge\mathrm{Sat}(P)$,
  $\bltwoWt=P.\bltwoWt$ and $l.c.\mathrm{ch}\approx\mathrm{nch}$;
\item $\bltwoN{geo}(c)=\bltwoN{or}(c)\times\mathrm{Good}_O$, where $\mathrm{Good}_O(w.D(Z),\mathrm{nch},x_d)$ is
  the goodness condition on norm choices and theta data.
\end{enumerate}
For a norm-choice extension $N$, a spectral extension $S$ and a geometric predicate $\Gamma$ we write
$\bltwop(N,S,\Gamma)$ for the parameter they define and $\bltwop(N,S,\Gamma).s$ for its slot family.
Let $S_0$ be the trivial spectral extension and $\Gamma_0$ the closure of the trivial geometric predicate;
put $\bltwop(N)=\bltwop(N,S_0,\Gamma_0)$, $\bltwop_{\mathrm{or}}(c)=\bltwop(\bltwoN{or}(c))$ and
$\bltwop_{\mathrm{geo}}(c)=\bltwop(\bltwoN{geo}(c))$.
\formallabel{GlobalTheory.HeightCounting.Up.GeometricNormExtension}
\end{definition}

\begin{definition}[Level-three data]\label{pf:baselat2:lv3def}
Let $\mathcal S$ be a slot family, $w$ a column witness and $Z$ a level of $w$. \emph{Level-two data}
$\mathrm{Lv}_2(\mathcal S,K,w,Z)$ consist of the $G$-condition of $(K,w,Z)$, coherence of $(B,T)$, a set
$\bltwoWt$ of places, a base-level row side $p_0$ of $w.D(Z)$, rows $t$ and an element of the row
requirement $\mathcal S.R_x$. \emph{Level-three data} $\mathrm{Lv}_3(\mathcal S,K,w,Z)$ add a norm choice
$\mathrm{nch}$, norm-choice data $n$, theta data $x_d$, an element of the norm requirement $\mathcal S.N_x$,
and the condition that for $k\in\{0,1\}$ and each place $v$ of type $k$ there is $u\in(\Z/g_k\Z)^\times$ of
order $g_k-1$ with $\chi_T(k,v,j)=u^{2(j+1)}$ for all $j<m$. A witness $w$ \emph{has level three},
$\mathrm{has}_3(\mathcal S)(w)$, if $\mathrm{Lv}_3(\mathcal S,K,w,Z)\ne\emptyset$ for some level $Z$ of $w$.
The \emph{level-three movement condition} $\mathrm{Mv}_3(\mathcal S)$ means: for all $K,B,f,T,w,Z$, all
$l\in\mathrm{Lv}_3(\mathcal S,K,w,Z)$, every spectral stage $S'$, every family $N=(N_u)_u$ of
$\OO_L$-submodules of the division modules of $Z$ and every function $h$ (the \emph{heights}) on the words
of $B$, there is $l'\in\mathrm{Lv}_3(\mathcal S,K,w,\mathrm{up}(Z,S',N,h))$ whose package is the update of the
package of $l$ by $(Z,S',N,h)$; here $\mathrm{up}(Z,S',N,h)$ is the level $Z$ with spectral stage $S'$,
division lattices $N$ and heights $h$.
\formallabel{GlobalTheory.Mem.Top.Bl.Lv2}
\formallabel{GlobalTheory.Mem.Top.Bl.Lv3}
\formallabel{GlobalTheory.Mem.Top.Bl.HasZ}
\formallabel{GlobalTheory.Mem.Top.Bl.Has3}
\formallabel{GlobalTheory.Mem.SRe.LevelThreeCondition}
\end{definition}

\begin{definition}[Stage statements]\label{pf:baselat2:stagedefs}
Let $h_L\in\sfname{RecTExists}$ and $c\in\R$. All statements are of the form $\forall^{\mathrm{st}}(\varphi)$, and we
only give $\varphi(x)$, for a state $x$ over a pre-datum $P$ over $(K,B,f,T,w)$:
\begin{enumerate}[label=(\arabic*)]
\item $\bltwoS{sg}$: if $x$ is unlocalised, there is a state $x'$ over $P$ with a same-Galois level and with
  the same Galois data as $x$;
\item $\bltwoS{mt}$: $x$ admits a module input;
\item $\bltwoS{rat}$: every same-Galois level of $x$ gives a rationality datum;
\item $\bltwoS{cusp}$: for every same-Galois level $l$ of $x$, the row of $l$ is cuspidal;
\item $\bltwoS{1}$, $\bltwoS{2}$, $\bltwoS{3}$: if $x$ has an intermediate (first, second) intermediate level,
  some state $x'$ over $P$ has a first (second, third) intermediate level;
\item $\bltwoS{sha}$: every third intermediate level $l$ of $x$ admits height-counting character data for
  $(l.\mathrm{mt},l.b,T.k,T.\psi)$;
\item $\bltwoS{fin}$: for every strict-Henselian-acyclic level $l$ of $x$ and every $(S,\Gamma)$ there are a
  level $Z$ of $w$ and $v\in\mathrm{Lv}_3(\bltwop(\bltwoN{law},S,\Gamma).s,K,w,Z)$ with $v.\bltwoWt=P.\bltwoWt$
  and $l.c.\mathrm{ch}\approx v.\mathrm{nch}$;
\item $\bltwoS{1}^{c}$: if $x$ has an intermediate level, some $x'$ over $P$ has a first intermediate level
  and satisfies $\mathrm{List}_c(x')$;
\item $\bltwoS{2}^{c}$, $\bltwoS{3}^{c}$: if $x$ has a first (second) intermediate level, some $x'$ over $P$
  has a second (third) intermediate level, and $\mathrm{List}_c(x)\Rightarrow\mathrm{List}_c(x')$;
\item $\bltwoS{fin}^{c}$: as (7), with $\bltwoN{list}(c)$ in place of $\bltwoN{law}$, under the premise
  $\mathrm{List}_c(x)$;
\item $\bltwoS{spec}^{c}$: as (7), with $\bltwoN{or}(c)$ in place of $\bltwoN{law}$, under the premise
  $\mathrm{List}_c(x)\wedge\mathrm{Sat}(P)$;
\item $\mathbf{S}_{\eta}$: for every strict-Henselian-acyclic level $l$ of $x$, if the automorphic side
  exists for the character-frame embedding, the module family, the norm choice and the tuple of $l$, then this
  side is forced by
  automorphic semisimple data $D'$ which are \emph{large} and satisfy $D'.\eta_J(j,v)=P.q.\eta(j,v)$ for
  $1\le j\le m$ and all places $v$;
\item $\mathbf{S}_{\theta}$: every second intermediate level of $x$ gives theta data $x_d$ for $D_x$ which
  are \emph{full}: for every index $(i,j,b)$ with $1\le j\le m$ they select all data.
\end{enumerate}
Finally, the \emph{second pass} $\mathbf{P}^{c}$ is the statement: for all $(S,\Gamma)$,
\[
  \mathrm{Pass}\Bigl(\kappa_0,\ \mathrm{Built}(\bltwop(\bltwoN{or}(c),S,\Gamma).s),\ (K,B,f,T)\mapsto
  \exists w,\ \mathrm{start}(w)\wedge\mathrm{has}_3(\bltwop(\bltwoN{geo}(c),S,\Gamma).s)(w)\Bigr),
\]
where $\mathrm{Pass}(\kappa,Q,r)$ is as in the section of
\lref{GlobalTheory.Mem.Top.Bl.exists_baseLattice_of_pair_alt}: eventually, for large untwisted
$T$ with $\kappa(B,f)$, for every core $c$ with $Q$ whose assembled datum satisfies the export condition,
for every finite exclusion function $\mathcal A$ and every finite set $S_2(\cdot)$ of algebraic numbers,
there are $B'$ enlarging $B$, the same frame $f'=f$ and $T'$ with $T'.\mathrm{pl}=T.\mathrm{pl}$, $T'$ avoiding
$\mathcal A$ and $S_2$ realised in $B'.\mathrm{LL}$, such that $r(K,B',f',T')$.
\formallabel{GlobalTheory.HeightCounting.Up.SameGaloisLevelExists}
\formallabel{GlobalTheory.HeightCounting.Up.RationalityStatementThirdPlaceSeries}
\formallabel{GlobalTheory.HeightCounting.Up.CuspStatement'}
\formallabel{GlobalTheory.HeightCounting.Up.FirstIntermediateLevelExists}
\formallabel{GlobalTheory.HeightCounting.Up.SecondIntermediateLevelExists}
\formallabel{GlobalTheory.HeightCounting.Up.ThirdIntermediateLevelExists}
\formallabel{GlobalTheory.HeightCounting.Up.StrictHenselianAcyclicLevelNonempty}
\formallabel{GlobalTheory.HeightCounting.Up.FinitenessStatement}
\formallabel{GlobalTheory.HeightCounting.Up.SettingsFormsExistList'}
\formallabel{GlobalTheory.HeightCounting.Up.SeriesExistList'}
\formallabel{GlobalTheory.HeightCounting.Up.ModuliPointsExistList'}
\formallabel{GlobalTheory.HeightCounting.Up.FinitenessListStatement}
\formallabel{GlobalTheory.HeightCounting.Up.SpectralFiniteness}
\formallabel{GlobalTheory.HeightCounting.Up.EtaStatement'}
\formallabel{GlobalTheory.HeightCounting.Up.ThetaDataExists}
\formallabel{GlobalTheory.HeightCounting.Up.SecondNonvanishing}
\end{definition}

% ---------------------------------------------------------------------------
\subsection{The main theorem}\label{sec:pf:baselat2:main}

The \emph{existence of unlocalised saturated states} is the statement
\[
  \mathbf{U}:=\mathcal{W}\bigl(\mathrm{start},\ w\mapsto\exists P\ \exists x\in\bltwoSt(P),\
  \mathrm{Sat}(P)\wedge x\text{ unlocalised}\wedge\forall j,v\ \ x.o.\eta(j,v)=0\bigr),
\]
proved in the section of \lref{GlobalTheory.T.large_unramified_exponent_bundles_exist_alt}. The
packages $\mathcal{H}_{A}$, $\mathcal{H}_{\mathrm{ell}}$, $\mathcal{H}_{\mathrm{mc}}$ and
$\mathcal{H}_{\mathrm{LL}}(h_L)$ are the conjunctions of literature statements and of statements
proved elsewhere described in that section: $\mathcal H_A$ collects statements $a_1,\dots,a_{23}$ on ray class
fields, principal ideals, CM elliptic curves, Harish-Chandra modules, rapid decay of cusp forms, the
Casselman--Wallach theorem, the Weil representation, discrete series, Siegel--Weil, formal degrees, the
Siegel mass ratio and global Artin reciprocity; $\mathcal H_{\mathrm{ell}}$ collects Galois-cohomological
and automorphic statements; $\mathcal H_{\mathrm{mc}}$ contains the discrete series cover, the published
theta-condition theorems ($h_1$), the continuation of partial Hecke $L$-functions ($h_3$) and two
multiplicity statements at pinned and at starting witnesses; and $\mathcal{H}_{\mathrm{LL}}(h_L)$ consists of
the published third-layer inputs $\mathcal I_3$, the series inputs $\mathcal I_{\mathrm{ser}}$, an archimedean
comparison $\mathcal C_\infty(h_L)$ (which contains a Hodge--Tate comparison $\mathcal C_{\mathrm{HT}}$), a
parameter statement $\mathrm{Par}(h_L)$, a point statement $\mathrm{Pt}(h_L)$ and a member statement for the
rows of layer points. We name the fields when they are used.

\begin{theorem}[Member existence along the origin extension]\label{pf:baselat2:main}
Assume $\mathcal{H}_{A}$, $\mathcal{H}_{\mathrm{ell}}$, $\mathcal{H}_{\mathrm{mc}}$ and
$\mathbf{U}$. Let $h_L\in\sfname{RecTExists}$ with
$\mathcal{H}_{\mathrm{LL}}(h_L)$, let $c\in\R$, and assume the multiplicity condition
$\mathbf{M}(\bltwop_{\mathrm{or}}(c))$. Then
\[
  \sfname{MemberE}(\bltwop_{\mathrm{or}}(c).s)\quad\text{and}\quad
  \sfname{ExBL}\bigl(\kappa_0,\ \mathrm{Built}(\bltwop_{\mathrm{or}}(c).s),\ \mathrm{Good}(\bltwop_{\mathrm{or}}(c).s),\
  \mathrm{Lk}(\bltwop_{\mathrm{or}}(c))\bigr).
\]
\formallabel{GlobalTheory.Mem.Top.Bl.exists_baseLattice_of_prime_alt}
\end{theorem}

\begin{proof}
\emph{Step 1: two auxiliary conditions.} The fields $h_1$ (published theta-condition theorems), $h_3$
(continuation of partial Hecke $L$-functions) and the multiplicity statement at starting witnesses of
$\mathcal{H}_{\mathrm{mc}}$ give the \emph{weight-case condition} $\mathbf{W}$ (eventually, for large twisting
data, at every starting witness every real datum which satisfies the group, $\kappa$-definition, geometric and
core conditions and whose places are those it determines satisfies the base-change condition for real data)
by \lref{GlobalTheory.T.w_of_weight_residue_char}, which derives $\mathbf W$ from exactly these three
statements. The discrete series cover contained in $\mathcal H_{\mathrm{mc}}$
and the fields $a_5$ (Harish-Chandra modules), $a_6$ (rapid decay of cusp forms), $a_7$
(Casselman--Wallach), $a_8$ (Weil representation of unitary pairs), $a_9$ (existence of discrete series)
and $a_{10}$ (Harish-Chandra finiteness) of $\mathcal{H}_{A}$ give the extension property
$\mathbf{X}$ of fine adelic data by \cref{pf:baselat2:datext}.

\emph{Step 2: the stages.} \Cref{pf:baselat2:allstages}, applied with $a_7,a_5,a_8,a_9$, the field $a_{12}$
(uniqueness of discrete series of $\Uni(p,q)$), $\mathcal{H}_{\mathrm{mc}}$ and the fields
$\mathcal I_3,\mathcal I_{\mathrm{ser}},\mathcal C_\infty(h_L),\mathrm{Pt}(h_L)$ of $\mathcal{H}_{\mathrm{LL}}(h_L)$, gives the
stage statements $\bltwoS{sg}$, $\bltwoS{mt}$, $\bltwoS{rat}$, $\bltwoS{cusp}$,
$\bltwoS{sha}$ and, for our $c$, $\bltwoS{1}^{c}$, $\bltwoS{2}^{c}$, $\bltwoS{3}^{c}$ and
$\bltwoS{fin}^{c}$.

\emph{Step 3: spectral finiteness and the second pass.} \Cref{pf:baselat2:sfino}, applied with
the discrete series cover, $a_7$, $a_5$, $\mathcal I_3$, $\mathcal I_{\mathrm{ser}}$, $\mathbf{X}$, $\mathbf{W}$,
$a_9$, the comparison $\mathcal{C}_{\mathrm{HT}}$, $\bltwoS{mt}$, $\bltwoS{rat}$, $\bltwoS{cusp}$ and
$\bltwoS{fin}^{c}$, gives $\bltwoS{spec}^{c}$ and $\mathbf{P}^{c}$.

\emph{Step 4: level three at starting witnesses.} \Cref{pf:baselat2:has3}, applied to
$\mathbf{U}$ and the statements of Steps 2 and 3, gives
$\mathcal{W}(\mathrm{start},\mathrm{has}_3(\bltwop_{\mathrm{or}}(c).s))$.

\emph{Step 5: movement.} \Cref{pf:baselat2:lv3mv} gives $\mathrm{Mv}_3(\bltwop_{\mathrm{or}}(c).s)$
and $\mathrm{Mv}_3(\bltwop_{\mathrm{geo}}(c).s)$, and
\lref{GlobalTheory.T.slot_mv_extension_with_origin_geo_cl}, which states the position-movement
property for $\bltwop(\bltwoN{or}(c),S_0,\overline\Gamma)$ for every geometric predicate $\Gamma$ with closure
$\overline\Gamma$ (positions of witnesses move along the replacement of division lattices and heights), gives it
for $\bltwop_{\mathrm{or}}(c)$.

\emph{Step 6: conclusion.} Apply the theorem on pairs of extensions,
\lref{GlobalTheory.Mem.Top.Bl.exists_baseLattice_of_pair_alt}, with $N=\bltwoN{or}(c)$ and
$N_g=\bltwoN{geo}(c)$. That theorem states: for norm-choice extensions $N,N_g$ such that every $N$-datum
yields a $\bltwoN{law}$-datum and every $N_g$-datum yields an $N$-datum together with $\mathrm{Good}_O$, the
packages $\mathcal{H}_{A}$, $\mathcal{H}_{\mathrm{ell}}$, $\mathcal{H}_{\mathrm{mc}}$,
$\mathcal{H}_{\mathrm{LL}}(h_L)$, the statement $\mathcal{W}(\mathrm{start},\mathrm{has}_3(\bltwop(N).s))$,
the two level-three movement conditions for $\bltwop(N).s$ and $\bltwop(N_g).s$, the position-movement property for $\bltwop(N)$,
the second pass
$\mathrm{Pass}(\kappa_0,\mathrm{Built}(\bltwop(N).s),\exists w,\ \mathrm{start}(w)\wedge\mathrm{has}_3(\bltwop(N_g).s)(w))$
and $\mathbf{M}(\bltwop(N))$ imply $\sfname{MemberE}(\bltwop(N).s)$ and
$\sfname{ExBL}(\kappa_0,\mathrm{Built}(\bltwop(N).s),\mathrm{Good}(\bltwop(N).s),\mathrm{Lk}(\bltwop(N)))$.
The two projection hypotheses hold by \cref{pf:baselat2:ext}: an element of
$\bltwoN{or}(c)=(\bltwoN{law}\times\mathrm{Pin}_c)\times\mathrm{Orig}_c$ has a first component in
$\bltwoN{law}$, and an element of $\bltwoN{geo}(c)=\bltwoN{or}(c)\times\mathrm{Good}_O$ consists of an element
of $\bltwoN{or}(c)$ and a proof of $\mathrm{Good}_O$. The remaining hypotheses are the assumptions of the
theorem and the statements of Steps 3--5 (the second pass is $\mathbf{P}^{c}$ at $(S_0,\Gamma_0)$).
\end{proof}

\begin{remark}\label{pf:baselat2:rem:main}
(1) The parameter statement $\mathrm{Par}(h_L)$ contained in $\mathcal{H}_{\mathrm{LL}}(h_L)$ is
not used, neither in this proof nor in \lref{GlobalTheory.Mem.Top.Bl.exists_baseLattice_of_pair_alt}
(which uses $\mathcal I_3$, $\mathcal I_{\mathrm{ser}}$ and the member statement); inside
\cref{pf:baselat2:allstages} the parameter statement is derived again from the published third-layer
inputs and $\bltwoS{mt}$. (2) $\mathcal{H}_{\mathrm{ell}}$ is used only through
\lref{GlobalTheory.Mem.Top.Bl.exists_baseLattice_of_pair_alt}. (3) The statements
$\bltwoS{1}$, $\bltwoS{2}$, $\bltwoS{3}$ and $\bltwoS{fin}$ produced by
\cref{pf:baselat2:allstages} are not used here; only their list versions are.
\end{remark}

% ---------------------------------------------------------------------------
\subsection{The chain of stages}\label{sec:pf:baselat2:chain}

\begin{proposition}[Level three from the chain of stages]\label{pf:baselat2:has3}
Let $h_L\in\sfname{RecTExists}$ and $c\in\R$. Assume $\mathbf{U}$,
$\bltwoS{sg}$, $\bltwoS{mt}$, $\bltwoS{rat}$, $\bltwoS{cusp}$, $\bltwoS{1}^{c}$,
$\bltwoS{2}^{c}$, $\bltwoS{3}^{c}$, $\bltwoS{sha}$ and $\bltwoS{spec}^{c}$. Then for every
spectral extension $S$ and geometric predicate $\Gamma$,
\[
  \mathcal{W}\bigl(\mathrm{start},\ \mathrm{has}_3(\bltwop(\bltwoN{or}(c),S,\Gamma).s)\bigr).
\]
\formallabel{GlobalTheory.Mem.Top.Bl.has3_neo_of_localLanglands_alt}
\end{proposition}

\begin{proof}
All hypotheses are statements $\mathcal{W}(\mathrm{start},\cdot)$; by the closure properties of
\cref{sec:pf:baselat2:conv} it suffices to show that $\mathrm{has}_3$ follows pointwise. So let $w$ be a
starting column witness over $(K,B,f,T)$ at which the conclusions of all hypotheses hold.
\begin{enumerate}[label=(\alph*)]
\item By $\mathbf{U}$ there are a pre-datum $P$ with $\mathrm{Sat}(P)$ and an
  unlocalised state $x$ over $P$.
\item By $\bltwoS{sg}$ there is a state $x_1$ over $P$ with a same-Galois level $l_{\mathrm{sg}}$.
  By $\bltwoS{mt}$, $x_1$ has a module input; by $\bltwoS{rat}$, $l_{\mathrm{sg}}$ gives a rationality datum;
  by $\bltwoS{cusp}$ the row of $l_{\mathrm{sg}}$ is cuspidal. These four objects form an intermediate level of
  $x_1$.
\item By $\bltwoS{1}^{c}$ there is a state $x_2$ over $P$ with a first intermediate level and
  $\mathrm{List}_c(x_2)$; by $\bltwoS{2}^{c}$ there is $x_3$ over $P$ with a second intermediate level
  and $\mathrm{List}_c(x_2)\Rightarrow\mathrm{List}_c(x_3)$; by $\bltwoS{3}^{c}$ there is $x_4$ over $P$
  with a third intermediate level $l_3$ and $\mathrm{List}_c(x_3)\Rightarrow\mathrm{List}_c(x_4)$. Hence
  $\mathrm{List}_c(x_4)$.
\item By $\bltwoS{sha}$ there are height-counting character data for $l_3$; together with $l_3$ they
  form a strict-Henselian-acyclic level $l$ of $x_4$.
\item Since $\mathrm{List}_c(x_4)$ and $\mathrm{Sat}(P)$ hold, $\bltwoS{spec}^{c}$ applied to $l$ and
  $(S,\Gamma)$ gives a level $Z$ of $w$ and an element of $\mathrm{Lv}_3(\bltwop(\bltwoN{or}(c),S,\Gamma).s,K,w,Z)$.
\end{enumerate}
Thus $\mathrm{has}_3(\bltwop(\bltwoN{or}(c),S,\Gamma).s)(w)$.
\end{proof}

The same composition, without the list condition, is recorded in the development as follows. Let
$\mathrm{bare}(w)$ be the predicate ``there are $P$ and an unlocalised state $x$ over $P$ with
$x.o.\eta=0$''.

\begin{proposition}[Forgetting lists]\label{pf:baselat2:stepspin}
Let $h_L\in\sfname{RecTExists}$ and $c\in\R$.
\begin{enumerate}[label=(\alph*)]
\item $\bltwoS{1}^{c}\Rightarrow\bltwoS{1}$, $\bltwoS{2}^{c}\Rightarrow\bltwoS{2}$ and
  $\bltwoS{3}^{c}\Rightarrow\bltwoS{3}$.
\item $\mathcal{W}(\mathrm{start},\mathrm{bare})$, $\bltwoS{sg}$, $\bltwoS{mt}$,
  $\bltwoS{rat}$, $\bltwoS{cusp}$, $\bltwoS{1}^{c}$, $\bltwoS{2}^{c}$, $\bltwoS{3}^{c}$,
  $\bltwoS{sha}$ and $\bltwoS{fin}^{c}$ imply, for all $(S,\Gamma)$,
  \[
    \mathcal{W}\bigl(\mathrm{start},\mathrm{has}_3(\bltwop(\bltwoN{list}(c),S,\Gamma).s)\bigr).
  \]
\end{enumerate}
\formallabel{GlobalTheory.HeightCounting.UpPin.steps_pin_alt}
\end{proposition}

\begin{proof}
(a) Each implication follows pointwise by forgetting the list condition in the conclusion. (b) As in the
proof of \cref{pf:baselat2:has3}, with $\mathrm{bare}$ in place of (a) there; in step (c) one obtains
$\mathrm{List}_c(x_4)$ from $\mathrm{List}_c(x_2)$ by the two implications, and in step (e) one uses
$\bltwoS{fin}^{c}$, whose premise is $\mathrm{List}_c(x_4)$ only.
\end{proof}

% ---------------------------------------------------------------------------
\subsection{All stages at members}\label{sec:pf:baselat2:allstages}

\begin{theorem}[All stages]\label{pf:baselat2:allstages}
Assume the Casselman--Wallach theorem, the theory of Harish-Chandra modules, the Weil representation of
unitary dual pairs, the existence of discrete series and the uniqueness of discrete series of
$\Uni(p,q)$, in the forms stated in the development, as well as $\mathcal{H}_{\mathrm{mc}}$, $\mathcal{I}_3$ and
$\mathcal{I}_{\mathrm{ser}}$. Let $h_L\in\sfname{RecTExists}$ with $\mathcal{C}_{\infty}(h_L)$ and
$\mathrm{Pt}(h_L)$. Then $\bltwoS{sg}$, $\bltwoS{mt}$, $\bltwoS{rat}$,
$\bltwoS{cusp}$, $\bltwoS{1}$, $\bltwoS{2}$, $\bltwoS{3}$, $\bltwoS{sha}$ and
$\bltwoS{fin}$ hold, and for every $c\in\R$ the statements $\bltwoS{1}^{c}$, $\bltwoS{2}^{c}$,
$\bltwoS{3}^{c}$ and $\bltwoS{fin}^{c}$ hold.
\formallabel{GlobalTheory.T.all_stages_at_member_of_specializations_alt}
\formallabel{GlobalTheory.HeightCounting.Up.all_stages_of_specializations_alt}
\end{theorem}

\begin{lemma}\label{pf:baselat2:mtpar}
$\mathcal I_3$ implies $\bltwoS{mt}$, and $\mathcal I_3$ together with $\bltwoS{mt}$ implies $\mathrm{Par}(h_L)$ for
every $h_L$.
\formallabel{GlobalTheory.T.sMt_of_published_inputs'}
\formallabel{GlobalTheory.T.parSt_of_third_layer_sMt}
\end{lemma}

\begin{proof}
The first statement is \lref{GlobalTheory.MtInMain.prN_sMt_alt_alt}, which derives the existence of
module inputs at all states from the published third-layer inputs. The second is
\lref{GlobalTheory.HeightCounting.Up.conjDualParam_stable_of_prime_not_dividing_level_alt}, which
derives the parameter statement (stability of conjugate-dual local parameters at primes not dividing the
level) from the same inputs and $\bltwoS{mt}$.
\end{proof}

\begin{proof}[Proof of \cref{pf:baselat2:allstages}]
\begin{enumerate}[label=(\arabic*)]
\item $\bltwoS{mt}$ and $\mathrm{Par}(h_L)$ hold by \cref{pf:baselat2:mtpar}.
\item $\bltwoS{rat}$ holds by \cref{pf:baselat2:rat,pf:baselat2:isot}.
\item $\mathbf{S}_{\eta}$ holds by \cref{pf:baselat2:eta}, from $\mathcal I_3$ and $\mathrm{Pt}(h_L)$.
\item $\bltwoS{fin}$ and $\bltwoS{fin}^{c}$ for all $c$ hold by \cref{pf:baselat2:fin}, from the Weil
  representation, the existence of discrete series, $\mathcal I_3$, $\mathcal{C}_{\infty}(h_L)$,
  $\mathbf{S}_{\eta}$ and $\mathbf{S}_{\theta}$ (\cref{pf:baselat2:sxd}).
\item $\bltwoS{sg}$ holds by \cref{pf:baselat2:sa}, from $\mathcal I_3$ and the comparison
  $\mathcal{C}_{\mathrm{HT}}$ contained in $\mathcal{C}_{\infty}(h_L)$.
\item $\bltwoS{cusp}$ holds by \cref{pf:baselat2:cusp}, from $\mathcal I_3$, $\mathrm{Par}(h_L)$
  and $\bltwoS{rat}$.
\item For every $c$, $\bltwoS{1}^{c}$, $\bltwoS{2}^{c}$ and $\bltwoS{3}^{c}$ hold by
  \cref{pf:baselat2:settings,pf:baselat2:series,pf:baselat2:points}, from $\mathcal I_3$ and $\mathcal I_{\mathrm{ser}}$.
  By \cref{pf:baselat2:stepspin}(a) with $c=0$ they give $\bltwoS{1}$, $\bltwoS{2}$ and
  $\bltwoS{3}$.
\item $\bltwoS{sha}$ holds by \cref{pf:baselat2:sha}, from the discrete series cover contained in
  $\mathcal{H}_{\mathrm{mc}}$, the Casselman--Wallach theorem, Harish-Chandra modules, $\mathcal I_3$ and
  $\mathcal I_{\mathrm{ser}}$.\qedhere
\end{enumerate}
\end{proof}

\begin{remark}\label{pf:baselat2:rem:upq}
The uniqueness of discrete series of $\Uni(p,q)$ is not used: in the formal proof it is passed to the
rationality statement, whose proof discards it. In this proof $\mathcal{H}_{\mathrm{mc}}$ is used only through
the discrete series cover and the multiplicity statement at pinned witnesses (the latter in
\cref{pf:baselat2:isot}).
\end{remark}

We now prove the stage statements. Each proof is a pointwise deduction, in the sense of
\cref{sec:pf:baselat2:conv}, from statements $\mathcal{W}(\mathrm{start},\cdot)$ provided by intermediate
results of the development, which we cite with their content.

\subsubsection{The same-Galois level}

\begin{proposition}\label{pf:baselat2:sa}
$\mathcal{I}_3$ and $\mathcal{C}_{\mathrm{HT}}$ imply $\bltwoS{sg}$ for every
$h_L$.
\formallabel{GlobalTheory.HeightCounting.Up.sa_of_prN_alt_alt}
\formallabel{GlobalTheory.T.stageASameGalois_of_inputs_phtCmp}
\end{proposition}

\begin{proof}
Three statements are used.
(i) \lref{GlobalTheory.T.isRayPrime_g_and_integral_range}: at every state $x$, the datum $D_x$ has an
elementary datum and a valuation datum, the map $\OO\to L$ is injective with range the integral elements
of a subfield, and there are a host tie $h$ and an embedding $j_E:E\to\Qbar$ agreeing with the frame of
$h$ such that each auxiliary prime $g_k$ is a \emph{ray prime} for the field cut out by $j_E$ and the
modulus $QS\cdot q\cdot d$.
(ii) \lref{GlobalTheory.T.depth_of_repOf_of_inputs}: at every state the rows have the depth property.
(iii) \cref{pf:baselat2:raytypes}, proved in \cref{sec:pf:baselat2:ingr} from the published inputs
contained in $\mathcal I_3$.
By the closure properties it suffices to argue pointwise. Let $x$ be an unlocalised state over $P$. The
types of $D_x$ are true for the frame $f$ (a condition contained in $x$), and the exponents are integral
by the elementary datum of (i). So (iii) applies and gives $(L,A)$. By
\lref{GlobalTheory.T.heckeGaloisStage_fits_of_same_centre_actions} the stage $N_L$ of $L$ fits the
level of $x$ (since types, places and central actions agree), and by
\lref{GlobalTheory.T.exists_transported_state_heckeGaloisStage} there is a state $x'$ over $P$ whose
over-datum is that of $x$ with $N$ replaced by $N_L$ and whose row representation is that of $x$. By
\lref{GlobalTheory.HeightCounting.omFrame_exists_with_jE} there is a frame for $\Omega$ compatible with
$j_E$, and by \lref{GlobalTheory.HeightCounting.PhtExt.exists_phtOn_all} ($\mathcal{C}_{\mathrm{HT}}$) the
Weil--Deligne side of $A$ can be replaced by one satisfying the Hodge--Tate condition on this frame. The
selection map $P_q$ exists because the selection of $w$ contains a prime above every prime of the maximal
real subfield. Collecting the data
gives a same-Galois level of $x'$, and $x'$ has the same Galois data as $x$ by construction (equality
of the localisation data, of the field $\Omega$ and of the Galois data recorded in (iii)).
\end{proof}

\subsubsection{Rationality}

\begin{lemma}\label{pf:baselat2:isot}
The Weil representation of unitary dual pairs, the existence of discrete series, $\mathcal{H}_{\mathrm{mc}}$
and $\mathcal{I}_3$ imply that at every state $x$ the isotypic decomposition of the
rows of $D_x$ is internal: for $1\le j\le m$ the isotypic parts of the $j$-th row form an internal direct
sum.
\formallabel{GlobalTheory.HeightCounting.Up.IsotAtStates.isotInternal_at_states}
\end{lemma}

\begin{proof}
By \lref{GlobalTheory.HeightCounting.Up.IsotAtStates.isotInternal_of_multEq} (transported to states by
\lref{GlobalTheory.HeightCounting.Up.RowModel.overRows_allSt}), at every state, a rational reading
without congruence of the real datum together with the multiplicity equality of the real datum (equality of
its twisted multiplicities) implies internality; the rational reading is the one contained in the read package of the state.
The multiplicity equality holds at pinned witnesses by the multiplicity statement contained in
$\mathcal{H}_{\mathrm{mc}}$, and starting witnesses are pinned. The conditions required there (group
conditions, existence at all places and the geometric condition) are part of the state.
\end{proof}

\begin{proposition}\label{pf:baselat2:rat}
$\mathcal{I}_3$, $\mathcal{I}_{\mathrm{ser}}$ and internality of the isotypic decompositions at
all states (\cref{pf:baselat2:isot}) imply $\bltwoS{rat}$ for every $h_L$. (The second formal statement below
also assumes the uniqueness of discrete series of $\Uni(p,q)$; this hypothesis is not used.)
\formallabel{GlobalTheory.HeightCounting.RatSide.sRat_of_isotInternal_alt_alt}
\formallabel{GlobalTheory.HeightCounting.RatSide.sRat_of_isotInternal_inputs}
\end{proposition}

\begin{proof}
For a state $x$ with a same-Galois level, a rationality datum consists of an embedding
$k_{K_s}:K_s\to L_0$ compatible with the inclusion into $L$, a rational input and a rational adaptation. The following statements, each of the form ``at every state with a same-Galois
level'', provide its components:
\lref{GlobalTheory.HeightCounting.RatSide.formsC_at_state_alt} (for $1\le j\le m$ the forms of the rational datum $S_x$ are $K_s$-forms, finite-dimensional over $K_s$, and
their integral lattice is spanned over $\OO_{K_s}$ by a finite set which also spans the forms over $K_s$);
\lref{GlobalTheory.HeightCounting.ThetaSt.main_alt_alt} (there are characters
$\theta_k:\Z/g_k\Z\to L_0^\times$, $\theta_k\neq1$, satisfying the level condition with respect to the real datum);
\lref{GlobalTheory.HeightCounting.RatSide.HeckeAtState.heckeAtState_alt} (rapid decay and two Hecke
inputs of $\mathcal{I}_{\mathrm{ser}}$ give Hecke operators $h_L'(j,k,\mathfrak P)$ on the forms which stabilise them, satisfy
$h'_L(j,0,\cdot)=1$ and are read by the real datum);
\lref{GlobalTheory.T.intC_state_of_lvlPsi_stabC_heckeRd} (integrality of these operators);
\cref{pf:baselat2:eigspan} (eigenvalue data and the spanning property); and
\lref{GlobalTheory.Mem.RatDatRExp.rhoStab_chiInLL0} (stability of $\rho$ and values of $\chi$ in $L_0$
for the rational datum of $x$). Assembling these components gives a rationality datum pointwise.
\end{proof}

\begin{lemma}\label{pf:baselat2:eigspan}
Under $\mathcal{I}_3$ and internality of isotypic decompositions at all states, eventually
for large twisting data, for every state $x$ with a same-Galois level and all $\theta,h'_L$ as in the
proof of \cref{pf:baselat2:rat} (with $\theta_k\ne1$, the level condition, stability, $h'_L(j,0,\cdot)=1$ and
the reading property) there is an eigenvalue system $\mathrm{eig}$ with values in $L_0$ which is adapted to
the read package of $x$ and such that the forms are spanned by the corresponding eigenvectors.
\formallabel{GlobalTheory.HeightCounting.RatSide.eigenspace_span_of_isotypicInternal_alt}
\end{lemma}

\begin{proof}
\Lref{GlobalTheory.HeightCounting.RatSide.exists_eigAd_of_state_alt} gives an adapted eigenvalue system
from stability and the reading property; \lref{GlobalTheory.T.span_c_of_stab_c_hecke_rd_eig_ad} gives the
spanning property for every adapted system, provided the isotypic decomposition is internal and an
occurrence condition on good types holds; and this condition follows from the goodness of all types
(\lref{GlobalTheory.T.allSt_goodFTAll_of_levelA} and \lref{GlobalTheory.EigRows.occGood_of_all}).
\end{proof}

\subsubsection{Row cuspidality}

\begin{proposition}\label{pf:baselat2:cusp}
$\mathcal{I}_3$, $\mathrm{Par}(h_L)$ and $\bltwoS{rat}$ imply $\bltwoS{cusp}$.
\formallabel{GlobalTheory.SCuspRows.sCusp_of_parSt_alt_alt}
\end{proposition}

\begin{proof}
Let $x$ be a state with a same-Galois level $l$. Row cuspidality asks, for $1\le j\le m$ and every type
$\sigma$ with nonzero isotypic part in the $j$-th row and good row, for a cuspidal representation $\pi$ of
$\GL_r$ over $E$ which is conjugate self-dual and regular algebraic, whose local components at split
places are the local representations of the row, and such that for all but finitely many places $w$ of
$E$ the characteristic polynomial $P_w$ of the local parameter of
$(\pi\otimes\eta_T^{-1})_w$, twisted by $|\cdot|^{(1-n_V)/2}$, has coefficients in $L_0$. The parameter
statement gives $\pi$ with the first three properties. The rationality datum of $l$ gives, by
\lref{GlobalTheory.SCuspRows.exists_eigTrue_alt}, a finite set $S_b$ of places and an eigenvalue system
which is true for the row. Then:
\begin{enumerate}[label=(\alph*)]
\item by \lref{GlobalTheory.CuspFrobSplit.galPol_lifts_ae}, for almost all split places $w$ the
  characteristic polynomial of the local parameter of $\pi_w$ twisted by $|\cdot|^{(1-r)/2}$ has coefficients in $L_0$ (the
  eigenvalues at split places lie in $L_0$, and the valuation datum of
  \lref{GlobalTheory.HeightCounting.Up.lowerLeftValuation_eq_zero_all_stages_alt} controls the
  normalisation);
\item by \lref{GlobalTheory.GalAllPlaces.galPol_lpTw_mem_lifts_eventually} (using that the
  Galois side of $l$ is true, finiteness of automorphic forms on $\GL_n$, a published statement on
  cuspidal representations of unitary groups and the
  Chebotarev density theorem for primes of degree one), the same holds for almost all places;
\item by \lref{GlobalTheory.TwistLift.galPol_lifts_tw_etaTinv}, it then holds for the twist of $\pi$ by
  $\eta_T^{-1}$ with the normalisation $(1-n_V)/2$.\qedhere
\end{enumerate}
\end{proof}

\subsubsection{The eta statement and finiteness}

\begin{proposition}\label{pf:baselat2:eta}
$\mathcal{I}_3$ and $\mathrm{Pt}(h_L)$ imply $\mathbf{S}_{\eta}$.
\formallabel{GlobalTheory.StEta.sEta_alt_alt_alt}
\end{proposition}

\begin{proof}
We use the notions of \cref{sec:pf:baselat2:ingr} (existence and forcing of the automorphic side, the
conditions $(\mathrm{O})$, $(\mathrm{P}_{41})$, $(\mathrm{GR})$ and largeness). Let $x$ be a state over $P$ with a
strict-Henselian-acyclic level $l$ for which the automorphic side exists. The following statements hold at
$l$: \cref{pf:baselat2:etaj} (the automorphic side is forced by the property ``$(\mathrm{O})$ implies
$D'.\eta_J(j,v)=P.q.\eta(j,v)$'');
\lref{GlobalTheory.T.star_mem_LL0_at_level_c} ($L_0$ is stable under complex conjugation);
\lref{GlobalTheory.T.eExp_integer_of_levelC} (the exponents are integers);
\lref{GlobalTheory.T.places_over_pde_isSplitPl_of_levelC} (the places over the auxiliary primes
$p_{de},p'_{de}$ are split);
\lref{GlobalTheory.T.chebSplit_of_published_inputs_at_level_c} (the Chebotarev splitting statement);
\lref{GlobalTheory.T.eventually_galoisRepsCusp_of_level_c} (the automorphic side is forced by
$(\mathrm{GR})$). The point statement $\mathrm{Pt}(h_L)$ and the last statement give, by the laws of forcing, that
the automorphic side is forced by the property ``$(\mathrm{O})$ implies $(\mathrm{P}_{41})$ and residual
irreducibility, and $(\mathrm{GR})$ holds''. With these, \cref{pf:baselat2:large} shows that the automorphic side
is forced by largeness. The laws of forcing combine this with the first statement; since largeness contains
$(\mathrm{O})$, the eta equality holds for the forcing data.
\end{proof}

\begin{lemma}\label{pf:baselat2:sxd}
$\mathbf{S}_{\theta}$ holds for every $h_L$.
\formallabel{GlobalTheory.HeightCounting.Up.sXd_all_alt}
\end{lemma}

\begin{proof}
Let $x$ be a state with a second intermediate level $l$. The series data of $l$ contain an export tie;
by \lref{GlobalTheory.ExportType.type_of_exportTie} (using that the real datum of $x$ is $L$-uniform
and models highest weights, both part of the state), the type maps of the export tie satisfy the kernel
relations required by theta data. The theta datum is then the one which, for every $(i,j,b)$, selects
all data if $1\le j\le m$ and none otherwise; it is full by construction.
\end{proof}

\begin{proposition}\label{pf:baselat2:fin}
The Weil representation of unitary dual pairs, the existence of discrete series, $\mathcal{I}_3$,
$\mathcal{C}_{\infty}(h_L)$, $\mathbf{S}_{\eta}$ and $\mathbf{S}_{\theta}$ imply $\bltwoS{fin}$ and, for every
$c$, $\bltwoS{fin}^{c}$.
\formallabel{GlobalTheory.HeightCounting.Fin3.sFin_and_sFinl_of_theta_data_exists}
\formallabel{GlobalTheory.T.sFin_and_sFinl_of_archCmp_sEta_thetaData}
\end{proposition}

\begin{proof}
By \lref{GlobalTheory.HeightCounting.Fin3.start_at_all_pre2} the rows are available at every
pre-datum over a starting witness. Both statements then follow pointwise from
\cref{pf:baselat2:lv3third}, applied with $N=\bltwoN{law}$ (the hypothesis on $N$ is trivial) and with
$N=\bltwoN{list}(c)$ (an $\bltwoN{law}$-datum together with the list pin, which is transported from $x$ to
the level by the hypothesis of \cref{pf:baselat2:lv3third}, is an $\bltwoN{list}(c)$-datum).
\end{proof}

\begin{lemma}\label{pf:baselat2:lv3third}
Assume the hypotheses of \cref{pf:baselat2:fin}. Let $x$ be a state over $P$ with a strict-Henselian-acyclic
level $l$ satisfying the conclusion of $\mathbf{S}_{\eta}$, let $x_d$ be full theta data for $D_x$,
and let $N$ be a norm-choice extension such that every $\bltwoN{law}$-datum at a level $Z$ with
$\bltwoWt=P.\bltwoWt$ for which the list pins of $x$ transport to $Z$ yields an $N$-datum. Then there are a
level $Z$ of $w$ and $v\in\mathrm{Lv}_3(\bltwop(N,S,\Gamma).s,K,w,Z)$ with $v.\bltwoWt=P.\bltwoWt$ and
$l.c.\mathrm{ch}\approx v.\mathrm{nch}$.
\formallabel{GlobalTheory.HeightCounting.Fin3.lv3_of_third_alt_alt}
\end{lemma}

\begin{proof}
By \cref{pf:baselat2:builtex} the automorphic side exists for $l$; by $\mathbf{S}_{\eta}$ it is forced by large data,
so (by the laws of forcing) there is a local datum $\mathcal L$ built for the relevant family $\mathrm{big}$, i.e.\
with a specialisation $\mathrm{Sp}$, an automorphic specialisation and semisimple residual representations, for
which the eta equality holds. By \lref{GlobalTheory.HeightCounting.Fin3.nonempty_allL_of_archCmp}
the archimedean comparison provides the archimedean local data. By
\lref{GlobalTheory.HeightCounting.Fin3.exists_base_level_rXs_and_yRel} there are a base-level row side
$p$, rows $t$, an element of the row requirement and the relations between the readings at the level
$Z=(o_{\mathcal L}).Z$ attached to $\mathcal L$ (using the Weil representation and discrete series). By
\lref{GlobalTheory.HeightCounting.Fin3.exists_core_and_listPin} there are a norm-choice core and the
transport of the list pins. The hypothesis on $N$ gives an $N$-datum. The norm choice is the local
character choice of $\mathcal L$, which is kept from
$l.c.\mathrm{ch}$ by \lref{GlobalTheory.HeightCounting.Fin3.kept_characterChoiceLocal}; the condition on
$\chi_T$ in \cref{pf:baselat2:lv3def} holds because the types of the level are true
(\lref{GlobalTheory.Mem.Top.TypesTrue.chiT_gen}: true types give a generator $u$ of $(\Z/g_k\Z)^\times$ with
$\chi_T(k,v,j)=u^{2(j+1)}$). These data form an element of $\mathrm{Lv}_3$ with $\bltwoWt=P.\bltwoWt$.
\end{proof}

\subsubsection{Settings and forms, series, moduli points}

\begin{proposition}\label{pf:baselat2:settings}
$\mathcal{I}_3$ and $\mathcal{I}_{\mathrm{ser}}$ imply $\bltwoS{1}^{c}$ for all $h_L$ and $c$.
\formallabel{GlobalTheory.HeightCounting.Up.PointwiseSettings.settings_forms_exist_list_of_published_inputs}
\formallabel{GlobalTheory.T.settingsFormsExistList_of_inputs_serN}
\end{proposition}

\begin{proof}
By \lref{GlobalTheory.T.exists_lam_grid_fieldHas_at_member}, at every state $x$ there is a grid datum
$G$ for the level characters of $x$ at level $(1,c,b)$ whose values lie in $L_0$, whose associated
$\lambda$-datum is admissible, has character frame that of $x$ and has $q_0$ equal to the one of the
absolutely irreducible span. Given an intermediate level $l$ of $x$, \cref{pf:baselat2:pointwise} turns
$(x,l,G)$ into a state with a first intermediate level and the list pin.
\end{proof}

\begin{lemma}\label{pf:baselat2:pointwise}
Under $\mathcal{I}_3$ and $\mathcal{I}_{\mathrm{ser}}$, for every state $x$ over $P$ with an
intermediate level $l$ and every grid datum $G$ as in the proof of \cref{pf:baselat2:settings} there is a
state $x'$ over $P$ with a first intermediate level and $\mathrm{List}_c(x')$.
\formallabel{GlobalTheory.HeightCounting.Up.PointwiseSettings.pointwise_alt_alt}
\end{lemma}

\begin{proof}
The Galois side of $l$ gives Galois representations
(\lref{GlobalTheory.HeightCounting.galEx_of_cgTrue_alt}: a true Galois side yields Galois
representations for $D_x$), and Chebotarev gives the splitting statement
for $D_x$. \Cref{pf:baselat2:secondstep} gives a size bound $c_0$ such that for all size bounds $b_0\ge c_0$, all
finite sets $S_x$ of places containing the set $S_b$ of the rational input, closed under conjugation and
containing the places over the low primes, and all ring specifications, there are a twisting frame $F$, an
expansion datum $\mathrm{ex}$ whose sup lattices are finitely generated and span the forms over the fraction
field, and side data. The twist rings, characters and class groups
are provided by \lref{GlobalTheory.T.exists_twistRings_characters_classGroups} (for the maximum of the two
bounds); the set $S_x$ by \lref{GlobalTheory.HeightCounting.Up.PointwiseSettings.exists_Sx}. The closures
$\mathrm{cl}_C$ of the cyclotomic integers are closed under quotients of elements of non-increasing size
(\lref{GlobalTheory.CycClosure.exists_cyclotomic_integers}), which provides the generators $i_C$ by
\lref{GlobalTheory.Mem.Df.exists_iC_of_clC}. The toric characters are placed in the twist coefficient
rings by \lref{GlobalTheory.T.exists_toricCharacters_in_twistCoefficientRings} (using Chebotarev, Artin
reciprocity, the ratio statement of $\mathcal{I}_{\mathrm{ser}}$, the grid datum $G$ and the fact that the conductor
primes lie in the low set, \lref{GlobalTheory.T.conductor_primes_in_low_set}); this gives a $\lambda$-tie and
side data. Theta data for the new twisting are produced by
\lref{GlobalTheory.HeightCounting.exists_nth_of_fg_full} from finite generation and full rank. Finally
\lref{GlobalTheory.T.side_data_over_set_levels_alt} turns these side data, the grid datum and the
$\lambda$-tie into a state $x'$ with a first intermediate level and $\mathrm{List}_c(x')$.
\end{proof}

\begin{proposition}\label{pf:baselat2:series}
$\mathcal{I}_3$ and $\mathcal{I}_{\mathrm{ser}}$ imply $\bltwoS{2}^{c}$ for all $h_L$ and $c$.
\formallabel{GlobalTheory.HeightCounting.Up.SeriesStep.sb1l_alt_alt}
\formallabel{GlobalTheory.T.seriesExistList_of_inputs_serN}
\end{proposition}

\begin{proof}
Let $x$ be a state over $P$ with a first intermediate level $l$, at a witness over $(K,B,f,T)$ with $K$ over
the pre-base and $T$ forced. By \cref{pf:baselat2:sercopy} there is a new twisting $N$ of $D_x$ which preserves
full $\lambda$-ties and such that series data exist at every state $x'$ with an intermediate level whose
over-datum is the twist $o_{\mathrm{tw}}(x,N)$ of that of $x$ and whose remaining data are the transports of those
of $l$ (for module inputs with identity frames). Such a state $x'$ with an intermediate level
$l'$ exists by \lref{GlobalTheory.T.twisted_state_transport_heq_levels'}; replacing the module
input of $l'$ by the one supplied at every state by
\lref{GlobalTheory.MtId.exists_matrix_input_identity_frame_of_inputs} gives a second intermediate level
of $x'$. The list condition is preserved by \lref{GlobalTheory.HeightCounting.Up.SeriesStep.listSt_oTw},
since $N$ preserves full $\lambda$-ties.
\end{proof}

\begin{proposition}\label{pf:baselat2:points}
$\mathcal{I}_3$ and $\mathcal{I}_{\mathrm{ser}}$ imply $\bltwoS{3}^{c}$ for all $h_L$ and $c$.
\formallabel{GlobalTheory.HeightCounting.Up.points_step_alt_alt}
\formallabel{GlobalTheory.T.moduliPointsExistList_of_inputs_serN}
\end{proposition}

\begin{proof}
By \cref{pf:baselat2:geolaws}, for every state $x$ with a second intermediate level $l$ and every twisted
index $i$, the series data of $l$ satisfy the geometric laws at the generator $i_C$. These laws provide theta
data $\mathrm{th}$ with a theta tie, and
\lref{GlobalTheory.T.exists_levelB2_twist_of_thetaTie} gives a state $x'$ over $P$ with a third
intermediate level whose over-datum is the twist of that of $x$ by $\mathrm{th}$ and whose level characters
are the transported ones. The list pin transports: the grid datum of $x$ is moved to the new level by
the grid transport of \lref{GlobalTheory.Mem.GridMove.gdC}, with the same values and roots, and the full
$\lambda$-tie is preserved.
\end{proof}

\subsubsection{The strict-Henselian-acyclic level}

\begin{definition}\label{pf:baselat2:thirdstage}
Let $h_L\in\sfname{RecTExists}$. The statement $\bltwoS{ch}$ is $\forall^{\mathrm{st}}(\varphi)$, where $\varphi(x)$ says: for
every third intermediate level $l$ of $x$ and every map $s$ from the places of $L$ to the twisted indices of
$D_x$ which reads every place and is invariant under conjugation of places, there are height-counting
character data $c$ for $l$ whose norm choice has twisted index map $s$.
\formallabel{GlobalTheory.HeightCounting.Up.ThirdStageExists}
\end{definition}

\begin{lemma}\label{pf:baselat2:sck}
$\bltwoS{ch}$ implies $\bltwoS{sha}$.
\formallabel{GlobalTheory.HeightCounting.Up.ChoiceEx.sc_of_sck_alt}
\end{lemma}

\begin{proof}
Pointwise: let $l$ be a third intermediate level of a state $x$. Since $L_0$ is stable under complex
conjugation (elementary datum of $l$) and the series data of $l$ cover every place (for every place $v$ there is
a twisted index $i$ with $\ell_i$ the residue characteristic of $v$ at which $v$ is a side place),
\lref{GlobalTheory.HeightCounting.Up.ChoiceEx.exists_settings_conj} gives a map $s$ which reads every place
and is conjugation invariant. Then $\bltwoS{ch}$ gives character data for $l$.
\end{proof}

\begin{proposition}\label{pf:baselat2:sha}
The discrete series cover, the Casselman--Wallach theorem, the theory of Harish-Chandra modules,
$\mathcal{I}_3$ and $\mathcal{I}_{\mathrm{ser}}$ imply $\bltwoS{sha}$ for every $h_L$.
\formallabel{GlobalTheory.HeightCounting.Up.ChoiceEx.stage_c_nonempty_of_published_inputs}
\formallabel{GlobalTheory.T.stage_c_nonempty_of_prints_third_layer}
\end{proposition}

\begin{proof}
\Cref{pf:baselat2:stagec,pf:baselat2:sck}.
\end{proof}

\begin{proposition}\label{pf:baselat2:stagec}
The discrete series cover, the Casselman--Wallach theorem, the theory of Harish-Chandra modules,
$\mathcal{I}_3$ and $\mathcal{I}_{\mathrm{ser}}$ imply $\bltwoS{ch}$ for every $h_L$.
\formallabel{GlobalTheory.LevCmp.stageCExists_of_harishChandra_serN}
\formallabel{GlobalTheory.T.stageCExists_of_casselmanWallach}
\end{proposition}

\begin{proof}
Eight statements of the form $\forall^{\mathrm{st}}(\cdot)$ are combined: \cref{pf:baselat2:choice} (character data
with prescribed twisted index exist for every norm choice satisfying the conditions listed there);
\lref{GlobalTheory.T.gPlaces_of_levelA}, \lref{GlobalTheory.T.twistConductors_avoid_excludedPrimes},
\lref{GlobalTheory.T.badV_subset_gPl_union_qPl_of_admTw},
\lref{GlobalTheory.T.paraShape_of_twisted_sideRead_readsHecke},
\lref{GlobalTheory.T.level_uniformity_of_cmDatum_alt}, \lref{GlobalTheory.T.joint_level_of_cm_datum_twOK}
and \lref{GlobalTheory.T.cellWt_of_sideRead_of_twOK} (local properties of admissible twists: the places
of type $g$, avoidance of excluded primes by the twist conductors, location of the bad places, the
parahoric shape, level uniformity, the joint level and the cell weights). Given a third intermediate level
$l$ of a state $x$ over $P$ and a map $s$ as in $\bltwoS{ch}$, the level-comparison choice
\lref{GlobalTheory.LevCmp.exists_level_comparison_choice_alt} produces a norm choice $C$ with twisted
index $s$, a tied tuple and a choice side, using the twisting datum $T$ (its characters $\psi_v$ satisfy
the condition on $(\ell_v,Q_v)$ required of twisting characters and have order dividing $\ell_v^{k_v}$, $L$
contains the corresponding roots of unity,
$Q_v=\emptyset$ at untwisted places, and the primes of $Q_v$ avoid the excluded places, since they lie outside the exceptional sets of the
conjugate places); its admissibility at every twisted place
is the comparison of levels (\cref{pf:baselat2:leveq}) together with the local statements above and the
nontriviality of local unit characters (\lref{GlobalTheory.T.exists_local_unit_character_nontrivial}). The
first statement then gives the character data.
\end{proof}

% ---------------------------------------------------------------------------
\subsection{Spectral finiteness and the second pass}\label{sec:pf:baselat2:second}

\begin{lemma}[Spectral finiteness]\label{pf:baselat2:sfinoalt}
For all $h_L$ and $c$, $\bltwoS{fin}^{c}$ implies $\bltwoS{spec}^{c}$.
\formallabel{GlobalTheory.HeightCounting.Up.sFino_of_alt}
\end{lemma}

\begin{proof}
Pointwise: let $x$ be a state over $P$ with a strict-Henselian-acyclic level $l$, let $(S,\Gamma)$ be given,
and assume $\mathrm{List}_c(x)\wedge\mathrm{Sat}(P)$. $\bltwoS{fin}^{c}$ gives a level $Z$ and
$v\in\mathrm{Lv}_3(\bltwop(\bltwoN{list}(c),S,\Gamma).s,K,w,Z)$ with $v.\bltwoWt=P.\bltwoWt$ and
$l.c.\mathrm{ch}\approx v.\mathrm{nch}$. Keep all components of $v$ and extend its norm-requirement
component, an element of $\bltwoN{list}(c)$, by the origin datum $(P,x,l)$: the conditions of
$\mathrm{Orig}_c$ are exactly $\mathrm{List}_c(x)\wedge\mathrm{Sat}(P)$, $v.\bltwoWt=P.\bltwoWt$ and
$l.c.\mathrm{ch}\approx v.\mathrm{nch}$. This is an element of
$\mathrm{Lv}_3(\bltwop(\bltwoN{or}(c),S,\Gamma).s,K,w,Z)$ with the same $\bltwoWt$ and norm choice.
\end{proof}

The second pass is obtained from the following criterion for $\mathrm{Pass}$.

\begin{lemma}[A criterion for passes]\label{pf:baselat2:passof}
Let $Q$ be a slot and $p,q,r$ predicates on $(K,B,f,T)$. Assume $\mathrm{Step}(\kappa_0,\top,p)$,
$\mathrm{Step}(\kappa_0,\top,q)$ and the eventual statement: \emph{eventually in the scale, for all
$B,f,T$ over $K$ with $\kappa_0(B,f)$, $T$ untwisted and $p(K,B,f,T)$, for every core $c$ satisfying
$Q$ whose assembled datum satisfies the export condition, for every finite exclusion function $\mathcal A$ and
every finite set $S_2$ of algebraic numbers (depending on the pre-base, the frame and a set of
characters), there are $B',f',T'$ with $B'$ enlarging $B$, $f'=f$, $T'.\mathrm{pl}=T.\mathrm{pl}$, $T'$
avoiding $\mathcal A$, $S_2$ realised in $B'.\mathrm{LL}$, $\kappa_0(B',f')$, and
$q(K,B',f',T')\Rightarrow r(K,B',f',T')$.} Then $\mathrm{Pass}(\kappa_0,Q,r)$.
\formallabel{GlobalTheory.HeightCounting.Up.pass_of}
\end{lemma}

\begin{proof}
Eventual statements in the scale are closed under finite conjunctions
(\lref{GlobalTheory.Mem.Top.EvG.inter}); fix a scale context $K$ at which the three statements hold.
Since forcedness is itself large (\lref{GlobalTheory.Mem.Top.GroupRetwisting.lg_forced}: for every $K$,
the prime sets of large twisting data contain the primes required by $\mathrm{Forced}$), the statement for
$q$ gives $\mathrm{lg}_2(K,\ \kappa_0\Rightarrow q)$, i.e.\ exceptional data $S_1,P_1,\mathcal A_0,S_{2,0}$
such that $q(K,B,f,T)$ holds whenever $\kappa_0(B,f)$, the fixed coefficient field $L_0$ of $B$ contains
$S_1$, $B.\mathrm{LL}$ contains $S_{2,0}$, $P_1\subseteq T.\mathrm{pl}$ and $T$ avoids $\mathcal A_0$. The
property ``$S_1\subseteq L_0$ and $P_1\subseteq T.\mathrm{pl}$'' is again large (take the same $S_1,P_1$, no
excluded primes, and the same $S_{2,0}$), and so is $p$; we show that the conclusion of $\mathrm{Pass}$ holds
for all $(B,f,T)$ in the intersection of these two large ranges.

Let $(B,f,T)$ be such, with $\kappa_0(B,f)$ and $T$ untwisted, let $c$ be a core with $Q$ and
the export condition, and let $\mathcal A$, $S_2$ be given. Apply the third hypothesis with the enlarged
exclusion function $\ell\mapsto\mathcal A(\ell)\cup\mathcal A_0(\ell)$ and the enlarged set
$S_2\cup S_{2,0}$ (both are again admissible); it gives $(B',f',T')$. Since $B'$ enlarges $B$ it has the
same pre-base, the same subfield $K_s$ and the same fixed coefficient field $L_0$, and $f'=f$; the
exceptional data $S_1,P_1,\mathcal A_0,S_{2,0}$ depend only on these, so they are the same for $(B',f')$ as for
$(B,f)$. Hence $S_1\subseteq L_0$, $S_{2,0}$ is realised in $B'.\mathrm{LL}$, $P_1\subseteq T.\mathrm{pl}=T'.\mathrm{pl}$
and $T'$ avoids $\mathcal A_0$; with $\kappa_0(B',f')$ this gives $q(K,B',f',T')$, hence
$r(K,B',f',T')$. Finally $T'$ avoids $\mathcal A$ and $S_2$ is realised in $B'.\mathrm{LL}$, as required by
$\mathrm{Pass}$.
\end{proof}

Built cores of the origin parameter remember a state. Recall that a \emph{seed} for a datum $D$ and a
family $A=(A_i)$ of finite sets of primes of $E$ indexed by the twisted indices $i$ of $D$
consists of: a map $s$ from the places of the coefficient
field to twisted indices which reads every place, is invariant under conjugation and is constant on
twisted places lying over one place; full theta data; for each twisted index $i$ a finite
set $Q_i$ of primes of $E$, an exponent $k_i$ and a character $\psi_i$ of the ray class group of $E$ of
modulus $Q_i$ (at the index $i$) with values in the units of the coefficient ring $\OO_i/\ell_i^{k_i}$;
such that at twisted places the twist inputs are admissible and all twists $\psi_i^a$, $a$ prime to
$\ell_i$, are good twists avoiding $A_i$; and such that the sets $Q_i$ for distinct indices at twisted
places are disjoint.

\begin{lemma}[States of built cores]\label{pf:baselat2:prepoint}
Let $c$ be a core at a level $Z$ of $w$ which is built for $\bltwop(\bltwoN{or}(c_0),S,\Gamma).s$ and
whose datum satisfies the export condition. Then there are a pre-datum $P$ over $(K,w)$, a state $x$ over $P$
and a strict-Henselian-acyclic level $l$ of $x$ with $\mathrm{List}_{c_0}(x)\wedge\mathrm{Sat}(P)$, and for
every family $A$ there is a seed for $(D_x,A)$.
\formallabel{GlobalTheory.HeightCounting.Up.exists_pre_point_with_henselian_level}
\end{lemma}

\begin{proof}
The norm requirement of a built core for $\bltwoN{or}(c_0)$ contains an element of
$\mathrm{Orig}_c$: a pre-datum $P$, a state $x$ over $P$ and a level $l\in\bltwoL{sha}(x)$ with
$\mathrm{List}_{c_0}(x)\wedge\mathrm{Sat}(P)$, $c.\bltwoWt=P.\bltwoWt$ and $l.c.\mathrm{ch}\approx c.\mathrm{nch}$.
A seed is then provided by \lref{GlobalTheory.HeightCounting.Up.nonempty_seed_of_kept}: for a norm
choice kept from the character data of $l$, with the same twisted places (which holds because, by
\lref{GlobalTheory.Mem.Top.TupTieW.tupTied_of_muIsW}, the twisted places of the core and of $x$ are those
of $c.\bltwoWt=P.\bltwoWt$), with a conjugation-invariant twisted index map from the side of $l$, with the
twisted place indices in the first pair
(\lref{GlobalTheory.Mem.Top.twistedPlaceIndex_mem_firstPair_of_compatibleChoice}), the theta data of the
core and the export condition, seeds exist for every $A$.
\end{proof}

New twisting data are produced from seeds by global class field theory.

\begin{lemma}[Twisting characters from seeds]\label{pf:baselat2:twlift}
Assume $\mathcal{I}_3$. Let $X$ be stages over a base $B$ with datum $D$, $f$ a
frame, $\mathrm{pl}$ a finite set of primes, $A$ a family of finite sets of primes of $E$, $\sigma$ a seed for
$(D,A)$ and $\mathcal A$ an exclusion function. Assume: the residue characteristic of every twisted place of
$D$ lies in $\mathrm{pl}$; a place lying over one place with a twisted place is twisted; and for every
twisted index $i$
and every prime $p$ which lies in the exceptional set $\mathcal X(B)$ of the pre-base or in $\mathcal A(\ell_i)$, all
primes of $E$ above $p$ lie in $A_i$. Then for every finite set $S_2$ of algebraic numbers there are an
enlargement $B'$ of $B$, a frame $f'=f$ and a twisting datum $T'$ over $B'$ with $T'.\mathrm{pl}=\mathrm{pl}$,
$T'$ avoiding $\mathcal A$, $S_2$ realised in $B'.\mathrm{LL}$, $\mathrm{coh}(B',T')$, and $T'$ \emph{twisted from
$\sigma$}: $B'.\mathrm{LL}$ contains the $\ell_{v'}^{k_{v'}}$-th roots of unity for every place $v'$; at a
place $v'$ whose restriction $v$ to $B$ is twisted, $Q_{v'}=Q_{s(v)}$, $k_{v'}=k_{s(v)}$ and $\psi_{v'}$ is a
Hecke character which corresponds to $\psi_{s(v)}$ under Artin reciprocity at Frobenius elements;
and $\psi_{v'}=1$ at the other places.
\formallabel{GlobalTheory.HeightCounting.Up.exists_twisted_cohomology_class_lift_alt}
\end{lemma}

\begin{proof}
The published inputs contain global Artin reciprocity: an Artin map
$\mathrm{art}:\A_E^\times\to\Gal(\Omega/E)^{\mathrm{ab}}$ with the reciprocity properties. For each index $i$
with admissible seed data, compose $\psi_i$ (a character of the ray class group, i.e.\ of the Galois group
of the ray class field) with a root isomorphism $\iota_i$ identifying $(\OO_i/\ell_i^{k_i})^\times$-valued
roots of unity with complex ones, and pull back along $\mathrm{art}$; this is a Hecke character $\chi_i$
of $E$ (for non-admissible $i$ put $\chi_i=1$). By
\lref{GlobalTheory.Mem.Top.TwChar.pickChar_spec} it reads $\psi_i$ at Frobenius elements, satisfies
the twisting condition for $(\ell_i,Q_i)$, has order dividing $\ell_i^{k_i}$, and is a twist character for the
exceptional set (finite order, trivial on the id\`eles of $F$, tame, ramified only at split primes of prime
norm outside the exceptional set). By \lref{GlobalTheory.Mem.Top.TwChar.exists_twDat_stagesOn} there is a
twisting datum $T_1$ over $B$ with $T_1.\mathrm{pl}=\mathrm{pl}$, avoiding $\mathcal A$, equal to
$(\chi_{s(v)},Q_{s(v)},k_{s(v)})$ at twisted places $v$ and untwisted elsewhere. Avoidance holds because a
prime of $Q_i$ has prime absolute norm and does not lie in $A_i$
(\lref{GlobalTheory.Mem.Top.TwChar.pick_not_over}), whereas every prime of $E$ above a prime of
$\mathcal A(\ell_i)$ lies in $A_i$. By \lref{GlobalTheory.Mem.Top.TwChar.exists_enl_roots} there is an enlargement $B'$ of $B$ whose
coefficient field contains $S_2$ and the required roots of unity; put $T'=$ the pull-back of $T_1$ to $B'$
and $f'=$ the transport of $f$. Avoidance is preserved by pull-back, since the sets $Q$ are unchanged. Coherence: if two places $v'_1,v'_2$ of $B'$ have the
same restriction to $L_0$, then their restrictions to $B$ lie over one place
(\lref{GlobalTheory.HeightCounting.Up.overOne_plUnder}); either both restrictions are twisted, in which case
the seed gives the same index and hence the same $(Q,k,\psi)$, or, by the hypothesis on places lying over
one place,
neither is, in which case both data are untwisted. Roots of unity are present by construction.
\end{proof}

\begin{lemma}[Carrying level three to the enlarged base]\label{pf:baselat2:carried}
Assume $\mathcal{I}_3$. Let $B'$ enlarge $B$, $x$ a state over $P$ at $(B,f,T,w)$, $\sigma$
a seed for $(D_x,A)$ with $A$ containing the design sets of $D_x$, and $T'$ twisted from $\sigma$. Let $P'$ be a
pre-datum at a witness $w'$ over $(B',f',T')$ such that $v'\in P'.\bltwoWt$ if and only if the restriction of
$v'$ to $B$ lies in $P.\bltwoWt$, and let $y$ be a state over $P'$ with a third intermediate level $l$, such
that $D_x$ is carried to $D_y$, $\mathrm{List}_c(y)\wedge\mathrm{Sat}(P')$ holds, $\bltwoS{ch}$ holds at
$y$ (for every admissible map $s$ there are character data for $l$ with twisted index map $s$), and
$\bltwoS{spec}^{c}$ holds at $y$. Then $\mathrm{has}_3(\bltwop(\bltwoN{geo}(c),S,\Gamma).s)(w')$.
\formallabel{GlobalTheory.HeightCounting.Up.has3_of_carried_levelB2}
\end{lemma}

\begin{proof}
Carrying preserves the set of twisted indices (\lref{GlobalTheory.TwistChoice.same_basic}). The map
$v'\mapsto\sigma.s(v'|_B)$ reads every place of $B'$
(\lref{GlobalTheory.HeightCounting.Up.readsAt_of_plUnder}) and is conjugation invariant (restriction to $B$ commutes with
conjugation and $\sigma.s$ is conjugation invariant); $\bltwoS{ch}$ gives character data $c$ for $l$
with this index map, and $(l,c)$ is a strict-Henselian-acyclic level of $y$. $\bltwoS{spec}^{c}$ gives a
level $Z$ of $w'$ and $v\in\mathrm{Lv}_3(\bltwop(\bltwoN{or}(c),S,\Gamma).s,K,w',Z)$ with $v.\bltwoWt=P'.\bltwoWt$ and
$v.\mathrm{nch}$ kept from $c.\mathrm{ch}$. The goodness condition $\mathrm{Good}_O(w'.D(Z),v.\mathrm{nch},v.x_d)$ holds
by \lref{GlobalTheory.HeightCounting.Up.goodO_of_alt_alt}, from the seed (whose twists are good), the
fact that $T'$ is twisted from $\sigma$, the correspondence of the sets $\bltwoWt$, and the tied tuples of
$v$ (\lref{GlobalTheory.Mem.Top.TupTieW.tupTied_of_muIsW}). Adding it to the norm requirement of $v$ gives
an element of $\mathrm{Lv}_3(\bltwop(\bltwoN{geo}(c),S,\Gamma).s,K,w',Z)$.
\end{proof}

\begin{theorem}[Second pass]\label{pf:baselat2:secondno}
Assume the discrete series cover, the Casselman--Wallach theorem, the theory of Harish-Chandra modules,
$\mathcal{I}_3$, $\mathcal{I}_{\mathrm{ser}}$, $\mathbf{X}$, $\mathbf{W}$,
the existence of discrete series and $\mathcal{C}_{\mathrm{HT}}$. Let $h_L\in\sfname{RecTExists}$ and assume $\bltwoS{mt}$,
$\bltwoS{rat}$ and $\bltwoS{cusp}$. Then $\bltwoS{fin}^{c}$ implies $\mathbf{P}^{c}$, for every $c$.
\formallabel{GlobalTheory.HeightCounting.Up.secondNo_of_sMt_sRat_sCusp_sFinl}
\end{theorem}

\begin{proof}
Fix $(S,\Gamma)$ and write $\bltwop_{\mathrm{or}}=\bltwop(\bltwoN{or}(c),S,\Gamma)$,
$\bltwop_{\mathrm{geo}}=\bltwop(\bltwoN{geo}(c),S,\Gamma)$. We apply \cref{pf:baselat2:passof} with $p=\top$,
$Q=\mathrm{Built}(\bltwop_{\mathrm{or}}.s)$, $r=\exists w,\ \mathrm{start}(w)\wedge\mathrm{has}_3(\bltwop_{\mathrm{geo}}.s)(w)$
and $q=q_3\wedge q_4$, where $q_3,q_4$ are the conclusions of the following two statements, each of the
form $\mathcal{W}(\cdot,\cdot)$ and hence of the form $\mathrm{Step}(\kappa_0,\top,\cdot)$.

(H3) \lref{GlobalTheory.Mem.R3.Rw.Second.exists_ext_twisted_places} (from Green's theorem on cuspidal
representations and $\mathbf{X}$): eventually, for large twisting data, for every pinned witness $w'$
over $(B',f',T')$, every pinned witness $w$ over $(B,f,T)$ with $B'$ enlarging $B$, $f'=f$, the same column
pin data and $T'.\mathrm{pl}=T.\mathrm{pl}$, and every fine real datum $z$ at a level of $w$ (with the group,
$\kappa$-definition and base-change conditions, $\bltwoWt$ the set of places determined by $z$, and residue
characteristics in $T.\mathrm{pl}$) such that $\psi'_{v'}=1$ at every place $v'$ of $B'$ none of whose
restrictions lies in $\bltwoWt$, there are a level $Z'$ of $w'$ and a fine real datum $z'$ at $Z'$ with the
corresponding conditions, which \emph{extends} $z$ ($\mathrm{Ext}(z,z')$) and such that $v'$ lies in the set
of places of $z'$ if and only if its restriction lies in $\bltwoWt$.

(H4) the conjunction of the following statements $\mathcal{W}(\mathrm{start},\cdot)$:
\lref{GlobalTheory.T.constructionStartingChoices_of_fine_adelicData} (from the Weil representation, the
existence of discrete series, Harish-Chandra finiteness, rapid decay and $\mathbf{W}$:
every fine real datum $z$ at a level $Z$ with the above conditions, untwisted outside its places, gives a
pre-datum $P$ with $P.Z_0=Z$, $P.\bltwoWt$ the places of $z$ and $\mathrm{Ext}(z,P.z)$);
\lref{GlobalTheory.T.bareK_of_discreteSeriesExists} (from the published inputs and the existence of
discrete series: the statement $\bltwoS{bare}$, which provides, for a pre-datum $P'$ at $w'$ and stages over
the same pre-base and frame whose $N$-stage satisfies the group condition, an unlocalised state $x_0$ over
$P'$ with the same Galois data as these stages); \lref{GlobalTheory.T.stageAExists_of_inputs_phtCmp} ($\bltwoS{sg}'$:
unlocalised states have states with a same-Galois level and the same Galois data); $\bltwoS{mt}$,
$\bltwoS{rat}$, $\bltwoS{cusp}$; $\bltwoS{mv}$ (\cref{pf:baselat2:mvbase}: if $B'$
enlarges $B$, $w'$ and a starting witness $w$ have the same pin data, $x$ over $P$ at $w$ has a third
intermediate level, $P'$ at $w'$ satisfies $\mathrm{Ext}(P.z,P'.z)$ and $x'$ over $P'$ has an intermediate level and
the same Galois data as $x$, then there is a state $y$ over $P'$ with a third intermediate level, $D_x$ is
carried to $D_y$, and $\mathrm{List}_{c'}(x)\Rightarrow\mathrm{List}_{c'}(y)$ for all $c'$);
$\bltwoS{ch}$ (\cref{pf:baselat2:stagec}); and $\bltwoS{spec}^{c}$ (\cref{pf:baselat2:sfinoalt}).

It remains to verify the third hypothesis of \cref{pf:baselat2:passof}. By
\lref{GlobalTheory.Mem.Top.ReW.extend_pinned} the following holds eventually in the scale: for every
pinned witness $w$ over $(B,f,T)$ there are an exclusion function $\mathcal A_1$ and a finite set $S_1$ of
algebraic numbers such that for every enlargement $B'$ of $B$, $f'=f$ and $T'$ avoiding $\mathcal A_1$ with
$S_1$ realised in $B'.\mathrm{LL}$, there is a pinned witness $w'$ over $(B',f',T')$ with the same column pin,
selection, $J$-datum and field datum as $w$. Fix $K$ in this range, $(B,f,T)$ with $\kappa_0(B,f)$ and $T$
untwisted, a built core satisfying the export condition, and $\mathcal A$, $S_2$.
\begin{enumerate}[label=(\arabic*)]
\item By the definition of positions the core
  is built at a level $Z$ of a witness $w$, which is a starting witness: it is pinned and coherent by
  the built data, and $T$ is untwisted.
\item By \cref{pf:baselat2:prepoint} there are $P$, $x$, $l\in\bltwoL{sha}(x)$ with
  $\mathrm{List}_c(x)\wedge\mathrm{Sat}(P)$ and seeds for every $A$. Let $\mathcal A_1,S_1$ be given by
  \lref{GlobalTheory.Mem.Top.ReW.extend_pinned} for $w$. Choose a seed $\sigma$ for the family
  $A_i$ consisting of the design set of $D_x$ at $i$ and of all primes of $E$ above primes in
  $\mathcal X(B)\cup(\mathcal A\cup\mathcal A_1)(\ell_i)$.
\item Apply \cref{pf:baselat2:twlift} to the stages of $x$, the frame $f$, $\mathrm{pl}=T.\mathrm{pl}$, $\sigma$,
  $\mathcal A\cup\mathcal A_1$ and $S_2\cup S_1$. Its hypotheses hold: twisted places of $D_x$ lie in
  $P.\bltwoWt$ (by the definition of states) and have residue characteristic in
  $T.\mathrm{pl}$; places lying over one place with a twisted place are twisted, because $P$ is saturated
  and $T$ untwisted; and
  primes above $\mathcal X(B)\cup(\mathcal A\cup\mathcal A_1)(\ell_i)$ lie in $A_i$ by construction. We obtain $(B',f',T')$ with $T'.\mathrm{pl}=T.\mathrm{pl}$,
  $T'$ avoiding $\mathcal A\cup\mathcal A_1$, $S_2\cup S_1$ realised, $\mathrm{coh}(B',T')$, and $T'$ twisted from $\sigma$.
\item By the choice of $\mathcal A_1,S_1$ there is a pinned witness $w'$ over $(B',f',T')$ with the pin data of
  $w$. Its field datum shows $\kappa_0(B',f')$ (the pre-base is unchanged).
\item Assume $q_3\wedge q_4$ at $(K,B',f',T')$. Apply (H3) to $w'$, $w$ and the real datum $P.z$ of $P$; the
  support condition holds because $T'$ is twisted from $\sigma$ (so $\psi'_{v'}=1$ unless the restriction
  of $v'$ is twisted, hence in $P.\bltwoWt$). We get $Z'$ and $z'$ with $\mathrm{Ext}(P.z,z')$ and the
  correspondence of places. These data show that $w'$ is a starting witness: pinned, with good data
  ($\mathrm{good}$ is witnessed by $(Z',z')$), and coherent.
\item Apply (H4) at $w'$. The starting-choice statement gives $P'$ with $P'.Z_0=Z'$, $P'.\bltwoWt$ the places of
  $z'$ and $\mathrm{Ext}(z',P'.z)$; by transitivity of extensions
  (\lref{GlobalTheory.T.standardReading_transport_and_extension_trans}) $\mathrm{Ext}(P.z,P'.z)$. $\bltwoS{bare}$,
  applied to $P'$ and the stages of $x$, gives an unlocalised state $x_0$ over $P'$ with the Galois data of
  $x$; $\bltwoS{sg}'$ gives a
  state $x_1$ over $P'$ with a same-Galois level and the same Galois data; $\bltwoS{mt}$, $\bltwoS{rat}$ and
  $\bltwoS{cusp}$ complete it to an intermediate level. $\bltwoS{mv}$, applied to $x$ with its third
  intermediate level (contained in $l$) and to $x_1$, gives $y$ over $P'$ with a third intermediate level,
  $D_x$ carried to $D_y$ and $\mathrm{List}_c(y)$.
\item By the correspondence of places, $v'\in P'.\bltwoWt$ if and only if $v'|_B\in P.\bltwoWt$; hence
  $\mathrm{Sat}(P')$ follows from $\mathrm{Sat}(P)$ (\lref{GlobalTheory.HeightCounting.Up.overOne_plUnder_of}).
  \Cref{pf:baselat2:carried}, with $\bltwoS{ch}$ and $\bltwoS{spec}^{c}$ at $y$, gives
  $\mathrm{has}_3(\bltwop_{\mathrm{geo}}.s)(w')$.
\end{enumerate}
Thus $r(K,B',f',T')$ holds, and $T'$ avoids $\mathcal A$ and realises $S_2$.
\end{proof}

\begin{theorem}[Spectral finiteness and second pass]\label{pf:baselat2:sfino}
Under the hypotheses of \cref{pf:baselat2:secondno}, $\bltwoS{fin}^{c}$ implies $\bltwoS{spec}^{c}$ and
$\mathbf{P}^{c}$.
\formallabel{GlobalTheory.HeightCounting.Up.sFino_and_secondNo_of_rat_cusp_finl}
\formallabel{GlobalTheory.T.sFino_and_secondNo_of_sFinl'}
\end{theorem}

\begin{proof}
\Cref{pf:baselat2:sfinoalt,pf:baselat2:secondno}.
\end{proof}

% ---------------------------------------------------------------------------
\subsection{Level-three movement and fine adelic data}\label{sec:pf:baselat2:lv3}

\begin{proposition}[Level-three movement]\label{pf:baselat2:lv3mv}
For all $h_L$, $c$, $S$ and $\Gamma$, $\mathrm{Mv}_3$ holds for $\bltwop(\bltwoN{or}(c),S,\Gamma).s$ and
for $\bltwop(\bltwoN{geo}(c),S,\Gamma).s$.
\formallabel{GlobalTheory.HeightCounting.Up.lv3Mv_extensionWithOrigin_and_nEog}
\formallabel{GlobalTheory.T.lv3Mv_extensionWithOrigin_and_NEog}
\end{proposition}

\begin{proof}
Let $l\in\mathrm{Lv}_3(\bltwop(\bltwoN{or}(c),S,\Gamma).s,K,w,Z)$, and let $S',N,h$ be a new spectral stage, division
lattices and heights; write $Z^\sharp=\mathrm{up}(Z,S',N,h)$. By
\lref{GlobalTheory.T.adelicRowsDivision_transport_spectralStage_lattices} there are a base-level row side
$p'_0$ and rows $t'$ at $Z^\sharp$ such that the package of $(p'_0,t')$ is the update of that of
$(l.p_0,l.t)$, the rows are the same and geometric, and the row requirement transfers.
By \lref{GlobalTheory.T.exists_kept_choice_upper_lattice_alt} there are a norm choice $\mathrm{nch}'$ and theta
data $x'_d$ at $Z^\sharp$ such that $\mathrm{nch}'$ is kept from $l.\mathrm{nch}$, norm-choice data transfer, goodness
$\mathrm{Good}_O$ transfers, and norm-choice cores transfer. The list pin transfers along rows that are the
same (\lref{GlobalTheory.Mem.GridMove.listPin_rowsSame}), and the origin datum transfers along kept
norm choices (if $\mathrm{Orig}_c$ holds for
$\mathrm{nch}$ and $\mathrm{nch}'$ is kept from $\mathrm{nch}$, it holds for $\mathrm{nch}'$, since keeping is
transitive). The condition on $\chi_T$, the $G$-condition, coherence and $\bltwoWt$ are unchanged. This gives
$l'\in\mathrm{Lv}_3(\bltwop(\bltwoN{or}(c),S,\Gamma).s,K,w,Z^\sharp)$ with the updated package. For
$\bltwoN{geo}(c)$ the same construction applies, and the additional goodness component transfers by the
goodness transfer above.
\end{proof}

\begin{proposition}[Extension of fine adelic data]\label{pf:baselat2:datext}
The discrete series cover, the theory of Harish-Chandra modules, rapid decay of cusp forms, the Casselman--Wallach theorem,
the Weil representation of unitary dual pairs, the existence of discrete series and Harish-Chandra finiteness imply
$\mathbf{X}$: eventually, for large twisting data, for every pinned witness $w'$ over $(B',f',T')$, every
pinned witness $w$ over $(B,f,T)$ with $B'$ enlarging $B$, $f'=f$ and the same pin data, and every fine real
datum $z$ at a level of $w$ (with the group conditions and the intersection-number condition), every level
$Z'$ of $w'$ satisfying the group conditions and the volume tie carries a fine real datum $z'$, on the
set of places of residue characteristic in $T'.\mathrm{pl}$, extending $z$.
\formallabel{GlobalTheory.Mem.R3.Rw.DatExt}
\formallabel{GlobalTheory.Mem.R3.Rw.datExt_of_hypotheses}
\formallabel{GlobalTheory.T.fineAdelicDatum_extends_to_larger_field_of_hypotheses}
\end{proposition}

\begin{proof}
\Lref{GlobalTheory.T.adelicData_extend_largerFields_of_hypotheses} proves, from the same hypotheses, the
variant $\mathbf{X}_0$ of this statement, and $\mathbf{X}_0\Rightarrow\mathbf{X}$ is a restatement (the two statements differ only in the order of quantifiers).
\end{proof}

% ---------------------------------------------------------------------------
\subsection{Ingredients of the stage constructions}\label{sec:pf:baselat2:ingr}

This subsection proves the principal ingredients of the stage constructions of
\cref{sec:pf:baselat2:allstages} and of the second pass (\cref{sec:pf:baselat2:second}). Each of them is
reduced, by an argument given in full, to statements of the development which are cited with their
content.

\paragraph{Local data and forcing.} Let $D$ be a global datum with a character frame $cf$, an embedding
$j_E:E\to\Qbar$ and a character-frame embedding $A$, let $M$ be a module family, $C$ a norm choice and
$\mathrm{tp}$ a tied tuple for $C$. Let $\mathrm{big}$ be a family of finite subsets of $\C$ indexed by the
places of $L$. A \emph{local datum built for $\mathrm{big}$} is a family $\mathcal L$ of local data attached to
the layer of $D$ (a localisation datum $X$, local twists, automorphic data $Y$, base-change data, type data
$Z$ and a further datum $B$, together with the base-change, residual and module-parameter ties between
$X$, $Y$ and $Z$) for which there are a specialisation $\mathrm{Sp}$ for $(A,M,C,\mathrm{big})$ and $X$ and its
twists, and an automorphic specialisation for $(A,M,\mathrm{tp},\mathrm{Sp},Y)$ and the base-change data, such
that all residual representations $\bar r(\sigma,\lambda)$ of $Z$ are semisimple. We write $D_{\mathcal L}$ for the
datum $D$ localised along $\mathcal L$. The \emph{automorphic side exists} for $(A,M,C,\mathrm{tp})$ if for every
$\mathrm{big}$ there is a local datum built for $\mathrm{big}$; it is \emph{forced by} a property $\Phi$ of global
data if there is $\mathrm{big}$ such that $\Phi(D_{\mathcal L})$ holds for every local datum $\mathcal L$ built for
$\mathrm{big}$. By \lref{GlobalTheory.CoefDuty.LargeRing.evO_laws_alt}, forcing is closed under
conjunction and monotone (if $\Phi$ is forced and $\Phi\Rightarrow\Psi$ for all data, then $\Psi$ is forced).

\paragraph{Conditions on semisimple data.} For a datum $D'$ with local indices $i$ (coefficient rings
$\OO_i$ with residue characteristic $\ell_i$) we use the following conditions of the development.
\begin{itemize}
\item[$(\mathrm{O})$] Every Hecke--Galois representation $G_E\to\GL_{n_V}(\Qbar_{\ell_i})$ at a point of the theta
  locus of $i$ is conjugate to the image of one with values in $\GL_{n_V}(\OO_i)$.
\item[$(\mathrm{Spl})$] For every $i$, every Hecke--Galois representation with values in $\GL_{n_V}(\OO_i)$ is split
  after extension of scalars to the fraction field of $\OO_i$.
\item[$(\mathrm{Dec})$] For every $i$, every point $\pi$ of the theta locus and every Hecke--Galois representation
  $r$ at $\pi$ with values in $\GL_{n_V}(\OO_i)$ there is $\rho:G_E\to\GL_r(\OO_i)$ such that $r$ decomposes as
  prescribed by $(\rho,\chi_i)$.
\item[$(\mathrm{Tk})$] At every live index $i$ there are a good residual representation and a representation
  $\rho$ for which the wild-inertia condition holds (after enlargement to a finite discrete valuation ring over
  $\Z_{\ell_i}$ inside $\OO_i$ there are a Hida setting and, at every point of the theta locus, a wild-inertia
  cyclotomic element with a choice whose roots lie in $\OO_i$).
\item[$(\mathrm{GR})$] Every conjugate self-dual regular algebraic cuspidal representation of $\GL_a$ ($a\ge1$)
  appearing in $D'$ has an attached $\ell_i$-adic Galois representation at every local index $i$.
\end{itemize}
$D'$ is \emph{large} if it satisfies $(\mathrm{O})$, $(\mathrm{Spl})$ and $(\mathrm{Tk})$. We also use the provision
$(\mathrm{P}_{41})$, residual irreducibility (for good indices) and the conclusions $(\mathrm{C}_{42})$ and
$(\mathrm{C}_{Z})$ of the development.

\subsubsection{True types over ray-class primes}

\begin{proposition}\label{pf:baselat2:raytypes}
Assume $h_L\in\sfname{RecTExists}$, finiteness of automorphic forms on $\GL_n$, Berger's theorem (de Rham
representations are potentially semistable), Green's theorem on cuspidal representations, global Artin
reciprocity, the existence of ray class fields and the Chebotarev density theorem for primes of degree one.
Let $D$ be a global datum without localisation primes, $cf$ a character frame, $j_E:E\to\Qbar$ an embedding
and $\rho$ a row representation. Assume that the exponents $e_{\exp}(D,j,g)$ are integers, that the types of
$D$ are true for $cf$, that $\rho$ has the depth property, and that each auxiliary prime $g_k$ is a ray prime
for the field cut out by $j_E$ and the modulus $QS\cdot q\cdot d$. Then there are a new layer $L$ of $D$
(with the conditions $Z$ and $\eta$) and a character-frame embedding $A$ of the modified datum $D_L$ such
that: the Galois side of $A$ is true; its automorphic side is realised; the representation parameters of
$\rho$ hold for $D_L$ and the local parameters of $A$; the types of $L$ are true for $cf$; the places of type
$g$, the type data and the central characters of $L$ agree with those of $D$; the generators $\zeta_k$ of $A$
are the label generators at $g_k$; the localisation data and the field $\Omega$ of $L$ are those of $D$; and
the labels of $A$ are compatible with complex conjugation.
\formallabel{GlobalTheory.T.types_true_of_ray_class_alt}
\formallabel{GlobalTheory.HeightCounting.TopA.exists_sideA_alt}
\end{proposition}

\begin{proof}
\emph{Step 1: a layer with tame generators.} By
\lref{GlobalTheory.T.exists_tame_generators_green_cuspidal_labelling} (from Green's theorem and Artin
reciprocity: for a datum without localisation primes, with integral exponents and true types, there is a
layer $L$ without localisation primes whose types are true, whose places of type $g$, type data, central
characters, localisation data and field $\Omega$ are those of $D$, such that $D_L$ has a group side $\Gamma$ whose
generators $\zeta_k$ are the label generators and $D_L$ satisfies the compatibility condition on the characters
$\chi_T$ for the frame $cf_L$ induced by $cf$) we obtain $L$ and $\Gamma$.

\emph{Step 2: depth.} The depth property of $\rho$ only involves the places of type $g$ and the type data,
which are the same for $D$ and $D_L$; so $\rho$ has the depth property for $D_L$.

\emph{Step 3: the conclusion for $D_L$.} We show that $D_L$, $cf_L$, $j_E$, $\rho$ satisfy the conclusion of
the proposition. (a) By \lref{GlobalTheory.RowPar.exists_rowWD_rowTrue_m_alt} ($h_L$ and the depth
property give a Weil--Deligne row datum $R_w$, with local parameters and saturations, for which the
representation parameters of $\rho$ hold) we obtain $R_w$. (b) By
\lref{GlobalTheory.HeightCounting.exists_setAut_wdTie_labConj} (for a datum without localisation primes
and a Weil--Deligne row datum there are an automorphic modification of the datum, with the same places of
type $g$, a Weil--Deligne tie for it whose labels are compatible with conjugation, and row data whose local
parameters, read through the tie, and saturations are those of $R_w$) we obtain a modification $D^{\mathrm{aut}}$
of $D_L$ with these data. (c) Since the places of type $g$ agree, the tame generators of $\Gamma$ provide local
generators at the places of type $g$ of $D^{\mathrm{aut}}$, the types of $D^{\mathrm{aut}}$ are true, the condition on
$\chi_T$ transfers, and by (a), (b) the representation parameters of $\rho$ hold for $D^{\mathrm{aut}}$ with the
local parameters of its row data. (d)
\Lref{GlobalTheory.HeightCounting.TopA.conclusion_of_local_parameters_true} (from $h_L$, finiteness of
automorphic forms, Berger's theorem, ray class fields and Chebotarev: under the ray-prime hypothesis, a
datum without localisation primes with true types, a group side with label generators, the condition on
$\chi_T$, a Weil--Deligne tie with conjugation-compatible labels and row data for which the representation
parameters hold satisfies the conclusion of the proposition) gives the conclusion for $D^{\mathrm{aut}}$. (e) The
automorphic modification is itself the modification of $D_L$ by a layer, layers compose, and the equalities
in (b) identify places of type $g$, type data and central characters; hence the conclusion for $D^{\mathrm{aut}}$ is
the conclusion for $D_L$.

\emph{Step 4.} The same transport along the layer $L$ of Step 1, using the equalities listed there, turns
the conclusion for $D_L$ into the conclusion for $D$.
\end{proof}

\subsubsection{Existence of the automorphic side}

\begin{proposition}\label{pf:baselat2:autside}
Assume global Artin reciprocity. Let $D,cf,j_E,A$ be as above, $\mathrm{low}$ a finite set of primes, $P_q$ a
selection map, $S$ a rational datum with Hecke operators $h'_L$, an eigenvalue system $\mathrm{eig}$ and
characters $\theta_0,\theta_1$, and $B_s$ side data for these; let $M$ be a module family, $C$ a norm choice and
$\mathrm{tp}$ a tied tuple for $C$. Assume:
\begin{enumerate}[label=(\alph*)]
\item the choice-side condition for $(M,B_s,C,\mathrm{tp})$ and the module input condition for $(M,\mathrm{tp},C.Q)$;
\item $L_0$ is stable under complex conjugation, the exponents $e_{\exp}(D,j,g)$ are integers, the places
  above $p_{de}$ and $p'_{de}$ split, and the Chebotarev splitting statement holds for $D$;
\item at every place $w$ with $\mu_w\neq\mu_0$ the character $\mathrm{tp}.\psi_w$ takes values in $L$;
\item the map $\OO\to L$ is injective with range the integral elements of a subfield of $L$;
\item the blocks of $D$ are unique.
\end{enumerate}
Then the automorphic side exists for $(A,M,C,\mathrm{tp})$.
\formallabel{GlobalTheory.T.side_rational_hecke_data_alt}
\formallabel{GlobalTheory.HeightCounting.locAll_exists_alt}
\end{proposition}

\begin{proof}
Fix $\mathrm{big}$. The local datum is constructed in six steps, each provided by a statement of the
development under the hypotheses of the proposition (Artin reciprocity and (a)--(d)).
(1) By \lref{GlobalTheory.T.side_rational_eigen_data_alt}, together with the standard facts on Hecke
double coset operators (\lref{GlobalTheory.HeckeStdFacts.heckeStdPlain}), there are a localisation datum
$X$, local twists $T_X$ of $X$ and a specialisation $\mathrm{Sp}$ for $(A,M,C,\mathrm{big},X,T_X)$.
(2) By \lref{GlobalTheory.T.side_rational_level_data_alt}, for these there are local automorphic data $Y$
and base-change data $C_Y$ with an automorphic specialisation for $(A,M,\mathrm{tp},\mathrm{Sp},Y,C_Y)$.
(3) By \lref{GlobalTheory.HeightCounting.locBr_of_autSpec_alt}, the automorphic specialisation yields the
further datum $B$.
(4) By \lref{GlobalTheory.T.side_rational_character_data_alt}, for every local index $i$ and every place
$v$ of type $k$ with $\ell_i\neq g_k$ there is $u\in(\Z/g_k\Z)^\times$, equal to $\chi_T(k,v,j_i-1)$ or to its inverse,
such that the image of $\chi_i(1,t_g(k,v))$ under the coefficient embedding of $i$ is the Teichm\"uller lift of $u$.
(5) By \lref{GlobalTheory.HeightCounting.MtIndep.typeV_iSup_indep_alt}, for local indices $i,i'$ with the
same tuple $\mu$ and the same $j$ and for every parameter $\sigma'$, the sums of the type spaces of the
base-change vectors with local parameters $\sigma'$ coincide.
(6) By \lref{GlobalTheory.HeightCounting.locTy_exists} (using (e), the Chebotarev splitting statement and
(4), (5)) there is a type datum $Z$ for $Y$ satisfying the base-change, residual and module-parameter ties,
all of whose residual representations are semisimple.
The data $(X,T_X,Y,C_Y,Z,B)$ with these ties form a local datum, and $\mathrm{Sp}$, the automorphic
specialisation of (2) and the semisimplicity of (6) show that it is built for $\mathrm{big}$.
\end{proof}

\begin{corollary}\label{pf:baselat2:builtex}
Assume $\mathcal I_3$. Let $x$ be a state over $P$ (at a witness over $(B,f,T)$) with a
strict-Henselian-acyclic level $l$. Then the automorphic side exists for
$(l.a.A,\ l.\mathrm{mt}.M,\ l.c.\mathrm{ch},\ l.c.\mathrm{tp})$.
\formallabel{GlobalTheory.HeightCounting.Fin3.builtEx_of_published_inputs'}
\end{corollary}

\begin{proof}
Apply \cref{pf:baselat2:autside} to $D_x$, its frame, the base data of $l$, the low set of $D_x$, the
selection map of $l$, the rationality datum of $l$ and its side data $l.b$, the module family of $l$ and the
norm choice and tied tuple of $l.c$; global Artin reciprocity is part of $\mathcal I_3$. The hypotheses hold:
(a) the choice side is part of $l.c$, and the module input condition is the one of $l.\mathrm{mt}$, rewritten with
$\mathrm{tp}.\xi=w.\xi$ and $\mathrm{ch}.Q=Q_P$ (both part of $l.c$), since the tuple of $\mathrm{tp}$ is $\mu_0\cdot\mathrm{tp}.\xi$;
(b) conjugation stability, integrality of the exponents and the splitting of the places above $p_{de},p'_{de}$
form the elementary datum of the same-Galois level, and the Chebotarev splitting statement follows from the
Chebotarev theorem contained in $\mathcal I_3$ (\lref{GlobalTheory.ChebSplitDense.chebSplit_nd});
(c) $\mathrm{tp}.\psi=T.\psi$ (part of $l.c$) and the characters of a coherent twisting datum take values in $L$
(\lref{GlobalTheory.twDat_psi_valuesIn_LL}); (d) is part of the same-Galois level; (e) uniqueness of blocks
is a component of the truth of the Galois side of the same-Galois level, given finiteness of automorphic
forms, the cuspidal statement for unitary groups, the Jacquet--Shalika theorem and Chebotarev, all
contained in $\mathcal I_3$.
\end{proof}

\subsubsection{Large data and the eta equality}

\begin{proposition}[Forcing by large data]\label{pf:baselat2:large}
Assume the Chebotarev density theorem in the form: every coset of an open normal subgroup of the
absolute Galois group of a number field contains the Frobenius of an unramified prime of degree one
outside any given finite set of primes. Let $D,cf,j_E,A,\mathrm{low},P_q,S,h'_L,\mathrm{eig},\theta,B_s,M,C,\mathrm{tp}$ be
as in \cref{pf:baselat2:autside} and satisfy its hypotheses (a) and (b). Assume that the automorphic side
exists for $(A,M,C,\mathrm{tp})$ and is forced by the property
``$(\mathrm{O})$ implies $(\mathrm{P}_{41})$ and residual irreducibility, and $(\mathrm{GR})$ holds''.
Then it is forced by largeness.
\formallabel{GlobalTheory.CoefDuty.large_of_large_rings_alt}
\end{proposition}

\begin{proof}
We use the laws of forcing throughout.
(1) $(\mathrm{O})$ is forced, by \lref{GlobalTheory.T.side_eigen_rationality_alt} (from the standard facts
on Hecke double coset operators and hypotheses (a), (b)).
(2) $(\mathrm{Spl})$ is forced: by \lref{GlobalTheory.T.side_multiplicity_cusp_alt} there is $\mathrm{big}$
such that $(\mathrm{Spl})$ holds for $D_{\mathcal L}$ whenever $\mathcal L$ has a specialisation for $\mathrm{big}$ and an
automorphic specialisation, in particular whenever $\mathcal L$ is built for $\mathrm{big}$.
(3) $(\mathrm{Dec})$ is forced: by hypothesis $(\mathrm{GR})$ is forced, and $(\mathrm{GR})$ implies $(\mathrm{Dec})$ by
\lref{GlobalTheory.BlockComplement.complement_at} (from Chebotarev and attached Galois representations:
every Hecke--Galois representation $r$ at a point of the theta locus has a complement $\rho$ realising the
decomposition of $r$).
(4) By (1) and the hypothesis, ``$(\mathrm{O})$ and [$(\mathrm{O})$ implies $(\mathrm{P}_{41})$ and residual
irreducibility]'' is forced, hence so is $(\mathrm{P}_{41})$. With (2) and (3), $(\mathrm{P}_{41})\wedge(\mathrm{Spl})\wedge(\mathrm{Dec})$
is forced, and \lref{GlobalTheory.T.side_eigen_hecke_match_alt} (from the Chebotarev theorem, the
standard facts on Hecke operators, (a), (b) and the existence of the automorphic side: if this conjunction is
forced then so is $(\mathrm{Tk})$) shows that $(\mathrm{Tk})$ is forced.
(5) By (1), (2), (4) and the conjunction law, $(\mathrm{O})\wedge(\mathrm{Spl})\wedge(\mathrm{Tk})$, i.e.\ largeness, is
forced.
\end{proof}

\begin{proposition}[Eta equality]\label{pf:baselat2:etaj}
Assume $\mathcal I_3$ and $\mathrm{Pt}(h_L)$ for $h_L\in\sfname{RecTExists}$; recall that $\mathrm{Pt}(h_L)$ states
$\forall^{\mathrm{st}}$ of: for every strict-Henselian-acyclic level $l$ of $x$ the automorphic side for
$(l.a.A,l.\mathrm{mt}.M,l.c.\mathrm{ch},l.c.\mathrm{tp})$ is forced by the property ``$(\mathrm{O})$ implies
$(\mathrm{P}_{41})$, residual irreducibility, $(\mathrm{C}_{42})$ and $(\mathrm{C}_Z)$; the eigenvectors span; and
the isotypic decompositions of the rows are internal''. Then $\forall^{\mathrm{st}}$ of: for every
strict-Henselian-acyclic level $l$ of $x$ (with $x$ over $P$) the automorphic side for
$(l.a.A,l.\mathrm{mt}.M,l.c.\mathrm{ch},l.c.\mathrm{tp})$ is forced by the property ``$(\mathrm{O})$ implies
$D'.\eta_J(j,v)=P.q.\eta(j,v)$ for all $1\le j\le m$ and all places $v$''.
\formallabel{GlobalTheory.T.eventually_etaJ_eq_eta_of_levelC}
\formallabel{GlobalTheory.EtaMem.exists_etaJ_alt}
\end{proposition}

\begin{proof}
Besides $\mathrm{Pt}(h_L)$ we use two statements of the same form:
\lref{GlobalTheory.T.allSt_goodFTAll_of_levelA} (under $\mathcal I_3$, at every state with a same-Galois
level all finite types of $D_x$ are good) and
\lref{GlobalTheory.T.twistedPlace_innerProduct_classSum_rows} (at every row configuration, every module
family adapted to the read package which satisfies the module input condition has the class-sum
property: at every twisted place the inner product on the module is a positive multiple of the class-sum
pairing attached to an invariant form on the weight representation and to class representatives of the
bounded level group). The Chebotarev theorem of $\mathcal I_3$ gives the Chebotarev splitting statement for
every datum (\lref{GlobalTheory.ChebSplitDense.chebSplit_nd}). By the closure properties it suffices to
argue pointwise.

Let $x$ be a state over $P$ and $l$ a strict-Henselian-acyclic level of $x$. Let $\mathrm{big}$ be given by
$\mathrm{Pt}(h_L)$; we show that it also forces the eta equality. Let $\mathcal L$ be built for $\mathrm{big}$, put
$D'=D_{x,\mathcal L}$ and assume $(\mathrm{O})$ for $D'$. Then:
\begin{enumerate}[label=(\roman*)]
\item $\mathrm{Pt}(h_L)$ gives $(\mathrm{C}_{42})$, $(\mathrm{C}_Z)$, the spanning of the eigenvectors and the internality of
  the isotypic decompositions for $D'$; the finite types of $D'$ are good, since they are those of $D_x$;
\item the module input condition holds for $(l.\mathrm{mt}.M,\mathrm{tp},\mathrm{ch}.Q)$ (as in the proof of
  \cref{pf:baselat2:builtex}); the $Q$-data of $D'$ are those of $\mathrm{ch}$ at the localisation places (part of
  the specialisation of $\mathcal L$), and the lattices of the localised module family are integral at every
  place $w$ with $\mu_w\ne\mu_0$ (part of the module input condition);
\item the class-sum property holds for the module family of $l$ (second statement above, at the row
  configuration of $x$) and is transported to the localised module family, whose modules are those of
  $l.\mathrm{mt}.M$ at the corresponding tuples;
\item by \lref{GlobalTheory.FitHecke.heckeOpen_withLoc_alt} (from the side data, the choice side, the
  module input, the elementary datum and the Chebotarev splitting statement) the Hecke operators of $D'$ on
  the localised module family are tied to the automorphic side, rational, and compatible with the
  automorphic vectors;
\item the specialisation of $\mathcal L$ provides a normalised local compatibility datum for
  $D'$ and the localised module family (its base, maximal ideals, maps $\Psi$ and points are those of the
  specialisation, and the coherence of the complex embeddings follows from the automorphic specialisation), and
  the standard facts for the Hecke operators of its base are those of the specialisation.
\end{enumerate}
By \lref{GlobalTheory.T.etaJ_local_compatibility_alt} (under $\mathcal I_3$, at a state with a
strict-Henselian-acyclic level, for a local datum built for $\mathrm{big}$ such that $D'$ satisfies $(\mathrm{O})$
and for a normalised local compatibility datum, the data (i)--(v) imply $D'.\eta_J(j,v)=P.q.\eta(j,v)$ for
$1\le j\le m$ and all places $v$) the eta equality holds.
\end{proof}

\subsubsection{Twisted forms and side data}

A \emph{size bound} for a datum $D$ is a family $b=(b(\ell,\iota))$ of finite extensions of $\Q_\ell$ (inside an
algebraic closure), indexed by primes $\ell$ and embeddings $\iota:L_0\to\Qbar_\ell$; $c\le b$ means
$c(\ell,\iota)\subseteq b(\ell,\iota)$ for all $(\ell,\iota)$. A \emph{ring specification} for $(cf,b_0,b_1,S_x,T)$,
where $T$ is a family of twisting rings for $D$ (twisted indices $i$ with primes $\ell_i$, rings $\OO_i$,
exceptional sets $S_0(i)$ and embeddings), identifies the twisted indices with the primes of $L_0$ and the
twisting rings with rings of integers of local fields bounded by $b_0,b_1$, and contains the characters and
embeddings of the twisting.

\begin{theorem}[Twisted forms and side data]\label{pf:baselat2:secondstep}
Assume global Artin reciprocity. Let $D$ be a global datum with a character frame $cf$, an embedding $j_E$ and
a character-frame embedding $A$; let $S$ be a rational datum with a rational input (a finite set $S_b$ of
places, Hecke operators, eigenvalues and characters $\theta$), $\mathrm{low}$ a finite set of primes and $P_q$ a
selection map with $P_q(v)$ above $v$ for all $v$, and let $M$, $\mathrm{tup}$, $Q$ satisfy the module input
condition. Assume that $D$ has Galois representations, that the rows of $A$ are cuspidal, the Chebotarev
splitting statement for $D$, that $L_0$ is stable under complex conjugation, that the exponents are
integers and that the places above $p_{de},p'_{de}$ split. Then there is a size bound $c_0$ with the following
property. Let $b_0\ge c_0$ and $b_1$ be size bounds, $S_x$ a finite set of places of $E$ which contains $S_b$,
is closed under conjugation and contains all places above the primes of $\mathrm{low}$, $T$ a family of
twisting rings and $\mathrm{Sp}$ a ring specification for $(cf,b_0,b_1,S_x,T)$. Then there are a twisting frame
$F$ and an expansion datum $\mathrm{ex}$ for the twisting $N$ defined by $T$, the bare part of $A$ and $F$, such
that every sup lattice $B^{\sup}(i,j,b)$ of $\mathrm{ex}$ is finitely generated and spans the forms over the
fraction field of $\OO_i$, and for every theta datum $\mathrm{th}$ the new datum $D_N$ (formed with $\mathrm{ex}$,
$\mathrm{th}$ and the base ring data of $\mathrm{Sp}$) carries side data for the transported frame, rows and
rational datum, with the same $\mathrm{low}$, $P_q$, Hecke operators, eigenvalues and characters $\theta$.
\formallabel{GlobalTheory.T.row_cusp_multiplicity_input_alt}
\formallabel{GlobalTheory.HeightCounting.FormsStep.second_step_alt}
\end{theorem}

\begin{remark}
The module family $M$, the tuple $\mathrm{tup}$, the set $Q$ and the module input condition are hypotheses of the
formal statement but are not used in its proof.
\end{remark}

\begin{proof}[Proof of \cref{pf:baselat2:secondstep}]
\emph{The bound.} By \lref{GlobalTheory.T.exists_galois_rep_family_bounded_alt} (from Artin reciprocity,
the Galois representations of $D$, row cuspidality, Chebotarev splitting and conjugation stability) there is
a size bound $c_1$ such that for all data as in the theorem with $b_0\ge c_1$ there is a family $\rho(i,\sigma,b)$
of representations with values in $\GL_r$ over the twisting rings such that, for every good row $\sigma$ with
nonzero isotypic part and every conjugate self-dual regular algebraic cuspidal $\pi$ of $\GL_r(\mathbb A_E)$
whose components at split places are the local components of the row, $\rho(i,\sigma,b)$ (pushed to $\Qbar_{\ell_i}$
compatibly with the embeddings of $L_0$) is the $\ell_i$-adic representation attached to $\pi\otimes\eta_T^{-1}$
with the normalisation $(1-n_V)/2$. By \lref{GlobalTheory.T.exists_weight_data_integral_alt} (from Artin
reciprocity, conjugation stability and integrality of the exponents) there is a size bound $c_2$ such that
for $b_0\ge c_2$ there are \emph{weight data}: for all $i,b$ and $1\le j\le m$ an irreducible representation $W(i,j,b)$
over the fraction field of $\OO_i$ which is algebraic of some weight (in the sense of the development), whose base
changes to the twisted rings $\OO_{i,k}$ remain irreducible, which has a highest weight at every infinite place,
and whose highest weights have the infinitesimal character prescribed by the archimedean parameters of the
$j$-th row. Let
$c_0(\ell,\iota)$ be the compositum of $c_1(\ell,\iota)$ and $c_2(\ell,\iota)$; then $b_0\ge c_0$ implies $b_0\ge c_1$ and
$b_0\ge c_2$.

\emph{The twisting.} Fix $b_0\ge c_0$, $b_1$, $S_x$, $T$, $\mathrm{Sp}$, and let $\rho$ and the weight data be as
above. The following hold: the twisted lattices of the rational datum are finitely generated and span
(\lref{GlobalTheory.T.twisted_lattice_finitely_generated_spanning}); they are stable under the Hecke
operators satisfying the Hecke condition (\lref{GlobalTheory.T.hkStab_of_subset_of_star_closed}, using
$S_b\subseteq S_x$ and conjugation stability); there is signature data $\mathrm{sd}$ spanning the forms, with
injective signatures, rows and eigenvalue data (\lref{GlobalTheory.T.exists_signature_data_span_alt});
the central idempotents of the group algebra are universal
(\lref{GlobalTheory.T.isCentralGroupAlgebraIdempotent_universal}); and the generic unipotent character
of the twisting is the character attached to the theta side of the characters $\theta$ read through the
twisting, which are nontrivial if $r\ge2$ (\lref{GlobalTheory.T.genericUnipotentCharacterAtSide_eq_thetaSide}). The twisting frame $F$ is formed
from the twisted lattices, the Hecke stability, $\mathrm{sd}$ and $\rho$; the expansion datum $\mathrm{ex}$
additionally uses the idempotents, $P_q$ and the weight data. Its sup lattices are the twisted lattices,
which proves the first assertion.

\emph{Side data.} Let $\mathrm{th}$ be a theta datum. The side data of $D_N$ consist of the following
components.
(i) The twist frame and the excluded places are those of $\mathrm{Sp}$, with $S_x$ as extra excluded places
(closed under conjugation and containing the places over $\mathrm{low}$ by hypothesis).
(ii) The tied expansion: the clauses for the unitary group and the characters $\psi_U$, from the generic
unipotent character above; continuity of the universal anticyclotomic character
(\lref{GlobalTheory.GalChar.universal_anticyclotomic_character_laws}); for admissible sets $Q$ and $\mathfrak P\in Q$,
its $\mathfrak P$-component is surjective on the inertia groups at $\mathfrak P$ and trivial on those at the places of $Q$
other than $\mathfrak P$ and its conjugate (\lref{GlobalTheory.T.inertia_surjective_of_admQ}); the block
characters are tame at the places of type $g$ not above $\ell_i$
(\lref{GlobalTheory.T.hasTameBlockCharacter_typePlace_not_over_l}); and real
module data at admissible twists, formed from the weight data and from the levels given by
\lref{GlobalTheory.T.places_over_real_subfield_split_alt} (for every admissible twist $(i,k,Q,\psi)$ and
$1\le j\le m$ there is a level satisfying the level laws whose type space in the base change of $W(i,j,b)$ is
finite-dimensional).
(iii) The Hecke splitting: for every place $\mathfrak P\notin S_0(i)$ the characteristic polynomial of the Hecke
operator at $\mathfrak P$ on the twisted lattice splits over the fraction field of $\OO_i$
(\lref{GlobalTheory.T.lattice_basis_free_module_place_alt}).
(iv) Every place $w$ of $L$ is covered by a twisted index: there is $i$ with $\ell_i$ the residue characteristic of
$w$ such that an element of $L_0$ is $w$-integral if and only if its image under the embedding of $L_0$ attached
to $i$ lies in the base ring $\OO^0_i$ of the twisting
(\lref{GlobalTheory.T.exists_twistIndex_of_coefficientField_place}).
(v) The forms of $D_N$ are the base change of the forms of the rational datum, compatibly with the sup
lattices, the group action and the Hecke operators
(\lref{GlobalTheory.T.twistLiftForms_eq_baseChange_rationalDatumForms}); for $\mathfrak P\notin S_0(i)$ and a
Frobenius element $\sigma$ at $\mathfrak P$, the Hecke comparison unit at $\mathfrak P$ is $\delta_i(\sigma)^{\pm1}$, the sign
depending on $b$ (\lref{GlobalTheory.T.heckeComparisonUnitAtFrobenius_eq_delta}).
(vi) The signature span, the rows with their eigenvalue data and the Galois conditions on the rows come
from $\mathrm{sd}$ and from the property of $\rho$ stated above.
The remaining components (characters, embeddings and their compatibilities) are copied from $\mathrm{Sp}$ and
from the rational input.
\end{proof}

\subsubsection{Series data on a twisted copy}

\begin{proposition}\label{pf:baselat2:sercopy}
Assume $\mathcal I_3$ and $\mathcal I_{\mathrm{ser}}$. Let $x$ be a state over $P$ at a starting witness over
$(K,B,f,T)$ with $K$ over the pre-base of $B$ and $T$ forced, and let $l$ be a first intermediate level of
$x$. Then there is a new twisting $N$ of $D_x$ which preserves full $\lambda$-ties (every $\lambda$-datum whose tie
with the twisting of $D_x$ is full has a full tie with $N$) with the following property. Let $x'$ be a state
over $P$ whose over-datum is the twist $o_{\mathrm{tw}}(x,N)$ of that of $x$, whose real datum, group condition
and geometric datum are the transports of those of $x$, and let $l'$ be an intermediate level of $x'$ whose
rows, selection map, rational datum, set $S_b$, Hecke operators, eigenvalues and characters $\theta$ are the
transports of those of $l$. Then for every module input $m$ of the read package of $x'$ whose adelic data
are identity frames, series data exist for $(x',l',m)$.
\formallabel{GlobalTheory.T.series_copy_embedding_of_cmDatum_alt}
\formallabel{GlobalTheory.HeightCounting.SerCopy.series_copy_embedding_alt}
\end{proposition}

\begin{proof}
(1) Since $K$ is over the pre-base and $w$ is a starting witness, the same-Galois level of $l$ satisfies the
conditions of a choice (\lref{GlobalTheory.T.isChoice_of_start_of_over}); in particular every root of
unity of order dividing $g_k(g_k-1)$ lies in $L_0$.
(2) By \lref{GlobalTheory.HeightCounting.SerCopy.exists_prepared} (Artin reciprocity: given side data, the
valuation datum and these roots of unity, there is a new twisting $N$ which preserves full $\lambda$-ties, each of
whose twisted indices has the exceptional set $S_0$ and the prime $\ell$ of a twisted index of $D$, which has no
off-index places and universal characters, and for which there are side data with true $\tau$), applied to the
side data of $l$, we obtain $N$ and side data $B'_s$ for $D_N$ with true $\tau$.
(3) For every twisted index $i'$ of $N$, $S_b\subseteq S_0(i')$ (as $S_0(i')=S_0(i)$ for a twisted index $i$ of $D_x$
and $S_b\subseteq S_0(i)$ by the side data of $l$), and $S_0(i')$ lies over $\{\ell_{i'}\}\cup\mathrm{low}$, by the
corresponding property of $l$.
(4) The $\lambda$-datum of $l$ has a full tie with the twisting of $D_x$, hence with $N$, which gives a
$\lambda$-datum for $N$.
(5) By \lref{GlobalTheory.HeightCounting.SerCopy.pre_at_twisted_copy_of_side_b_tau_true} (Artin
reciprocity: from side data with true $\tau$, (3) and a $\lambda$-datum, there is a \emph{pre-series datum} $p$ at
the twisted datum, i.e.\ side data for the twisted real datum with true $\tau$, the transported rows, selection
map, Hecke operators, eigenvalues and characters $\theta$, together with a pre-tie with the geometric datum) we
obtain $p$. The twisting of the proposition is $N_p$: the twisting $N$ with the expansion data and theta data
replaced by those of $p$; it has the same twisting as $N$, so it preserves full $\lambda$-ties.
(6) Let $x',l',m$ be as in the statement. The moved real datum of $x'$ is the twisted real datum of $x$
(\lref{GlobalTheory.HeightCounting.Up.moved_eq_datTw}), and
\lref{GlobalTheory.HeightCounting.SerCopy.series_state_b1_twist_alt} constructs series data for
$(x',l',m)$ from $p$, the truth of $\tau$, this identification, the read package of $x'$, the data of $l'$, the
identity frames, the $\lambda$-datum of (4) and the inclusions of (3).
\end{proof}

\subsubsection{Geometric laws at the generator}

\begin{proposition}\label{pf:baselat2:geolaws}
Assume $\mathcal I_3$ and $\mathcal I_{\mathrm{ser}}$. Then $\forall^{\mathrm{st}}$ of: for every second intermediate
level $l$ of $x$ and every twisted index $i$ of $D_x$, the series data of $l$ satisfy the \emph{geometric laws}
at the generator $i_C(i)$ of the theta data of $D_x$, with respect to the frame tie of the geometric datum
of $x$: for every sign $b$ and $j<m$,
\begin{enumerate}[label=(\roman*)]
\item (detection) a form in $F_{i,j+1,b}$ whose values at all points over $i_C(i)$ vanish is zero;
\item (sup descent) the sup lattice of the series side at $(i,j+1,b)$ is the $\OO_i$-span of the unit ball of
  the reading at the points over $i_C(i)$;
\item (pin) a form of the expansion datum lies in its sup lattice if and only if its image in the forms of
  the series side lies in the sup lattice of the series side;
\item (Hecke) for every Hecke index $(k,\mathfrak P)$ of $i$ there is $s$ in the coefficient field, of absolute
  value $1$ under the embedding attached to $(i_C(i),b)$, such that the Hecke operator at $(k,\mathfrak P)$ is given
  on points by the Hecke relation with scalar $s$;
\item (stability) for every Hecke index $(k,\mathfrak P)$, every point $z$ over $i_C(i)$ and every subgroup $C$ in the
  Hecke kernel at $z$, some point over $i_C(i)$ is the Hecke quotient of $z$ by $C$;
\item (count) for every Hecke index $(k,\mathfrak P)$ there is $n$ such that the Hecke kernel at every point over
  $i_C(i)$ has $n$ elements;
\item the forms $F_{i,j,b}$ vanish for $j\notin[1,m]$.
\end{enumerate}
\formallabel{GlobalTheory.T.geoLaws_iC_at_levelB1_of_serN}
\formallabel{GlobalTheory.HeightCounting.Up.geoLaws_at_levelB1_states}
\end{proposition}

\begin{proof}
By \lref{GlobalTheory.T.heckeGeo_at_levelB1_states}, (iv) holds at every second intermediate level and
every twisted index, under $\mathcal I_3$ and $\mathcal I_{\mathrm{ser}}$, in the form $\forall^{\mathrm{st}}(\cdot)$. By the
closure properties it suffices to deduce the other laws pointwise. The generator $i_C(i)$ is integral (part of
the theta data). (i) is \lref{GlobalTheory.PointsDetectAlg.detects} (at an integral generator, the points
detect forms), and (ii) is \lref{GlobalTheory.GeoTie.supDescent_side} (at an integral generator, the sup
lattice of the series side is spanned by the unit ball of the reading). (iii) and (vii) are part of the series
data of $l$. (v) is \lref{GlobalTheory.T.heckeQuotient_integral_framed_point_complete_field} (from the
Hecke integrality input of $\mathcal I_{\mathrm{ser}}$: Hecke quotients of framed points over integral generators are
again such points). For (vi), a Hecke index $(k,\mathfrak P)$ makes $\mathfrak P$ a split Hecke prime off the level
(\lref{GlobalTheory.T.isSplitHeckePrimeOffLevel_of_heckeIndexPlace}), and for such primes the number of
elements of the Hecke kernel does not depend on the point
(\lref{GlobalTheory.HeckeCount.card_hkKer_const}).
\end{proof}

\subsubsection{Character data and the comparison of levels}

Let $\mathrm{ch}$ be a norm choice of a datum $D$. For a place $w$ with $\mu_w\neq\mu_0$ and $1\le j\le m$, $\mathrm{ch}$
determines a discrete valuation ring $\OO_e(w)$ with fraction field $K_e(w)$, an embedding $e_w:L\to K_e(w)$, a
twisted index $\mathrm{ch}.\iota(w)$, and the \emph{target} $T_{\mathrm{ch}}(w,j)$, the base change to $K_e(w)$ of the twisted
module of the expansion at $(\mathrm{ch}.\iota(w),j)$ with the data $(k'_w,Q_w,\psi'_w)$ of $\mathrm{ch}$, with a lattice
$\Lambda_T(w,j)$. Let $M_L(w,j)$ be the $L$-module of the tuple $\mu_w$ at $j$, with the lattice $\Lambda_S(w,j)$ spanned
by its generators. A \emph{comparison map} at $(w,j)$ is an $e_w$-semilinear map $e:M_L(w,j)\to T_{\mathrm{ch}}(w,j)$
which maps every $L$-basis to a $K_e(w)$-basis; the norm choice contains comparison maps $e_{\mathrm{ch}}(w,j)$ with
$\mathrm{span}_{\OO_e(w)}e_{\mathrm{ch}}(w,j)(\Lambda_S(w,j))=\Lambda_T(w,j)$.

\begin{proposition}[Character data with prescribed twisted index]\label{pf:baselat2:choice}
Assume the discrete series cover, the Casselman--Wallach theorem, the theory of Harish-Chandra modules,
$\mathcal I_3$ and $\mathcal I_{\mathrm{ser}}$. Then $\forall^{\mathrm{st}}$ of the following. Let $x$ be a state over $P$ at a
witness $w$ over $(B,f,T)$ and $l$ a third intermediate level of $x$; let $\mathrm{ch}$ be a norm choice of $D_x$ and
$\mathrm{tp}$ a tied tuple for $\mathrm{ch}$ such that $\mathrm{tp}.\xi$ is the character family of $w$, $\mathrm{tp}.\psi=T.\psi$,
$\mathrm{ch}.Q$ is the $Q$-datum of $P$, $\mathrm{ch}.k'_v=T.k_v$ at twisted places $v$, the choice-side condition holds for
$(l.\mathrm{mt}.M,l.b,\mathrm{ch},\mathrm{tp})$, and the twist of $\mathrm{ch}$ is admissible at every twisted place. Then there are
height-counting character data $c$ for $(l.\mathrm{mt},l.b,T.k,T.\psi)$ whose norm choice has the twisted index map
of $\mathrm{ch}$.
\formallabel{GlobalTheory.T.nonGalois_exists_choice_alt}
\formallabel{GlobalTheory.ChoiceRead.exists_stC_choice_alt}
\end{proposition}

\begin{proof}
By \lref{GlobalTheory.T.local_component_at_place_determined_alt} (same hypotheses; a statement of the
form $\forall^{\mathrm{st}}(\cdot)$), in the situation of the proposition there are, for every twisted place $v$ and every
$i<m$, a comparison map $e_{v,i}$ at $(v,i+1)$ and a unit $u_{v,i}\in K_e(v)^\times$ such that
$\mathrm{span}_{\OO_e(v)}e_{v,i}(\Lambda_S(v,i+1))=u_{v,i}\Lambda_T(v,i+1)$ and the theta condition holds with the factor
$u_{v,i}$: for every $x$ in the theta part of $M_L(v,i+1)$ and every expression of the reading of the theta vector of
$x$ as a combination $\sum_ka_kf_k$ of rational forms, the theta image of $e_{v,i}(x)$ is $u_{v,i}$ times the
corresponding combination $\sum_ke_v(a_k)\otimes f_k$ in the target. By the closure properties it suffices to argue
pointwise. Put $e'_{v,i}=u_{v,i}^{-1}e_{v,i}$. Multiplication by a unit maps bases to bases, so $e'_{v,i}$ is a
comparison map; $\mathrm{span}\,e'_{v,i}(\Lambda_S)=u_{v,i}^{-1}u_{v,i}\Lambda_T=\Lambda_T$; and the theta condition holds for
$e'_{v,i}$ with the factor $1$. Let $\mathrm{ch}'$ be the norm choice obtained from $\mathrm{ch}$ by replacing its comparison
maps at the twisted places $v$ and $1\le j\le m$ by $e'_{v,j-1}$. The tied tuple $\mathrm{tp}$ is a tied tuple for
$\mathrm{ch}'$, the conditions on $\xi$, $\psi$, $Q$, $k'$ and admissibility only involve data which are unchanged, the
choice-side condition transfers to $\mathrm{ch}'$, and the theta condition with factor $1$ is the remaining
requirement of height-counting character data. Hence $\mathrm{ch}'$ and $\mathrm{tp}$ form character data $c$, and
$\mathrm{ch}'$ has the twisted index map of $\mathrm{ch}$.
\end{proof}

For the comparison of levels we use the following notation. For a split place $s$ of $E$ (a place $s.v$ of the
maximal real subfield $F$ which splits, with a chosen place $s.w$ of $E$ above it), $\eta_s$ denotes the
component at $s$ of a Hecke character $\eta$, a character of $F_{s.v}^\times$. A \emph{level} is a pair $(J,\chi)$ of a
subgroup $J$ of $\GL_{n_V}(\mathbb A_{E,f})$ and a character $\chi:J\to\C^\times$; two levels are \emph{equal},
$(J,\chi)\equiv(J',\chi')$, if $J=J'$ and $\chi=\chi'$ on $J$. Let $g(D)$ be the set of places of $F$ below the places of
$E$ above $g_0$ or $g_1$, and $q(\eta)$ the set of places of $F$ outside $g(D)$ below a place of $E$ at which $\eta$ is
ramified; both sets are finite.

\begin{proposition}[Comparison of levels]\label{pf:baselat2:leveq}
Let $D$ be a global datum, $\eta$ a Hecke character of $E$, $m_0$ a real module datum at $(\mu,j')$ with level
$(J,\chi)$, and $X$ tied expansion data for $D$. Let $(i,k,Q,\psi)$ be an admissible twist, $b$ a sign, $j<m$,
let $(J',\chi'_0)$ be the level of the cell of $X$ at $(b,(i,k,Q,\psi),j)$ and let $e:K_{i,k}\to\C$ be an embedding
which is read by the side at $(i,b,k)$; put $\chi'=e\circ\chi'_0$. Assume:
\begin{enumerate}[label=(\alph*)]
\item the places of type $g$ lie above the $g_k$ and are the chosen places of split places, every place of
  $g(D)$ arises in this way, and some set of places of type $g$ is nonempty;
\item at every split place $s$ with $s.v\in q(\eta)$, $\eta_s$ is nontrivial on the units of $F_{s.v}$;
\item $J$ is a product level; at every place $v\in q(\eta)$, $(J,\chi)$ has the parahoric shape at a split place
  over $v$ with the component of $\eta$; at every place of type $g$, $(J,\chi)$ is a unipotent level for an
  additive character and the component of $\eta$; and the outside hypothesis holds for $(J,\chi)$ off
  $g(D)\cup q(\eta)$;
\item the bad places of $(i,Q)$ lie in $g(D)\cup q(\eta)$, and $(J',\chi')$ has the parahoric shape at the places of
  $q(\eta)$ and is a unipotent level at the places of type $g$ (for every embedding read by the side);
\item the joint condition: there are a frame-preserving automorphism $\theta$ of $\GL_{n_V}(\mathbb A_{E,f})$, a
  subgroup $J_1$, an embedding $e'$ read by the side, and an injective $e'$-semilinear map $\omega$ from the space
  of the cell to a complex vector space, such that $J_1$ contains the parahoric radicals at split places over
  the places of $q(\eta)$ and the unipotent radicals at the places of type $g$, and for $\kappa\in J_1$ the realisations of
  $m_0$ and of the cell (through $\omega$) are eigenvectors of $\theta(\kappa)$ with eigenvalues $c,c'$ satisfying
  $c^{g_k-1}=c'^{\,g_k-1}$ when $\kappa$ is supported at a place of type $g$ above $g_k$.
\end{enumerate}
Then $(J,\chi)\equiv(J',\chi')$.
\formallabel{GlobalTheory.LevCmp.levEq}
\end{proposition}

\begin{proof}
\emph{Shapes.} By (c), $J$ has the \emph{group shape}: it is a product level, at a split place over each $v\in
q(\eta)$ its elements supported at $v$ are those in the parahoric radical, at the places of type $g$ its elements
supported there are those in the unipotent radical, and off $g(D)\cup q(\eta)$ it coincides with the maximal compact
subgroup. The same holds for $J'$: it is a product level by
\lref{GlobalTheory.T.isProductLevel_definiteFormsRealisationAtTwist}, and (d) together with the outside
hypothesis for the cell (which follows from the one of the development and the location of the bad places)
gives the other properties.

\emph{Eigenvectors.} By \lref{GlobalTheory.LevCmp.exists_form} (using (a) and the group shape), each of the two
levels has a nonzero eigenvector in the corresponding representation space. For $m_0$ the eigenvector lies in
the range of the realisation, which is described by the level and its openness, so it is the realisation
$f=m_0.\mathrm{real}(y)$ of some $y$, and $f(g\kappa)=\chi(\kappa)f(g)$ for $\kappa\in J$. For the cell the eigenvector is
the realisation of some $y'$, and $F'=\omega\circ(\text{realisation of }y')$ is nonzero by injectivity of $\omega$ and
satisfies $F'(g\kappa)=e'(\chi'_0(\kappa))F'(g)$ for $\kappa\in J'$.

\emph{Local agreement.} By \lref{GlobalTheory.T.productLevel_eq_of_local_agreement}, two product levels
with group shapes are equal as soon as, at every place $u$ of the finite set $g(D)\cup q(\eta)$, they contain the
same elements supported at $u$ and their characters agree on them. Note $2\le r<n_V$, by the rank inequalities
of $D$.

Let $u\in g(D)$; by (a), $u=s.v$ for a split place $s$ with $s.w$ a place of type $g$ above $g_k$. By (c) and (d) both
levels are unipotent levels at $s$, for additive characters $\psi_0,\psi'_0$ and components $\xi,\xi'$ of $\eta$; by
\lref{GlobalTheory.T.heckeCharacter_splitPlaceComponent_unique} $\xi=\xi'$. By
\lref{GlobalTheory.T.additiveCharacter_eq_of_common_unipotentLevel_eigenvector} (given $\theta$, the joint
condition (e) and the nonzero eigenvectors $f,F'$; $g_k$ is prime, $s$ lies above it and $g_k-1$ is prime to $g_k$)
the two unipotent levels have the same character on the relevant elements. Elements supported at $s$ lie in $J$ if
and only if they lie in the unipotent radical, i.e.\ if and only if they lie in $J'$; for such $\kappa$ the lower
determinant is a unit and the superdiagonal block is integral
(\lref{GlobalTheory.T.upperParahoric_lowerDet_unit_and_superdiag_integral}), and both $\chi(\kappa)$ and
$\chi'(\kappa)$ are given by the same expression in the additive character, $\xi$ and this determinant; hence
$\chi(\kappa)=\chi'(\kappa)$.

Let $u\in q(\eta)$. By (c) and (d) there are split places $s,s'$ over $u$ at which $J$ and $J'$ have the parahoric
shape, with components $\xi,\xi'$ of $\eta$; by (e) there is a split place $s_1$ over $u$ whose parahoric radical
lies in $J_1$. By \lref{GlobalTheory.LevCmp.para_rigid} (using (b): an element of the maximal compact
subgroup supported at $s$ which acts by a scalar on a nonzero eigenvector of a parahoric level at $s$ lies in the
parahoric radical at $s$), applied to $f$ and to $F'$ and to the images under $\theta$ of the elements of the
parahoric radical at $s_1$ (which act by scalars by (e)), these images lie in the parahoric radicals at $s$ and at
$s'$; by \lref{GlobalTheory.T.splitPlace_eq_of_automorphism_maps_parahoric} (for split places over the same
place, this inclusion for a frame-preserving automorphism forces equality of the split places) $s'=s$. Then
$\xi'=\xi$ by uniqueness of components. Elements supported at $u$ lie in $J$ if and only if they lie in the parahoric radical,
i.e.\ in $J'$, and for them the lower determinant is a unit; both characters are $\xi$ of this determinant,
hence agree.
\end{proof}

\subsubsection{Movement to the base}

\begin{lemma}\label{pf:baselat2:tob1}
Assume $\mathcal I_3$ and $\mathcal I_{\mathrm{ser}}$, and let $h_L\in\sfname{RecTExists}$. Then the statement
$\mathcal{W}(\mathrm{start},q)$ holds, where $q(K,B',f',T',w')$ says: for every enlargement $B'$ of $B$, frame $f=f'$,
starting witness $w$ over $(B,f,T)$ with the column pin, selection, $J$-datum and field datum of $w'$, every
state $x$ over a pre-datum $P$ at $w$ with a third intermediate level, every pre-datum $P'$ at $w'$ with
$\mathrm{Ext}(P.z,P'.z)$ and every state $x'$ over $P'$ with an intermediate level and the same Galois data as $x$,
there is a state $x''$ over $P'$ with a second intermediate level such that $D_x$ is carried to $D_{x''}$ and
$\mathrm{List}_c(x)\Rightarrow\mathrm{List}_c(x'')$ for all $c$.
\formallabel{GlobalTheory.T.levelI_same_gal_transfer_to_levelB1_list}
\formallabel{GlobalTheory.HeightCounting.Up.exists_levelB1_carried_of_levelI_same_gal}
\end{lemma}

Here $D$ is \emph{carried} to $D'$ if they have the same export data and the embeddings of $L_0$ attached to
corresponding twisted indices agree.

\begin{proof}
Two eventual statements of the development are used: the host hypothesis holds for all stages over the
pre-base (\lref{GlobalTheory.Mem.HostFrames.hostHyp_frames_ev}), and $8fr+4\le n_V$ for these stages
(\lref{GlobalTheory.Mem.StdSize.vstd_size_ev}). Fix the data of $q$.
(1) By \lref{GlobalTheory.T.exists_eigenform_copy_of_eigensystem_alt} there are a new twisting $N$ of
$D_{x'}$, an embedding $\iota$ of the coefficient field of $D_x$ into that of $D_{x'}$ and an isomorphism of
the fields $L_0$ such that $N$ is a \emph{copy} of the twisting of $D_x$: all twisting data and expansion
characters of $D_x$ agree with those of $N$, transported along $\iota$ and the isomorphism.
(2) By \lref{GlobalTheory.T.twisted_state_transport_heq'} there are a state $x''$ over $P'$ whose
over-datum is the twist $o_{\mathrm{tw}}(x',N)$ and an intermediate level $l''$ of $x''$ whose level characters, rows,
selection map, rational datum, set $S_b$, Hecke operators, eigenvalues and characters $\theta$ are the
transports of those of the intermediate level of $x'$.
(3) By \lref{GlobalTheory.T.multiplicity_pos_of_restricted_eigenform_copy_alt} (under $\mathcal I_3$ and
$\mathcal I_{\mathrm{ser}}$, from the host hypothesis and the size bound, in the situation of (1), (2)) $x''$ has a
second intermediate level.
(4) $D_{x''}=D_{x',N}$, and $D_x$ is carried to it by the copy property.
(5) By \lref{GlobalTheory.T.hecke_eigenvalue_list_transfer_upper_alt}, in the same situation
$\mathrm{List}_c(x)\Rightarrow\mathrm{List}_c(x'')$ for all $c$.
\end{proof}

\begin{lemma}\label{pf:baselat2:tob2}
Assume $\mathcal I_3$ and $\mathcal I_{\mathrm{ser}}$. Then $\forall^{\mathrm{st}}$ of: for every datum $D_0$, if $x$ has a
second intermediate level and $D_0$ is carried to $D_x$, then there is a state $x'$ over the same pre-datum
with a third intermediate level such that $D_0$ is carried to $D_{x'}$ and $\mathrm{List}_c(x)\Rightarrow\mathrm{List}_c(x')$
for all $c$.
\formallabel{GlobalTheory.T.exists_levelB2_carried_of_levelB1}
\formallabel{GlobalTheory.HeightCounting.Up.points_preserved_under_level_transfer_alt}
\end{lemma}

\begin{proof}
This is the argument of \cref{pf:baselat2:points}, with one more property recorded. By
\cref{pf:baselat2:geolaws} the series data of a second intermediate level $l$ satisfy the geometric laws at
the generator; they provide theta data $\mathrm{th}$ with a theta tie, and
\lref{GlobalTheory.T.exists_levelB2_twist_of_thetaTie} gives a state $x'$ over the same pre-datum with a
third intermediate level, whose over-datum is the twist of that of $x$ by $\mathrm{th}$ and whose level characters
are the transported ones. Twisting by theta data does not change the export data nor the embeddings of $L_0$,
so $D_0$ is carried to $D_{x'}$. The list pin moves with
the grid transport as in \cref{pf:baselat2:points}.
\end{proof}

\begin{proposition}[Movement to the base]\label{pf:baselat2:mvbase}
$\mathcal I_3$ and $\mathcal I_{\mathrm{ser}}$ imply $\bltwoS{mv}$ for every $h_L$: in the situation of
\cref{pf:baselat2:tob1} there is a state $y$ over $P'$ with a third intermediate level such that $D_x$ is
carried to $D_y$ and $\mathrm{List}_c(x)\Rightarrow\mathrm{List}_c(y)$ for all $c$.
\formallabel{GlobalTheory.T.moveToBase_of_inputs_serN}
\formallabel{GlobalTheory.HeightCounting.Up.multiplicity_bound_of_local_components_alt}
\end{proposition}

\begin{proof}
By \cref{pf:baselat2:tob1,pf:baselat2:tob2} and the closure properties (the conclusion of
\cref{pf:baselat2:tob2} is of the form $\forall^{\mathrm{st}}(\cdot)$, hence of the form $\mathcal W(\mathrm{start},\cdot)$, and
$w'$ is a starting witness), it suffices to argue pointwise: given the data, \cref{pf:baselat2:tob1} gives
$x''$ with a second intermediate level, $D_x$ carried to $D_{x''}$ and the list transfer, and
\cref{pf:baselat2:tob2}, applied with $D_0=D_x$, gives $y$ with a third intermediate level, $D_x$ carried to
$D_y$ and $\mathrm{List}_c(x'')\Rightarrow\mathrm{List}_c(y)$. Compose the two list implications.
\end{proof}

% ---------------------------------------------------------------------------
\subsection{Theta series and Hecke operators on real sides}\label{sec:pf:baselat2:theta}

This subsection concerns twisted theta series at the series side of a datum, and Hecke operators on the forms
of a real side (\cref{pf:baselat2:hecke}); the results are used in the construction of the moduli points
(\cref{pf:baselat2:points}) and in the export statement of the section of
\lref{GlobalTheory.Mem.Top.Bl.exists_baseLattice_of_pair_alt}. We use the following objects from the
development. $D$ is a global datum, $X$ a \emph{real side} of $D$,
$E$ a family of integral Hecke forms and $Y$ a \emph{series side} over $(X,E)$.
For a twisted index $i$ the twisting provides a coefficient ring
$\OO_i$, which is a discrete valuation ring, with fraction field $K_i$ and a prime $\ell_i$ in its maximal
ideal, and for $j<m$ and $b\in\{\pm\}$ embeddings $e_{i,b}:K_i\to\C$. The series side provides: $K_i$-vector
spaces $F_{i,j,b}$ of forms with $\OO_i$-submodules $B^{\sup}_{i,j+1,b}$ (the \emph{sup lattices}); a scalar
$\nu_{i,j,b}\neq0$ and the \emph{reading} $\mathrm{rd}_{i,j,b}=\nu_{i,j,b}^{-1}\,\mathrm{ev}_K\circ e_F:F_{i,j,b}\to\C^{G}$, an
$e_{i,b}$-semilinear map into functions on an adelic group; and, for every \emph{datum} $(s,h,P)$ at
$(i,j,b)$, every finite set $Q$ of primes of $E$ in the \emph{range} of $i$ (split primes satisfying
the class condition and outside the exceptional set of $i$), every $M$ and every $a$, a function
$\xi_{i,j,b}(Q,M,h,P,a)\in\C^G$: the coefficients of the twisted universal series.

\begin{definition}[Bounded coefficients]\label{pf:baselat2:xibdd}
$\mathrm{Bd}_{\xi}(\mathrm{rd})$ means: for every twisted index $i$ with nonzero index place, every $j,b$ and every datum
$(s,h,P)$ at $(i,j,b)$ there is $c\in\N$ such that for all $Q$ in the range of $i$ whose primes do not
contain the norm product of $s$, all $M$, $a$ and all $x\in F_{i,j,b}$:
\[
  \mathrm{rd}_{i,j,b}(x)=\xi_{i,j,b}(Q,M,h,P,a)\ \Longrightarrow\ \ell_i^{\,c}\,x\in B^{\sup}_{i,j+1,b}.
\]
\formallabel{GlobalTheory.TwInputs.XiBdd}
\end{definition}

\begin{theorem}[Rationality and coefficient bound]\label{pf:baselat2:thetarat}
Assume rapid decay of cusp forms, global Artin reciprocity and the statements $\mathbf{R}_{\mathrm{ser}}$ (rationality
of the twisted series on good covers) and $\mathbf{I}_{\mathrm{ser}}$ (integrality of the twisted series and of their
descents) of the development. Let $D,X,E,Y$ be as above and assume $\mathbf{F}(X)$ ($L_0$ is Galois over
$\Q$ and stable under complex conjugation, the valuation datum holds, $L_0$ contains the ray class field of
modulus $g_0g_1p_{de}p'_{de}$ and the $g_k(g_k-1)$-th roots of unity, and $p_{de},p'_{de}$ split). Then the
rows of the series side are \emph{rational} (some normalised reading is rational) at every
twisted index with nonzero index place and all $j,b$, and $\mathrm{Bd}_{\xi}(\mathrm{rd})$ holds.
\formallabel{GlobalTheory.T.thetaSeries_rationality_and_coefficientBound}
\formallabel{GlobalTheory.Unframe.rat_and_xiBdd}
\end{theorem}

\begin{proof}
By \lref{GlobalTheory.Unframe.exists_framed} (which uses rapid decay) there are a real side $X'$ and a
series side $Y'$ over $(X',E)$ with $\mathbf{F}(X')$, with frame pins at all $j$, and such that
rationality of the rows for $Y'$ implies it for $Y$ and $\mathrm{Bd}_{\xi}$ for $Y'$ implies $\mathrm{Bd}_{\xi}$ for
$Y$. Rationality for $Y'$ is \lref{GlobalTheory.T.exists_rational_thetaSeries_of_twistSeriesSideData}
(from Artin reciprocity and $\mathbf{R}_{\mathrm{ser}}$), and the bound for $Y'$ is \cref{pf:baselat2:xibound}.
\end{proof}

\begin{theorem}[Coefficient bound]\label{pf:baselat2:xibound}
Assume global Artin reciprocity, $\mathbf{R}_{\mathrm{ser}}$ and $\mathbf{I}_{\mathrm{ser}}$. Let $D,X,E,Y$ be as above with
$\mathbf{F}(X)$ and frame pins at all $j$. Then $\mathrm{Bd}_{\xi}(\mathrm{rd})$ holds.
\formallabel{GlobalTheory.XiBound.main}
\formallabel{GlobalTheory.T.universalSeriesCoeff_boundedInSup_of_twistSeriesSideData}
\end{theorem}

The proof rests on the integral combination formula (\cref{pf:baselat2:xicomb} below) and on the following
two lemmas.

\begin{lemma}\label{pf:baselat2:powmul}
Let $\OO$ be a discrete valuation ring with fraction field $K$ and $\ell$ an element of its maximal ideal.
For every $z\in K$ there are $n\in\N$ and $o\in\OO$ with $o=\ell^nz$ in $K$.
\formallabel{GlobalTheory.XiBound.exists_pow_mul_mem}
\end{lemma}

\begin{proof}
If $\ell=0$ take $n=1$, $o=0$. Otherwise write $z=a/d$ with $a,d\in\OO$, $d\ne0$, and fix a uniformiser
$\varpi$. Then $d=v\varpi^k$ and $\ell=w\varpi^{e}$ with units $v,w$, and $e\ge1$ because $\ell$ is not a unit.
Hence $d$ divides $\ell^k=w^k\varpi^{ek}$, say $\ell^k=dt$, and $\ell^kz=ta\in\OO$.
\end{proof}

\begin{lemma}[Faithful semilinear readings]\label{pf:baselat2:faithful}
Let $\OO$ be a discrete valuation ring with fraction field $K$, $\ell$ in the maximal ideal, $e:K\to\C$ a
ring homomorphism, $F$ a $K$-vector space with an $\OO$-submodule $B$, $W$ a complex vector space and
$\rho:F\to W$ an $e$-semilinear map which is \emph{faithful}: it maps $K$-linearly independent families to
$\C$-linearly independent families. Let $\lambda\in\C^\times$. Then there is $N\in\N$ such that for all
$x\in F$: if $\rho(x')=\lambda\rho(x)$ for some $x'\in B$, then $\ell^Nx\in B$.
\end{lemma}

\begin{proof}
A faithful map is injective (a nonzero vector is a linearly independent family of length one).
Suppose first that $\lambda=e(\mu)$ for some $\mu\in K$; then $\mu\ne0$. By \cref{pf:baselat2:powmul} there are
$N$ and $o\in\OO$ with $o=\ell^N\mu^{-1}$. If $x'\in B$ and $\rho(x')=e(\mu)\rho(x)=\rho(\mu x)$, then
$x'=\mu x$ by injectivity, and $\ell^Nx=\ell^N\mu^{-1}\cdot\mu x=o\,x'\in B$.
Suppose now that $\lambda\notin e(K)$, and take $N=0$. Let $x'\in B$ with $\rho(x')=\lambda\rho(x)$. If
$\rho(x)=0$ then $x=0\in B$. Otherwise $x\neq0$, and $x,x'$ are $K$-linearly dependent (else $\rho(x)$ and
$\rho(x')=\lambda\rho(x)$ would be $\C$-linearly independent); so $x'=\mu x$ with $\mu\in K$, and then
$\lambda\rho(x)=e(\mu)\rho(x)$ with $\rho(x)\ne0$ gives $\lambda=e(\mu)$, a contradiction. So this case does
not occur, and the assertion holds.
\end{proof}

For the next two results we use more of the structure of a series side. For a twisted index $i$, $j<m$ and a
sign $b$, the series side provides a \emph{base datum} at $(i,j,b)$, consisting of a weight with signature
$(p_\tau,q_\tau)_\tau$ satisfying $p_\tau+q_\tau\le\operatorname{rank}W_j$, a complex space $\mathcal V$ of harmonic
polynomials (the \emph{value space}), and two maps into functions on a common set: the \emph{harmonic
function} $\mathrm{hf}(P)$ of a polynomial $P$, and the \emph{coefficient function} $\mathrm{cf}(F)$ of a polynomial $F$
with coefficients in $L_0$; the lattice of the base datum at the finite point determines the $\OO_i$-module
$\mathcal C$ spanned, through $e_{i,b}$, by the $\mathrm{cf}(F)$ for $F$ in the lattice. A datum $(s,h,P)$ at $(i,j,b)$
consists of a finite set $s$ of places and an element $h$ satisfying the splitting, support and finiteness
conditions of the development, and of $P\in\mathcal V$ with $\mathrm{hf}(P)\in\mathcal C$; with $(s,h,P)$, also $(s,h,0)$ is a
datum. By its definition, $\xi_{i,j,b}(Q,M,h,P,a)$ depends $\C$-linearly on $P$. We write $M_K$ for the integral
lattice of the $K_s$-form of the forms of the series side at $j+1$, and $\mathrm{ev}_W$ for the evaluation of its
elements through the series side.

For a prime $\lambda$ of the coefficient field $F_c$ of the core of $X$, a function on the adelic group of the
core is \emph{base-integral at $\lambda$} if it is a form of the core of weight $\kappa_j$ which is integral at
$\lambda$ in the sense of the development. A \emph{tower cover} of the core at a prime $\ell$ is a cover of its moduli
datum by a datum over a number field $F_C$ with an embedding $\sigma_2:F_C\to\C$, of level coprime to the
auxiliary data, as defined in the development; $\mathbf I_{\mathrm{ser}}$ states that twisted theta series on tower
covers are integral up to a bounded power of $\ell$ (precisely: for an extremal vector $u$ and a coefficient map
$c_0$ there is a phase $\phi$ such that for every prime $\ell$, local datum $s$ and finitely generated $\Z$-module
$\Lambda$ of coefficient vectors there is $c$ such that for every tower cover at $\ell$ which is large enough for
$\Lambda$, every set $S$ of primes of $F_C$ over $\ell$, every twist character $\xi$ and every level vector $(\Phi,f)$ with
$S$-integral values, and every $v\in\Lambda$, the function $g\mapsto\ell^c\phi\,\xi(\det g)\,\mathrm{series}(u,s,f,c_0v)(g)$ is an
$S$-integral form on the cover), together with a descent statement.

\begin{theorem}[Integrality on core forms]\label{pf:baselat2:xicore}
Assume global Artin reciprocity, $\mathbf{R}_{\mathrm{ser}}$ and $\mathbf{I}_{\mathrm{ser}}$. Let $D,X,E,Y$ be as above with
$\mathbf{F}(X)$, frame pins at all $j$, and $p_{de},p'_{de}$ split in $E$. Let $i$ be a twisted index with nonzero
index place, $j<m$, $b$ a sign, and let $(s,h,0)$ be a datum at $(i,j,b)$. Then there is $u\in\C^\times$ such that
for every $P\in\mathcal V$ and every polynomial $F$ with coefficients in $L_0$ with $\mathrm{hf}(P)=\mathrm{cf}(F)$ there is
$c\in\N$ with the following property: for all $Q$ in the range of $i$ whose primes do not contain the norm
product of $s$, all $M$, all $a$ and every identification of the space of the core at $\sigma_0$ with $W_j$, there
are a prime $\lambda$ of $F_c$ such that an element $x\in F_c$ is $\lambda$-integral if and only if $\sigma_0(x)=e_{i,b}(o)$
for some $o\in\OO_i$, elements $o_1,\dots,o_n\in\OO_i$ and functions $\psi_1,\dots,\psi_n$ which are base-integral at
$\lambda$, such that
\[
  \ell_i^{\,c}\,u\cdot\xi_{i,j,b}(Q,M,h,P,a)=\sum_k e_{i,b}(o_k)\,\psi_k .
\]
\formallabel{GlobalTheory.T.universal_series_coefficients_integral_core_forms}
\formallabel{GlobalTheory.XiIntSteps.xi_int_core_forms}
\formallabel{GlobalTheory.XiIntSteps.xi_int_of_ties}
\end{theorem}

\begin{remark}
Of the input $\mathbf I_{\mathrm{ser}}$ only the integrality of twisted theta series on tower covers is used below; its
descent statement is not used, and the descent is instead carried out by
\lref{GlobalTheory.XiIntSteps.intForm_descent_below}, at the cost of a bounded power of $\ell_i$.
\end{remark}

\begin{proof}[Proof of \cref{pf:baselat2:xicore}]
\emph{Preparation.} The real side $X$ provides the tie between its core and its real datum, the extremal
vector condition (\lref{GlobalTheory.IoExtremal.ioExt}), the equality of the primes of the core with those of
the cover of the real datum (\lref{GlobalTheory.CoverHalf.primes_eq_of_halfIs}), standardness and the
identification of the space of the core at $\sigma_0$ with $W_j$; Artin reciprocity gives the existence of ray
class fields (\lref{GlobalTheory.Published.GlobalArtinReciprocity.rayClassFieldSt}). Let
$\Lambda_O=e_{i,b}(\OO_i)\subseteq\C$; it is a discrete valuation ring isomorphic to $\OO_i$, $\ell_i$ is not a unit in it, and
for every $z$ in the image of the fraction field of $\OO_i$ either $z\in\Lambda_O$ or $z^{-1}\in\Lambda_O$.

\emph{Step 1: the constants.} The extremal vector $u_0$ of the real side, normalised by a suitable scalar
$c_0\neq0$, is a \emph{scale}: for a coefficient map $c_{0,L}$ on polynomials with coefficients in $L$ (with Gram data),
the image of $u_0\otimes c_{0,L}(v)$ under the extremal map is $c_0P$ whenever $P\in\mathcal V$ and $\mathrm{hf}(P)$ is the
coefficient function of $v$. Extend $c_{0,L}$ to a complex coefficient map $c_{\C}$
(\lref{GlobalTheory.ExtCoeff.exists_ext_map}). The integrality input $\mathbf I_{\mathrm{ser}}$, applied to $u_0$ and $c_{\C}$,
gives a phase $\phi$ and the integrality statement recalled above. Put $u=\phi c_0\neq0$. By
\lref{GlobalTheory.TowerPlace.placeIn_unique} there is $e\in\N$ such that for every number field $F'$ with an
embedding $\sigma:F'\to\C$ whose image lies in the fraction field of $\Lambda_O$ there is a prime $\lambda_\sigma$ of $F'$ with
$\sigma(x)\in\Lambda_O$ if and only if $x$ is $\lambda_\sigma$-integral, $\ell_i\in\lambda_\sigma$ and $\ell_i\notin\lambda_\sigma^{e+1}$. By
\lref{GlobalTheory.XiIntSteps.intForm_descent_below} (for $\ell_i$ and $e$) there is $c_1$ such that for every
uniformised pinned cover of the core (with the weight $\kappa_j$, whose last weights are positive), every prime
$\lambda'$ of the cover field with $\ell_i\in\lambda'$, $\ell_i\notin\lambda'^{\,e+1}$, and the prime $\mathfrak p$ of $F_c$ below it, every
$\lambda'$-integral form $\varphi$ on the cover which is invariant under the level of the core satisfies
$\ell_i^{c_1}\varphi=\sum_k\sigma_2(x_k)\psi_k$ with $x_k$ in the cover field, integral at $\lambda'$, and $\psi_k$ base-integral at
$\mathfrak p$; this uses that the holomorphic cusp forms of the core are spanned by sections
(\lref{GlobalTheory.OntoCore.holSub0_le_span_cuspA}, which needs $8fr+4\le n_V$).

\emph{Step 2: the exponent.} Let $P\in\mathcal V$ and $F$ with $\mathrm{hf}(P)=\mathrm{cf}(F)$. Since the weight of the base datum
is the lifted coefficient weight, $F$ is the image of a vector $v$ with coefficients in $L_0$, and by the scale
property and the injectivity of $\mathrm{hf}$ on $\mathcal V$ (\lref{GlobalTheory.T.harmonic_polynomial_eq_zero_of_function},
which uses $p_\tau+q_\tau\le\operatorname{rank}W_j$) the image of $u_0\otimes c_{0,L}(v)$ is $c_0P$. Let $c_2$ be the exponent
given by $\mathbf I_{\mathrm{ser}}$ for $\ell=\ell_i$, the local datum $h$ and the module $\Z v'$, where $v'$ is $v$ with complex
coefficients. Put $c=c_1+c_2$.

\emph{Step 3: a tower cover.} Let $Q,M,a$ be given. The admissible conjugate pairs of $(i,Q,M)$ determine a set
$H\subseteq Q$ (one place of each pair), on which the universal family of the twisting is good
(\lref{GlobalTheory.UnivGood.univOf_good}) and the base datum is admitted
(\lref{GlobalTheory.ChoiceAdm.adm_half_of_good_of_split_disjoint}). By
\lref{GlobalTheory.XiIntSteps.exists_cover_coprime} (from the existence of ray class fields and the
conditions $\mathbf F(X)$: $L_0$ Galois, values, ray class field and roots of unity, and the splitting of
$p_{de},p'_{de}$) there is a tower cover $C$ at $\ell_i$ whose level is $M_{\mathrm{l}}\,p_d$ with $p_d$ divisible by the
places of $H$ and of the datum, coprime to $QS\cdot q\cdot N(\mathfrak c)$, whose field contains the values of the relevant
characters and the ray class field, and whose embedding $\sigma_2$ extends $\sigma_0$. Let $\lambda'$ be the prime of the cover field attached to $\sigma_2$
by Step 1 and $\lambda$ the prime of $F_c$ attached to $\sigma_0$; then $\lambda=\lambda'\cap F_c$, and $\lambda$ has the integrality
property of the theorem.

\emph{Step 4: integral series on the cover.} By \lref{GlobalTheory.XiIntSteps.theta_twist} there are a twist
character $\xi_C$ for the cover (built from the character $\delta$ of the real datum at $\ell_i$, which is finite,
tame and unramified outside $p_d$, and whose values lie in the cover field) and a root of unity $z$ in the cover
field such that the twisted theta functions with level vectors equal $\xi_C(\det g)$ times the theta series. By
\lref{GlobalTheory.XiIntSteps.isLevelFn_datumFn}, for each fibre index $a'$ the level function of the datum is
a level vector for the level of $C$ with values in $\sigma_2(F_C)\cap\Lambda_O$, i.e.\ integral at $\lambda'$; the field of $C$ is
large enough for $\Z v'$ (it contains the values of the character quotients,
\lref{GlobalTheory.Member.character_quotient_values_in_embedding_field.small_quot_values}, and the ray
class field). Hence $\mathbf I_{\mathrm{ser}}$ gives a $\lambda'$-integral form $\varphi_{a'}$ on $C$ with
$\varphi_{a'}=\ell_i^{c_2}u\cdot\mathrm{series}(\dots,a')$ (after multiplication by the $\lambda'$-unit $z$). Summing over the
fibres (\lref{GlobalTheory.XiIntSteps.exists_sum_fibres}) gives a $\lambda'$-integral form $\Phi$ on $C$ with
$\Phi=\ell_i^{c_2}u\cdot\xi_{i,j,b}(Q,M,h,P,a)$.

\emph{Step 5: descent.} $\Phi$ is invariant under the level of the core: for $k$ in the level,
$\xi_{i,j,b}(\dots)(gk)=\xi_{i,j,b}(\dots)(g)$, by the compatibility of the level characters with the embedding
(\lref{GlobalTheory.T.level_character_embedding_compat}) and the right translation property of the
coefficients of the universal series. By Step 1, $\ell_i^{c_1}\Phi=\sum_k\sigma_2(x_k)\psi_k$ with $x_k$ integral at $\lambda'$
and $\psi_k$ base-integral at $\lambda$. By the choice of $\lambda'$, $\sigma_2(x_k)\in\Lambda_O$, i.e.\ $\sigma_2(x_k)=e_{i,b}(o_k)$ with
$o_k\in\OO_i$, and therefore $\ell_i^{c_1+c_2}u\cdot\xi_{i,j,b}(Q,M,h,P,a)=\sum_ke_{i,b}(o_k)\psi_k$.
\end{proof}

\begin{proposition}[Integral combination formula]\label{pf:baselat2:xicomb}
Assume global Artin reciprocity, $\mathbf{R}_{\mathrm{ser}}$ and $\mathbf{I}_{\mathrm{ser}}$. Let $D,X,E,Y$ be as above with
$\mathbf{F}(X)$ and frame pins at all $j$. Let $i$ be a twisted index with nonzero index place, $j<m$, $b$ a sign and
$(s,h,P)$ a datum at $(i,j,b)$. Then there are $u\in\C^\times$ and $c\in\N$ such that for all $Q$ in the range of $i$
whose primes do not contain the norm product of $s$, all $M$ and all $a$ there are $o_1,\dots,o_n\in\OO_i$ and
$m_1,\dots,m_n\in M_K$ with
\begin{equation}\label{eq:pf:baselat2:xicomb}
  \ell_i^{\,c}\,u\cdot\xi_{i,j,b}(Q,M,h,P,a)=\sum_k e_{i,b}(o_k)\,\mathrm{ev}_W(m_k).
\end{equation}
\formallabel{GlobalTheory.XiIntComb.xi_int_comb}
\formallabel{GlobalTheory.T.thetaSeries_integral_core_supLattice_combination}
\formallabel{GlobalTheory.XiIntCore.xi_int_core}
\end{proposition}

\begin{proof}
For $u\in\C^\times$ say that $P'\in\mathcal V$ is \emph{good} if there is $c$ such that the conclusion holds for $P'$ in place of
$P$. The condition $\mathbf F(X)$ contains the splitting of $p_{de},p'_{de}$, and $(s,h,0)$ is a datum; let $u$ be given
by \cref{pf:baselat2:xicore} for $(s,h,0)$.

(1) If $P'\in\mathcal V$ and $\mathrm{hf}(P')=\mathrm{cf}(F)$ for a polynomial $F$ with coefficients in $L_0$, then $P'$ is good.
Indeed, \cref{pf:baselat2:xicore} gives $c$ and, for given $Q,M,a$, a prime $\lambda$, elements $o_k\in\OO_i$ and
base-integral functions $\psi_k$ at $\lambda$ with $\ell_i^cu\,\xi(Q,M,h,P',a)=\sum_ke_{i,b}(o_k)\psi_k$; and by
\lref{GlobalTheory.XiIntCore.supLattice_comb_of_baseInt} every such combination of base-integral functions
equals $\sum_ke_{i,b}(o'_k)\,\mathrm{ev}_W(m_k)$ for suitable $o'_k\in\OO_i$ and $m_k\in M_K$.

(2) The good polynomials form an $\OO_i$-submodule of $\mathcal V$ (with $\OO_i$ acting through $e_{i,b}$): $0$ is good (with
the empty combination), and since $\xi$ is $\C$-linear in $P'$: if $P_1,P_2$ are good with exponents $c_1,c_2$, then
$P_1+P_2$ is good with exponent $c_1+c_2$ (multiply the two combinations by $\ell_i^{c_2}$ and $\ell_i^{c_1}$, which lie in
$\OO_i$), and $e_{i,b}(o)P_1$ is good with exponent $c_1$ (multiply the coefficients $o_k$ by $o$).

(3) By \lref{GlobalTheory.T.harmonic_polys_reach_small_field_coefficient_functions}, for every $F$ in the
lattice of the base datum there is $P'\in\mathcal V$ with $\mathrm{hf}(P')=\mathrm{cf}(F)$; such $P'$ is good by (1). Since $\mathrm{hf}$ is linear,
induction over the $\OO_i$-span shows, with (2): every $\varphi\in\mathcal C$ is $\mathrm{hf}(P')$ for some good $P'\in\mathcal V$.

(4) Apply (3) to $\varphi=\mathrm{hf}(P)\in\mathcal C$: $\mathrm{hf}(P)=\mathrm{hf}(P')$ with $P'$ good. Since $\mathrm{hf}$ is injective on $\mathcal V$
(\lref{GlobalTheory.T.harmonic_polynomial_eq_zero_of_function}, using $p_\tau+q_\tau\le\operatorname{rank}W_j$) and
$P,P'\in\mathcal V$, $P=P'$ is good.
\end{proof}

\begin{proof}[Proof of \cref{pf:baselat2:xibound}]
Let $i$ have nonzero index place and let $(s,h,P)$ be a datum at $(i,j,b)$. By \cref{pf:baselat2:xicomb} there are
$u\in\C^\times$ and $c\in\N$ such that for all admissible $Q,M$ and all $a$ there are $o_1,\dots,o_n\in\OO_i$ and
$m_1,\dots,m_n\in M_K$ satisfying \eqref{eq:pf:baselat2:xicomb}.
The reading $\mathrm{rd}_{i,j,b}$ is faithful by \lref{GlobalTheory.SeriesTie.realRd_faithful} (whose
hypotheses, $L$-uniformity, highest weights, the extremal vector and the $K_s$-form property, are part of
$X$, $Y$ and \lref{GlobalTheory.T.rational_datum_small_field_form_hecke_scalar_unit}). Put
$\lambda=\nu_{i,j,b}^{-1}\ell_i^{\,c}u\neq0$ and let $N$ be given by \cref{pf:baselat2:faithful} for
$(\OO_i,K_i,\ell_i,e_{i,b},F_{i,j,b},B^{\sup}_{i,j+1,b},\mathrm{rd}_{i,j,b},\lambda)$. Given $x$ with
$\mathrm{rd}(x)=\xi(Q,M,h,P,a)$, let $y=\sum_k o_k(1\otimes m_k)$; it lies in the sup lattice of the tensor model
of the forms, which is the $\OO_i$-span of $1\otimes M_K$, and $\mathrm{ev}_K(y)=\sum_ke(o_k)\mathrm{ev}_W(m_k)$. The series
side identifies $B^{\sup}_{i,j+1,b}$ with the sup lattice of the tensor model through $e_F$, so
$x'=e_F^{-1}(y)\in B^{\sup}_{i,j+1,b}$, and by \eqref{eq:pf:baselat2:xicomb}
\[
  \mathrm{rd}(x')=\nu^{-1}\mathrm{ev}_K(y)=\nu^{-1}\ell_i^{\,c}u\,\xi(Q,M,h,P,a)=\lambda\,\mathrm{rd}(x).
\]
Hence $\ell_i^Nx\in B^{\sup}_{i,j+1,b}$, with $N$ depending only on $(i,j,b,s,h,P)$.
\end{proof}

\subsubsection{Hecke operators and sup lattices}

The following result is used in the section of
\lref{GlobalTheory.Mem.Top.Bl.exists_baseLattice_of_pair_alt}. Let $X$ be a real side of $D$ and $i<m$. The
core of $X$ has a coefficient field $F_c$ (a number field mapping to $L$), and the real datum of $X$ provides a
\emph{form datum} at $i$: the $L$-vector space $W_L=L^{n_W(i)}$, an $F_c$-subspace $W_K\subseteq W_L$ of rational forms,
an action $\rho$ of a finite group $G$ on $W_L$, and for every prime $\lambda$ of $F_c$ a set $Y_\lambda$ of points and norms
$\|f\|_{\lambda,y}\in\R$ for $f\in W_K$, $y\in Y_\lambda$. Put $\mathrm{Ball}_\lambda=\{f\in W_K:\|f\|_{\lambda,y}\le1\text{ for all }y\in Y_\lambda\}$.
For $e_C:G\to L$ let $e(e_C)=\sum_ge_C(g)\rho(g)$; for a subgroup $J\le G$ and a character $\vartheta:J\to L^\times$ let
$W^{J,\vartheta}=\{x:\rho(p)x=\vartheta(p)x\text{ for }p\in J\}$, $W^{e_C}=e(e_C)(W_L)$ and $W^{J,\vartheta,e_C}=W^{J,\vartheta}\cap W^{e_C}$.
For a field $K'$ over $F_c$ and a valuation ring $\OO'\subseteq K'$ \emph{over $\lambda$} (an element of $F_c$ lies in $\OO'$ if
and only if it is $\lambda$-integral) let $\Lambda_{K'}\subseteq K'\otimes_{F_c}W_K$ be the $\OO'$-span of the $1\otimes f$,
$f\in\mathrm{Ball}_\lambda$, and for an $L$-subspace $V$ let $\Lambda_{K'}(V)=\Lambda_{K'}\cap\bigl(K'\otimes_{F_c}(V\cap W_K)\bigr)$; for a field
$K'$ over $L$ define $\Lambda^L_{K'}\subseteq K'\otimes_LW_L$ and $\Lambda^L_{K'}(V)=\Lambda^L_{K'}\cap(K'\otimes_LV)$ in the same way.

\begin{definition}[Keeping the sup lattice]\label{pf:baselat2:keeps}
An $L$-linear endomorphism $h$ of $W_L$ with $h(W_K)\subseteq W_K$ \emph{keeps the $\lambda$-sup lattice} if:
(sup) $h(\mathrm{Ball}_\lambda)\subseteq\mathrm{Ball}_\lambda$; ($\rho$) $h\rho(g)=\rho(g)h$ for all $g\in G$; (e) $h$ commutes with
every $e(e_C)$; (type) and (good) $h$ preserves every $W^{J,\vartheta}$ and every $W^{e_C}$; (K) for every $K'$ over
$F_c$ and every valuation ring $\OO'$ of $K'$ over $\lambda$, the base change of $h|_{W_K}$ preserves $\Lambda_{K'}$,
$\Lambda_{K'}(W^{J,\vartheta})$ and $\Lambda_{K'}(W^{J,\vartheta,e_C})$ for all $J,\vartheta,e_C$; (L) the same for the base change of $h$
to fields $K'$ over $L$ and the modules $\Lambda^L_{K'}$, $\Lambda^L_{K'}(W^{J,\vartheta})$, $\Lambda^L_{K'}(W^{J,\vartheta,e_C})$.
\formallabel{GlobalTheory.IndefiniteRowSideFullHecke.LatticesOverBalls.Keeps}
\end{definition}

\begin{theorem}[Hecke operators keep sup lattices]\label{pf:baselat2:hecke}
Assume rapid decay of cusp forms, the integrality of Hecke values, the transform statement for sections and
cusp forms, and the reading tie for Hecke operators, in the forms stated in the development. Let $X$ be a
real side of $D$ and $i<m$. Let $\mathfrak P$ be a prime of $E$ which splits off the finite set $S_b$ of the real datum
and does not contain $g_0g_1$, let $k\le r$, let $\alpha$ be a Hecke element at $(\mathfrak P,k)$ for the group of the
$i$-th row and $R_c$ a set of representatives of the double coset of $\alpha$ for the level group of the real datum.
Let $h$ be an $L$-linear endomorphism of $W_L$ with $h(W_K)\subseteq W_K$ which is the Hecke operator of $R_c$ scaled
by the Hecke constant $c(\mathfrak P,k)$ of the real datum: $\mathrm{ev}(j(hx))=c(\mathfrak P,k)\,T_{R_c}(\mathrm{ev}(j(x)))$ for all
$x\in W_L$, where $j$ is the reading of the forms at the base embedding, $\mathrm{ev}$ the evaluation and $T_{R_c}$ the
Hecke operator of $R_c$ on functions. Then $h$ keeps the $\lambda$-sup lattice for every prime $\lambda$ of $F_c$ which
does not lie above the residue characteristic of $\mathfrak P$.
\formallabel{GlobalTheory.T.heckeOperator_preserves_supLattice_realSide}
\formallabel{GlobalTheory.HeckeSupReal.hkSupAt_realSide}
\formallabel{GlobalTheory.HeckeSup.hkSupAt_of_clauses}
\formallabel{GlobalTheory.HeckeSup.hkSupAt_of_readsCusp}
\end{theorem}

\begin{proof}
\emph{Step 1: it suffices to prove (sup) and ($\rho$).} Assume them. (e) holds since $e(e_C)$ is an $L$-linear
combination of the $\rho(g)$. (type): for $x\in W^{J,\vartheta}$ and $p\in J$, $\rho(p)h(x)=h(\rho(p)x)=\vartheta(p)h(x)$. (good):
$h(e(e_C)y)=e(e_C)h(y)$. Hence $h$ also preserves the $W^{J,\vartheta,e_C}$. (K): the base change $h_{K'}$ of $h|_{W_K}$ is
$\OO'$-linear and maps a generator $1\otimes f$ of $\Lambda_{K'}$ to $1\otimes h(f)$, which is again a generator by (sup); so
$h_{K'}(\Lambda_{K'})\subseteq\Lambda_{K'}$. If $V$ is $h$-stable then so is $V\cap W_K$, hence $K'\otimes_{F_c}(V\cap W_K)$ is $h_{K'}$-stable,
and $\Lambda_{K'}(V)$ is $h_{K'}$-stable; apply this to $V=W^{J,\vartheta}$ and $V=W^{J,\vartheta,e_C}$. (L) is proved in the same way.

\emph{Step 2: ($\rho$).} This is \lref{GlobalTheory.HeckeSup.comm_rho_cl} (under the frame, family and model
conditions of the real datum, if the reading $j$ at the base embedding is injective and $\rho$ is the action
transported through $j$, every Hecke operator of Hecke data as in the theorem commutes with all $\rho(g)$);
the injectivity of $j$ and the description of $\rho$ are part of the reading at the base point provided by the
real side (the real datum is $L$-uniform, so $j$ maps linearly independent families to linearly independent
families and is injective, and $\rho$ is the action transported through $j$).

\emph{Step 3: the reading property.} The real side satisfies the hypotheses of
\lref{GlobalTheory.HeckeSup.readsCusp_of_clauses}: the frame, family and model conditions, the frame tie,
the reading at the base point, the pinning of its constants (a repackaging of the one-map rows,
\lref{GlobalTheory.OneMapTies.oneMapRow0_of_ties_ofFive}, and of the constant tie of the real side, using
the extremal vector condition \lref{GlobalTheory.IoExtremal.ioExt}), the spanning statement
\lref{GlobalTheory.AdSpan.spanAt} (which uses rapid decay), the condition on orthonormal data
(\lref{GlobalTheory.AdTh.ioExt_oneMap_ofFive}), and the reading tie for Hecke operators for all modifications
of the constants of the core. That statement gives the \emph{reading property}: if cusp forms $f_0,g_0$ of weight
$\kappa_i$ satisfy the geometric Hecke relation of $(\mathfrak P,p,k)$ with the normalised scalar, then their adelic
functions satisfy $\varphi_{g_0}(x)=c(\mathfrak P,k)\sum_{a\in R_c}\varphi_{f_0}(xa)$; consequently, if $h$ is as in the theorem,
then $h(e(f_0))=e(g_0)$, where $e$ denotes the map from cusp forms of weight $\kappa_i$ to $W_K$
(\lref{GlobalTheory.HeckeSup.apply_eq_of_read_cl}, using the injectivity of $j$).

\emph{Step 4: (sup).} Let $p$ be the rational prime below $\mathfrak P$, so that $N\mathfrak P=\#k_{\mathfrak P}=p^m$ for some $m$. Since
$g_0g_1\notin\mathfrak P$, $\mathfrak P$ does not divide the level of the core, and since $\mathfrak P$ splits off $S_b$ it is split; so $\mathfrak P$ is
a split Hecke prime off the level. Let $f\in\mathrm{Ball}_\lambda$. The map $e$ is surjective, so $f=e(f_0)$ for a cusp form
$f_0$. By \lref{GlobalTheory.HeckeSup.exists_transform} (from the transform statement and the Hecke point
laws \lref{GlobalTheory.HeckePointLaw.heckePointLaws}) there is a cusp form $g_0$ of the same weight such
that $(f_0,g_0)$ satisfy the geometric Hecke relation of $(\mathfrak P,p,k)$ with the scalar
\[
  s=c_q\,(N\mathfrak P)^{-\lfloor kn_V/2\rfloor},\qquad c_q=(\#k_{\mathfrak P})^{-\lfloor kn_V/2\rfloor}:
\]
at every geometric point $y$ the value of $g_0$ is $s$ times the Hecke sum of the values of $f_0$ at the
quotients of $y$ by the subgroups in the Hecke kernel at $y$. By Step 3, $h(f)=e(g_0)$. Since $s=p^{-2m\lfloor kn_V/2\rfloor}$
and $\lambda$ does not lie above $p$, $s$ is a $\lambda$-adic unit. The norms $\|\cdot\|_{\lambda,y}$ are at most $1$ exactly when the
corresponding point values are $\lambda$-integral (frame tie). By \lref{GlobalTheory.HeckeSup.transform_integral}
(from the integrality of Hecke values: each term of the Hecke sum of a form with integral point values is
integral over the valuation ring of the point), the values of $g_0$ at all points over $\lambda$ are therefore
integral. Hence $h(f)\in\mathrm{Ball}_\lambda$.
\end{proof}


% ---------------------------------------------------------------------------
\subsection{Differents and traces}\label{sec:pf:baselat2:diff}

For a finite extension of number fields $L/K$ write $\mathfrak D_{L/K}$ for the different of $\OO_L/\OO_K$.

\begin{theorem}[Bounded order of the relative different]\label{pf:baselat2:different}
For every prime $\ell$ and every $e\in\N$ there is $c\in\N$ (one can take $c=e^2$) such that for all number
fields $K\subseteq L$, all primes $\mathfrak p$ of $K$ and $\lambda$ of $L$ with $\lambda\cap\OO_K=\mathfrak p$,
$\ell\in\lambda$ and $\ell\notin\lambda^{e+1}$, the ideal $\lambda^{c+1}$ does not divide $\mathfrak D_{L/K}$.
\formallabel{GlobalTheory.T.relative_different_order_bounded_absolute_ramification}
\formallabel{GlobalTheory.TowerDifferent.relDifferent_order_bounded}
\end{theorem}

The hypothesis says that the absolute ramification index $e(\lambda/\ell)$ is at most $e$. Classically
$v_\lambda(\mathfrak D_{L/\Q})\le e(\lambda/\ell)-1+v_\lambda(e(\lambda/\ell))\le e^2-1$
\cite[Ch.~III, \S6, Rem.~1 after Prop.~13]{Serre1979}, and the theorem follows from the transitivity of the
different. The formal proof is different and self-contained; we give it.

\begin{lemma}[Truncated local rings]\label{pf:baselat2:quotpow}
Let $L$ be a number field, $\lambda$ a maximal ideal of $\OO_L$ and $J$ an ideal with $\lambda+J=\OO_L$. Let
$\ell$ be prime, $e,b\ge1$, and assume $\ell\OO_L=\lambda^eJ$. Then $B=\OO_L/\lambda^{eb}$ is a finite ring of
characteristic $\ell^b$, and there is $\pi\in B$ such that: $x\in B$ is a non-unit if and only if $\pi\mid x$;
$\ell B=\pi^eB$; and $|B|=|B/\pi B|^{be}$.
\formallabel{GlobalTheory.T.quotient_pow_finite_uniformizer_card_eq}
\end{lemma}

\begin{proof}
$B$ is finite since $\lambda^{eb}\ne0$. Its only prime ideal is $\lambda/\lambda^{eb}$, so $B$ is local. Choose
$\varpi\in\lambda\smallsetminus\lambda^2$ and let $\pi$ be its image. Since $\lambda/\lambda^2$ is one-dimensional over
$\OO_L/\lambda$, $\lambda=\varpi\OO_L+\lambda^2$; multiplying by $\lambda^{k-1}$ and substituting gives
$\lambda=\varpi\OO_L+\lambda^{k}$ for all $k\ge1$, in particular $\lambda=\varpi\OO_L+\lambda^{eb}$, so the maximal ideal of $B$ is
$\pi B$, which is the set of non-units. The image of $J$ in $B$ is the unit ideal (as $J+\lambda=\OO_L$, $J$ is not
contained in the unique maximal ideal), so $\ell B=\lambda^eB=\pi^eB$. Since $\ell^b\OO_L=\lambda^{eb}J^b\subseteq\lambda^{eb}$
and $v_\lambda(\ell^{b-1})=e(b-1)<eb$, the characteristic of $B$ is $\ell^b$. Finally $\lambda^k/\lambda^{k+1}$ is a
one-dimensional vector space over $\OO_L/\lambda=B/\pi B$, so $|B|=|\OO_L/\lambda|^{eb}$.
\end{proof}

\begin{lemma}[Non-vanishing of the trace]\label{pf:baselat2:trace}
Let $\ell$ be prime, $b,e\in\N$ with $e\ge1$, and let $B$ be a finite commutative $\Z/\ell^b$-algebra with an
element $\pi$ such that: the non-units of $B$ are the multiples of $\pi$; $\ell B=\pi^eB$; and
$|B|=|B/\pi B|^{be}$. If $\ell^b\nmid e$, there is $c\in B$ with $\Tr_{B/(\Z/\ell^b)}(c)\ne0$.
\formallabel{GlobalTheory.T.exists_trace_ne_zero_finite_local_zmod_prime_pow}
\formallabel{GlobalTheory.LocalTraceNonzero.exists_trace_ne_zero}
\end{lemma}

\begin{proof}
We have $b\ge1$, since $\ell^0=1$ divides $e$. $B$ is nonzero (otherwise $0=1$ would be a unit, but $0$ is a
multiple of $\pi$). The non-units are closed under addition, so $B$ is local with maximal ideal $\pi B$; and
$\ell\in\pi^eB\subseteq\pi B$, so the residue field $k=B/\pi B$ is a finite field of characteristic $\ell$.
\begin{enumerate}[label=(\alph*)]
\item Write $k=\F_\ell(\alpha)$ where $\alpha$ is a root of a monic irreducible $p\in\F_\ell[X]$ with $p'(\alpha)\ne0$, so
  that $|k|=\ell^{\deg p}$ (finite fields are simple separable extensions of $\F_\ell$).
\item Let $h\in(\Z/\ell^b)[X]$ be a monic lift of $p$ of the same degree. A finite local ring is Henselian
  (its maximal ideal is nilpotent, so it is complete). Since $\alpha$
  is a simple root of the reduction of $h$, there is $\omega\in B$ with $h(\omega)=0$ and $\omega\equiv\alpha$ modulo
  $\pi$.
\item Let $C=(\Z/\ell^b)[X]/(h)$ and $\varphi:C\to B$ the homomorphism $X\mapsto\omega$; $B$ is a $C$-module via $\varphi$.
  Every $x\in B$ is congruent modulo $\pi$ to some $\varphi(y)$, because $k=\F_\ell[\alpha]$. Moreover
  $\varphi(\ell)^b=\ell^b=0$ in $B$ and $\pi^e\in\varphi(\ell)B$. Let $M=\{\sum_{j<e}\varphi(c_j)\pi^j:\ c_j\in C\}$, a
  $C$-submodule. Applying the residue approximation $e$ times gives $B=M+\pi^eB\subseteq M+\varphi(\ell)B$; if
  $B=M+\varphi(\ell)^tB$, then $\varphi(\ell)^tB\subseteq\varphi(\ell)^tM+\varphi(\ell)^{t+1}B\subseteq M+\varphi(\ell)^{t+1}B$, so
  $B=M+\varphi(\ell)^{t+1}B$. For $t=b$ this gives $B=M$.
\item The $C$-linear map $T:C^e\to B$, $(c_j)\mapsto\sum_j\varphi(c_j)\pi^j$, is therefore surjective. As $C$ is free of
  rank $\deg p$ over $\Z/\ell^b$, $|C^e|=\ell^{be\deg p}=|k|^{be}=|B|$, so $T$ is bijective and $1,\pi,\dots,\pi^{e-1}$ is a
  $C$-basis of $B$.
\item By transitivity of traces, for $y\in C$:
  $\Tr_{B/(\Z/\ell^b)}(\varphi(y))=\Tr_{C/(\Z/\ell^b)}(\Tr_{B/C}(\varphi(y)))=\Tr_{C/(\Z/\ell^b)}(e\,y)=e\,\Tr_{C/(\Z/\ell^b)}(y)$.
\item Reduction modulo $\ell$ maps $C$ onto $C/\ell C=\F_\ell[X]/(p)$, a finite separable field extension of $\F_\ell$,
  and commutes with traces (the trace of an element of a free algebra is the trace of its multiplication
  matrix, which is compatible with base change). The trace of a finite separable extension is not
  identically zero (by Dedekind's independence of characters); lifting an element of nonzero
  trace gives $y\in C$ with $\Tr_C(y)\not\equiv0\pmod\ell$, so $\Tr_C(y)$ is a unit of $\Z/\ell^b$.
\end{enumerate}
Then $\Tr_{B/(\Z/\ell^b)}(\varphi(y))=e\cdot\Tr_C(y)\ne0$, since $\ell^b\nmid e$ and $\Tr_C(y)$ is a unit.
\end{proof}

\begin{lemma}[The absolute different]\label{pf:baselat2:absdiff}
Let $L$ be a number field, $\ell$ a prime, $e\in\N$ and $\lambda$ a nonzero prime of $\OO_L$ with $\ell\in\lambda$ and
$\ell\notin\lambda^{e+1}$. Then $\lambda^{e^2+1}$ does not divide $\mathfrak D_{L/\Q}$.
\formallabel{GlobalTheory.TowerDifferent.not_pow_dvd_differentIdeal_int}
\end{lemma}

\begin{proof}
By unique factorisation of ideals in $\OO_L$ there are
$1\le e'\le e$ and an ideal $J$ with $\lambda+J=\OO_L$ and $\ell\OO_L=\lambda^{e'}J$. Apply \cref{pf:baselat2:quotpow}
with $b=e'$: $B=\OO_L/\lambda^{e'^2}$ satisfies the hypotheses of \cref{pf:baselat2:trace} with $(b,e)=(e',e')$, and
$\ell^{e'}\nmid e'$ because $\ell^{e'}>e'$. So there is $c_0\in\OO_L$ whose image in $B$ has nonzero trace over
$\Z/\ell^{e'}$. Since $\lambda^{e'^2}+J^{e'}=\OO_L$, write $1=i+q$ with $i\in\lambda^{e'^2}$, $q\in J^{e'}$, and put
$x=qc_0\in J^{e'}$; then $x\equiv c_0$ modulo $\lambda^{e'^2}$. As $\lambda^{e'^2}J^{e'}=\ell^{e'}\OO_L$ and the two
factors are coprime, the Chinese remainder theorem gives $\OO_L/\ell^{e'}\cong B\times\OO_L/J^{e'}$, in which $x$
maps to $(\bar c_0,0)$; hence $\Tr_{\OO_L/\Z}(x)\bmod\ell^{e'}=\Tr_{B/(\Z/\ell^{e'})}(\bar c_0)\ne0$
(the trace of $x$ over $\Z$ reduces modulo $\ell^{e'}$ to the trace of $x$ on $\OO_L/\ell^{e'}$, as $\OO_L$ is free
over $\Z$, and the latter is the sum of the traces on the two factors). Now if $\lambda^{e'^2}$ divided $\mathfrak D_{L/\Q}$, then
$\lambda^{-e'^2}=\ell^{-e'}J^{e'}\subseteq\mathfrak D_{L/\Q}^{-1}$, the trace dual of $\OO_L$, and
$\Tr(\ell^{-e'}x)\in\Z$, i.e.\ $\Tr(x)\in\ell^{e'}\Z$, a contradiction. Since $e'^2\le e^2$,
$\lambda^{e^2+1}\nmid\mathfrak D_{L/\Q}$.
\end{proof}

\begin{proof}[Proof of \cref{pf:baselat2:different}]
Take $c=e^2$. By \cref{pf:baselat2:absdiff}, $\lambda^{e^2+1}\nmid\mathfrak D_{L/\Q}$. The different $\mathfrak D_{L/K}$
divides $\mathfrak D_{L/\Q}=\mathfrak D_{L/K}\,\mathfrak D_{K/\Q}\OO_L$ (transitivity of the different, \cite[Ch.~III, \S4]{Serre1979}), so $\lambda^{e^2+1}\nmid\mathfrak D_{L/K}$.
\end{proof}

% ---------------------------------------------------------------------------
\subsection{Congruence lengths and lattices}\label{sec:pf:baselat2:len}

The following three results are used in the proof of the equality of local and rational congruence
lengths (\cref{pf:baselat2:etasum}, at the end of this subsection), which enters the eta equality
(\cref{pf:baselat2:etaj}) through \lref{GlobalTheory.T.etaJ_local_compatibility_alt}. Lengths of modules
take values in $\N\cup\{\infty\}$.

For an $\OO$-module $V$ and submodules $L,G,B$ of $V$ put
\[
  \mathrm{cl}(L;G,B)=\len_{\OO}\Bigl(\bigl(G\cap(L+B)\bigr)\big/\bigl(L\cap G\bigr)\Bigr),
\]
the \emph{congruence length} of $L$ with respect to $(G,B)$ (note $L\cap G\subseteq G\cap(L+B)$).

\begin{proposition}[Additivity of congruence lengths]\label{pf:baselat2:cl}
Let $\OO$ be a commutative ring, $V$ an $\OO$-module and $(L_i,G_i,B_i)_{i\in F}$ a finite family of triples of
submodules such that the family $(G_i+B_i)_{i\in F}$ is independent and $L_i\subseteq G_i+B_i$ for all $i$. Then
\[
  \mathrm{cl}\Bigl(\sum_iL_i;\ \sum_iG_i,\ \sum_iB_i\Bigr)=\sum_{i\in F}\mathrm{cl}(L_i;G_i,B_i).
\]
\formallabel{GlobalTheory.T.congruenceLength_finset_sup_eq_sum}
\formallabel{GlobalTheory.congruence_length_finset_sup}
\end{proposition}

\begin{remark}
The formal statement also assumes $G_i\cap B_i=0$ for all $i$; this hypothesis is not used in the proof.
\end{remark}

\begin{proof}
We use three facts. (a) If $P_1\cap P_2=0$, $x_1,y_1\in P_1$, $x_2,y_2\in P_2$ and $x_1+x_2=y_1+y_2$, then $x_1=y_1$ and
$x_2=y_2$ (since $x_1-y_1=y_2-x_2\in P_1\cap P_2$).
(b) If $(G_1+B_1)\cap(G_2+B_2)=0$ and $L_k\subseteq G_k+B_k$ ($k=1,2$), then
\begin{gather*}
  (G_1+G_2)\cap\bigl((L_1+L_2)+(B_1+B_2)\bigr)=G_1\cap(L_1+B_1)+G_2\cap(L_2+B_2),\\
  (L_1+L_2)\cap(G_1+G_2)=L_1\cap G_1+L_2\cap G_2 .
\end{gather*}
For the first, ``$\supseteq$'' is clear; if $g_1+g_2=(l_1+b_1)+(l_2+b_2)$ with $g_k\in G_k$, $l_k\in L_k$, $b_k\in B_k$,
then $g_k$ and $l_k+b_k$ lie in $G_k+B_k$, so (a) gives $g_k=l_k+b_k\in G_k\cap(L_k+B_k)$. The second is proved
in the same way.
(c) If $S_k\subseteq T_k$ ($k=1,2$) and $T_1\cap T_2=0$, then
$\len\bigl((T_1+T_2)/(S_1+S_2)\bigr)=\len(T_1/S_1)+\len(T_2/S_2)$: the map
$T_1\times T_2\to(T_1+T_2)/(S_1+S_2)$, $(a,b)\mapsto a+b$, is surjective with kernel $S_1\times S_2$ by (a) (as
$S_1+S_2\ni a+b$ forces $a\in S_1$, $b\in S_2$), so the quotient is isomorphic to $T_1/S_1\times T_2/S_2$, and length is
additive on products.

Now induct on $F$. For $F=\emptyset$ both sides vanish. For $F'=F\cup\{j\}$, $j\notin F$, put $G=\sum_{i\in F}G_i$ and
similarly $L,B$. Independence gives $(G_j+B_j)\cap(G+B)=0$, and $L\subseteq G+B$. By (b) the module whose length is
the left-hand side for $F'$ is $T_1+T_2$ modulo $S_1+S_2$ with $T_1=G_j\cap(L_j+B_j)$, $S_1=L_j\cap G_j$,
$T_2=G\cap(L+B)$, $S_2=L\cap G$; we have $S_k\subseteq T_k$ and $T_1\cap T_2\subseteq(G_j+B_j)\cap(G+B)=0$. By (c) and the
induction hypothesis the left-hand side for $F'$ is $\mathrm{cl}(L_j;G_j,B_j)+\sum_{i\in F}\mathrm{cl}(L_i;G_i,B_i)$.
\end{proof}

\begin{proposition}[Intersections with subspaces under base change]\label{pf:baselat2:latbc}
Let $K$ be a number field, $\OO$ a discrete valuation ring with fraction field $F$, and $e:K\to F$ a ring
homomorphism with $e(\OO_K)\subseteq\OO$. Let $V$ be a finite-dimensional $K$-vector space, $V'$ an $F$-vector
space, $\iota:V\to V'$ an $e$-semilinear map, and $(b_a)$ a $K$-basis of $V$ such that $(\iota(b_a))$ is
$F$-linearly independent. Then for every $\OO_K$-submodule $A\subseteq V$ and every $K$-subspace $U\subseteq V$
\[
  \mathrm{span}_{\OO}\,\iota(A\cap U)=\mathrm{span}_{\OO}\,\iota(A)\ \cap\ \mathrm{span}_F\,\iota(U).
\]
\formallabel{GlobalTheory.T.lattice_inf_subspace_baseChange_dvr}
\formallabel{GlobalTheory.LatticeInterBaseChange.span_inf_baseChange}
\end{proposition}

\begin{remark}
The proof uses only that $\OO$ is an integral domain with fraction field $F$; the discrete valuation
hypothesis is not used.
\end{remark}

\begin{proof}
The inclusion ``$\subseteq$'' is clear. For the converse we need two auxiliary facts.

(i) \emph{Every $K$-linear form $\varphi$ on $V$ has an $F$-linear extension $\psi$ on $V'$ with
$\psi\circ\iota=e\circ\varphi$.} Define $\psi$ on the span of the $\iota(b_a)$ (a basis of it) by
$\psi(\iota(b_a))=e(\varphi(b_a))$ and extend to $V'$; the two $e$-semilinear maps $\psi\circ\iota$ and $e\circ\varphi$
agree on the basis $(b_a)$.

(ii) \emph{For $c_1,\dots,c_n\in K$ there are $\beta_1,\dots,\beta_n\in K$ with $\beta_jc_k\in\OO_K$ for all $j,k$ and
$c_k\sum_j\beta_jc_j=c_k$ for all $k$.} If all $c_k$ vanish take $\beta=0$. Otherwise the fractional ideal
$I=\sum_k\OO_Kc_k$ is invertible, $1\in II^{-1}$, so $1=\sum_j\beta_jc_j$ with $\beta_j\in I^{-1}$, and $\beta_jc_k\in I^{-1}I=\OO_K$.

(iii) \emph{Key step.} Let $R$ be a commutative ring with a homomorphism to $F$ such that $e(\OO_K)$ lies in the
image of $R$, let $\varphi,\psi$ be as in (i), and let $y\in\mathrm{span}_R\,\iota(A)$ with $\psi(y)=0$. Then
$y\in\mathrm{span}_R\,\iota(A\cap\ker\varphi)$. Indeed, write $y=\sum_kr_k\iota(a_k)$ with $a_k\in A$, put $c_k=\varphi(a_k)$, and
choose $\beta_j$ by (ii). Then $\rho_{jk}=\beta_jc_k\in\OO_K$ and $w_k=a_k-\sum_j\rho_{jk}a_j\in A$ satisfy
$\varphi(w_k)=c_k-c_k\sum_j\beta_jc_j=0$. Moreover
\[
  \sum_kr_k\iota(w_k)=y-\sum_j e(\beta_j)\Bigl(\sum_kr_k\,e(c_k)\Bigr)\iota(a_j)=y,
\]
because $\sum_kr_ke(c_k)=\sum_kr_k\psi(\iota(a_k))=\psi(y)=0$. Hence $y\in\mathrm{span}_R\,\iota(A\cap\ker\varphi)$.

Now choose linear forms $f_1,\dots,f_r$ on $V$ with $U=\bigcap_j\ker f_j$ (coordinates on $V/U$) and extensions
$\psi_j$ by (i). Let $x\in\mathrm{span}_\OO\,\iota(A)\cap\mathrm{span}_F\,\iota(U)$; then $\psi_j(x)=0$ for all $j$, as
$\psi_j$ vanishes on $\iota(U)$. Applying (iii) (with $R=\OO$) successively to $f_1,\dots,f_r$ and the submodules
$A\cap\ker f_1\cap\dots\cap\ker f_{t}$ gives $x\in\mathrm{span}_\OO\,\iota(A\cap U)$.
\end{proof}

\begin{lemma}[Finite length of torsion modules]\label{pf:baselat2:finlen}
Let $R$ be a Noetherian ring of Krull dimension at most one and $X$ a finitely generated torsion $R$-module.
Then $X$ has finite length.
\formallabel{GlobalTheory.isFiniteLength_of_isTorsion}
\end{lemma}

\begin{proof}
Since $X$ is finitely generated and torsion, its annihilator contains a non-zero-divisor $r$ (a product of
non-zero-divisors killing the generators). $X$ is a finitely generated module over $R/rR$. The ring $R/rR$ is
Noetherian of dimension $0$: a prime containing $r$ is not minimal, because in a Noetherian ring the minimal primes
consist of zero-divisors, so it has height $\ge1$ and is maximal as $\dim R\le1$. Hence $R/rR$ is Artinian, and a
finitely generated module over an Artinian ring is Noetherian and Artinian, i.e.\ of finite length
\cite[Ch.~4, Ch.~6 and Thm.~8.5]{AtiyahMacdonald1969}.
\end{proof}

We now prove the equality of local and rational congruence lengths. Let $D$ be a global datum and $M$ a
module family. For a tuple $\mu$ and $1\le j\le m$ let $M_L(\mu,j)$ be the $L$-module of the development, with the
injective semilinear map $j_{\mathrm{Mt}}$ into the complex module $M(\mu,j)$ and the $\OO_L$-lattice $\mathrm{lat}$ spanned by
the elements mapped to the generators of the module lattice. Let $Z_0\subseteq M(\mu,j)$ be the distinguished
subspace of the development, $U=j_{\mathrm{Mt}}^{-1}(Z_0)$ and $U'=j_{\mathrm{Mt}}^{-1}(Z_0^{\perp})$ ($L$-subspaces), and
\[
  \mathrm{up}=U\cap(\mathrm{lat}+U'),\qquad \mathrm{low}=\mathrm{lat}\cap U .
\]
The \emph{rational congruence length} $\eta_{\mathrm{rat}}(\mu,j,w)$ at a place $w$ of $L$ is the normalised length of
$\mathrm{up}/\mathrm{low}$ at $w$. For a local index $i$ (with tuple $\mu_i$, index $j_i$, coefficient ring $\OO_i$, integral
local module $M_{\mathrm{in}}(i)$ and projection $\mathrm{pr}_i$ of $M_L(\mu_i,j_i)$ to the local module) the local module
decomposes into the sum $E^+(i)$ of the eigenspaces at the good points of the theta locus and the sum $E^-(i)$ of
the others; the \emph{local congruence length} $\eta_{\mathrm{loc}}(i)$ is defined in the development, and
$\eta_J(j,w)=\sum_{i\in\mathrm{locs}(j,w)}\eta_{\mathrm{loc}}(i)$, the sum over the local indices at $(j,w)$. A \emph{normalised local
compatibility datum} $T$ for $(D,M)$ provides, for each $(j,w)$, a discrete valuation ring $\OO=\OO_c(j,w)$ with
fraction field $K$, a $K$-vector space $V$ with an $e_c$-semilinear map $\iota:M_L(\mu_w,j)\to V$ mapping bases to bases,
a Hecke algebra acting on $V$, and for $i\in\mathrm{locs}(j,w)$ distinct maximal ideals $\mathfrak m_i$ of the Hecke algebra
and injective maps $\Psi_i$ from the local modules to $V$ whose ranges are the parts of $V$ at $\mathfrak m_i$; the
\emph{local lattice} at $(j,w)$ is the $\OO$-span $\Lambda_C$ of the images under $\iota$ of the generators of $\mathrm{lat}$.

\begin{theorem}[Local and rational congruence lengths]\label{pf:baselat2:etasum}
Let $T$ be a normalised local compatibility datum for $(D,M)$, $1\le j\le m$ and $w$ a place of $L$ with
$\mu_w\neq\mu_0$. Assume the standard facts for the Hecke operators of $T$, that the local lattice at $(j,w)$ is
stable under the Hecke operators at the generating places, that $U+U'=M_L(\mu_w,j)$, and that for every
$i\in\mathrm{locs}(j,w)$ the projection $\mathrm{pr}_i$ maps $U$ into $E^+(i)$ and $U'$ into $E^-(i)$. Then
\[
  \eta_J(j,w)=\eta_{\mathrm{rat}}(\mu_w,j,w).
\]
\formallabel{GlobalTheory.T.sum_localCongruenceLength_eq_rationalLength}
\formallabel{GlobalTheory.EtaJoinSum.sum_etaLoc_eq_etaRat}
\end{theorem}

\begin{proof}
Write $F=\mathrm{locs}(j,w)$, let $\ell$ be the residue characteristic of $w$, and put, for $i\in F$,
\begin{gather*}
  G_i=\Psi_i\bigl(E^+(i)\bigr),\quad B_i=\Psi_i\bigl(E^-(i)\bigr),\quad L_i=\Psi_i\bigl(M_{\mathrm{in}}(i)\bigr),\\
  G'=\mathrm{span}_K\iota(U),\quad B'=\mathrm{span}_K\iota(U').
\end{gather*}

\emph{Step 1: the local terms.} By \lref{GlobalTheory.EtaJoinSum.etaLoc_eq_loc},
$\eta_{\mathrm{loc}}(i)=\mathrm{cl}(L_i;G_i,B_i)/\len_\OO(\OO/\ell\OO)$ in the notation of \cref{pf:baselat2:cl}, provided the
congruence length is finite. It is: by \lref{GlobalTheory.EtaJoinSum.length_loc} it equals the congruence
length of $M_{\mathrm{in}}(i)$ with respect to $E^+(i),E^-(i)$ in the local module, which is finite by
\lref{GlobalTheory.T.finiteLength_lattice_projection_quotient}, since $M_{\mathrm{in}}(i)$ is finitely generated and
spans the local module over the fraction field, $E^+(i)\cap E^-(i)=0$
(\lref{GlobalTheory.T.heckeEigenspaceSum_disjoint_compl}) and $E^+(i)+E^-(i)$ is the whole local module. For the
last point: the local module is spanned by the projections $\mathrm{pr}_i(y)$, $y\in M_L(\mu_i,j_i)$, and writing
$y=u+u'$ with $u\in U$, $u'\in U'$ (hypothesis $U+U'=M_L$) we get $\mathrm{pr}_i(y)=\mathrm{pr}_i(u)+\mathrm{pr}_i(u')\in E^+(i)+E^-(i)$ by the
hypothesis on $\mathrm{pr}_i$.

\emph{Step 2: the global term.} $U\cap U'=0$, since $j_{\mathrm{Mt}}$ is injective (it maps a basis to a linearly
independent family) and $Z_0\cap Z_0^\perp=0$; with the hypothesis, $M_L(\mu_w,j)=U\oplus U'$. By
\lref{GlobalTheory.T.coeffEmbedding_integers_resChar_unit}, $e_c(\OO_L)\subseteq\OO$; hence the $\OO$-span of
$\iota(\mathrm{lat})$ equals the $\OO$-span of the images of the generators of $\mathrm{lat}$, i.e.\ $\mathrm{span}_\OO\iota(\mathrm{lat})=\Lambda_C$.
\Lref{GlobalTheory.T.latticeProjection_baseChange_dvr_comm} gives
$\mathrm{span}_\OO\iota(\mathrm{up})=G'\cap(\Lambda_C+B')$, and \cref{pf:baselat2:latbc} gives $\mathrm{span}_\OO\iota(\mathrm{low})=\Lambda_C\cap G'$.
The lattice $\mathrm{lat}$ is finitely generated (its generating set is the preimage of a finite set under the
injective map $j_{\mathrm{Mt}}$), hence so is $\mathrm{up}$, which is the image of $\mathrm{lat}$ under the projection onto $U$ along
$U'$; and $\mathrm{low}\subseteq\mathrm{up}$. By \lref{GlobalTheory.T.normalizedLength_latticeQuotient_baseChange_eq} the
normalised length of $\mathrm{up}/\mathrm{low}$ at $w$ is
$\len_\OO\bigl(\mathrm{span}_\OO\iota(\mathrm{up})/\mathrm{span}_\OO\iota(\mathrm{low})\bigr)/\len_\OO(\OO/\ell\OO)$; hence
\[
  \eta_{\mathrm{rat}}(\mu_w,j,w)=\mathrm{cl}(\Lambda_C;G',B')/\len_\OO(\OO/\ell\OO).
\]

\emph{Step 3: decomposition.} (a) The family $(G_i+B_i)_{i\in F}$ is independent: $G_i$ and $B_i$ lie in the part of
$V$ at $\mathfrak m_i$, the $\mathfrak m_i$ are distinct, and the part at a maximal ideal meets the sum of the parts at the other
maximal ideals trivially, because these are sums of common eigenspaces of the commuting Hecke operators (by the
standard facts) over disjoint sets of eigensystems, and such eigenspaces for distinct eigensystems are
independent (\lref{Formalism.L2.NCov.iSupIndep_alt}).
(b) $G_i\cap B_i=0$, since $\Psi_i$ is injective and $E^+(i)\cap E^-(i)=0$. (c) $L_i\subseteq G_i+B_i$, since
$E^+(i)+E^-(i)$ is the whole local module. (d) $G'\subseteq\sum_iG_i$ and $B'\subseteq\sum_iB_i$: every $y\in M_L(\mu_w,j)$ satisfies
$\iota(y)=\sum_{i\in F}\Psi_i(\mathrm{pr}_iy)$ (\lref{GlobalTheory.T.rationalForm_eq_sum_localisation_projections}),
and $\mathrm{pr}_iy\in E^+(i)$ for $y\in U$, $\mathrm{pr}_iy\in E^-(i)$ for $y\in U'$. (e) $G'+B'=V$, since $U+U'=M_L(\mu_w,j)$ and
$\iota$ maps a basis to a basis. (f) $\bigl(\sum_iG_i\bigr)\cap\bigl(\sum_iB_i\bigr)=0$: if $\sum_ig_i=\sum_ib_i$ with $g_i\in G_i$,
$b_i\in B_i$, then $\sum_i(g_i-b_i)=0$ with $g_i-b_i\in G_i+B_i$, so $g_i=b_i$ for all $i$ by (a), and $g_i\in G_i\cap B_i=0$ by
(b). Now use the following fact about submodules:
\emph{if $a\subseteq c$, $d\subseteq e$, $a+d=V$ and $c\cap e=0$, then $c=a$}; indeed, by the modular law
$c=(a+d)\cap c=a+(d\cap c)$, and $d\cap c\subseteq e\cap c=0$. Applied to $(a,c,d,e)=(G',\sum G_i,B',\sum B_i)$ and to
$(B',\sum B_i,G',\sum G_i)$ it gives $\sum_iG_i=G'$ and $\sum_iB_i=B'$, by (d), (e), (f). Finally (g)
$\Lambda_C=\sum_{i\in F}L_i$, by \lref{GlobalTheory.T.integralFormLattice_eq_sum_localised_lattices} (which uses the
stability of the local lattice under the Hecke operators).

\emph{Step 4.} The same identities hold for the $\OO$-modules obtained by restriction of scalars. By (a), (c),
(g) and \cref{pf:baselat2:cl},
\[
  \mathrm{cl}(\Lambda_C;G',B')=\mathrm{cl}\Bigl(\sum_iL_i;\sum_iG_i,\sum_iB_i\Bigr)=\sum_{i\in F}\mathrm{cl}(L_i;G_i,B_i).
\]
All terms are finite by Step 1; dividing by $\len_\OO(\OO/\ell\OO)$ and using Steps 1 and 2 gives
$\eta_{\mathrm{rat}}(\mu_w,j,w)=\sum_{i\in F}\eta_{\mathrm{loc}}(i)=\eta_J(j,w)$.
\end{proof}

% ---------------------------------------------------------------------------
\subsection{Auxiliary lemmas on local fields, Galois groups and characters}\label{sec:pf:baselat2:aux}

The lemmas of this subsection are used in the stage constructions and in the section of
\lref{GlobalTheory.Mem.Top.Bl.exists_baseLattice_of_pair_alt}.

\begin{lemma}\label{pf:baselat2:resunit}
\textup{(a)} Let $A$ be a local ring, and $q,\ell$ coprime natural numbers with $\ell\in\mathfrak m_A$. Then $q$ is a unit of
$A$.
\textup{(b)} Let $E$ be a number field, $\mathfrak P$ a prime of $\OO_E$ and $\ell$ a prime number with $\ell\notin\mathfrak P$. Then
$|\OO_E/\mathfrak P|$ is prime to $\ell$.
\textup{(c)} Let $\OO\to A$ be a homomorphism of local rings making $A$ a finite $\OO$-module, and $\ell\in\N$ with
$\ell\in\mathfrak m_{\OO}$. Then $\ell\in\mathfrak m_A$.
\formallabel{GlobalTheory.ResidueUnit.isUnit_of_coprime}
\formallabel{GlobalTheory.ResidueUnit.card_residueField_coprime}
\formallabel{GlobalTheory.ResidueUnit.natCast_mem_maximalIdeal_of_finite}
\end{lemma}

\begin{proof}
(a) Write $aq+b\ell=1$ with $a,b\in\Z$. If $q\in\mathfrak m_A$ then $1\in\mathfrak m_A$. (b) $|\OO_E/\mathfrak P|=p^n$ where
$p$ is the characteristic of the residue field; if $p=\ell$, then $\ell=0$ in $\OO_E/\mathfrak P$, i.e.\ $\ell\in\mathfrak P$.
(c) $A$ is integral over $\OO$, so $\mathfrak m_A\cap\OO$ is a maximal ideal, hence equal to $\mathfrak m_\OO$, which
contains $\ell$.
\end{proof}

In the next two lemmas $E$ is a number field, $\Omega$ an algebraic closure of $E$, $G_E=\Gal(\Omega/E)$ with
its Krull topology, and for a valuation subring $A$ of $\Omega$ the decomposition group is
$D_A=\{\sigma:\ \sigma(A)=A\}$.

\begin{lemma}\label{pf:baselat2:powapply}
For $\sigma\in G_E$ and $y\in\Omega$ there is $d\ge1$ with $\sigma^d(y)=y$.
\formallabel{GlobalTheory.PrimeToInertia.exists_pow_apply_eq}
\end{lemma}

\begin{proof}
The stabiliser of $y$ is an open subgroup, hence of finite index; some positive power of $\sigma$ lies in it.
\end{proof}

\begin{lemma}\label{pf:baselat2:valapply}
If $\sigma\in D_A$, then $v_A(\sigma y)=v_A(y)$ for all $y\in\Omega$.
\formallabel{GlobalTheory.PrimeToInertia.valuation_apply_eq}
\end{lemma}

\begin{proof}
For $y=0$ there is nothing to prove. Elements of $D_A$ preserve the maximal ideal of $A$, hence the condition
$v_A(z)<1$; applied to $z=a/b$ this shows that every $\tau\in D_A$ preserves strict inequalities
$v_A(a)<v_A(b)$ ($b\ne0$). Let $d\ge1$ with $\sigma^d y=y$ (\cref{pf:baselat2:powapply}). If $v_A(\sigma y)<v_A(y)$, then
applying $\sigma^i\in D_A$ shows that $i\mapsto v_A(\sigma^iy)$ is strictly decreasing, contradicting $\sigma^dy=y$; the
case $v_A(\sigma y)>v_A(y)$ is symmetric.
\end{proof}

\begin{lemma}\label{pf:baselat2:embcompl}
Let $E$ be a number field, $v$ a prime of $\OO_E$, and $C$ a field with a non-archimedean absolute value.
\textup{(a)} Let $a,b:E\to C$ be embeddings such that for $x\in\OO_E$: $|a(x)|<1\iff x\in v\iff|b(x)|<1$. Then
$|a(x)|<1\iff|b(x)|<1$ for all $x\in E$.
\textup{(b)} Let $C=\Qlbar$ and let $a,b:E\to\Qlbar$ be embeddings with $|a(x)|<1\iff|b(x)|<1$ for all $x\in E$. Then
every $\Q_\ell$-linear relation $\sum_kc_ka(y_k)=0$ ($c_k\in\Q_\ell$, $y_k\in E$) implies $\sum_kc_kb(y_k)=0$, and there is
$g\in\Aut(\Qlbar/\Q_\ell)$ with $g\circ a=b$.
\formallabel{GlobalTheory.EmbThroughCompletion.norm_lt_one_iff}
\formallabel{GlobalTheory.EmbThroughCompletion.sum_eq_zero}
\formallabel{GlobalTheory.EmbThroughCompletion.exists_conj}
\end{lemma}

\begin{proof}
(a) Integers have absolute value $\le1$ under any embedding into $C$ (they are integral over $\Z$ and $C$ is
non-archimedean). Write $x=n/d$ or $x=d/n$ with $n\in\OO_E$ and $d\in\OO_E\smallsetminus v$ (possible since
$\OO_{E,v}$ is a discrete valuation ring). Then $|a(d)|=|b(d)|=1$. In the first case
$|a(x)|<1\iff n\in v\iff|b(x)|<1$. In the second case $|a(x)|\cdot|a(n)|=1$ with $|a(n)|\le1$, so $|a(x)|\ge1$, and
likewise for $b$.

(b) For $c\in\Q_\ell^n$ consider $s_a(c)=\sum_kc_ka(y_k)$ and $s_b(c)$. The sets $\{c:|s_a(c)|<1\}$ and $\{c:|s_b(c)|<1\}$
are clopen (open balls are closed in an ultrametric space), so $\{c:\ |s_a(c)|<1\Rightarrow|s_b(c)|<1\}$ is closed; it
contains the dense subset $\Q^n$ because for $c\in\Q^n$ we have $s_a(c)=a(\sum_kc_ky_k)$ and similarly for $b$. Hence
$|s_a(c)|<1$ implies $|s_b(c)|<1$ for all $c$. If $s_a(c)=0$ but $s_b(c)\ne0$, choose $r\in\Q_\ell$ with
$|r|>|s_b(c)|^{-1}$; then $s_a(rc)=0$ and $|s_b(rc)|>1$, a contradiction. For the last assertion let $\gamma$ be a
primitive element of $E/\Q$. The minimal polynomial of $a(\gamma)$ over $\Q_\ell$ is a $\Q_\ell$-linear relation among
the powers $a(\gamma^k)$, so it annihilates $b(\gamma)$; hence there is $g\in\Aut(\Qlbar/\Q_\ell)$ with $g(a(\gamma))=b(\gamma)$,
and $g\circ a=b$ since both are ring homomorphisms $E\to\Qlbar$ agreeing on $\gamma$.
\end{proof}

\begin{lemma}\label{pf:baselat2:antiladic}
Let $E/F$ be a quadratic extension of number fields with non-trivial automorphism $c$, such that for every
prime $\mathfrak P$ of $E$ and $k\in F^\times$ the normalised valuations satisfy
$v_{\mathfrak P}(k)=v_{\mathfrak P\cap F}(k)$. Let $\chi$ be a Hecke character of $E$ trivial on the id\`eles of $F$, and let
$v,v'$ be primes of $E$ with $v'=c(v)$ at which $\chi$ is unramified. Then $\chi(\varpi_v)\chi(\varpi_{v'})=1$, where
$\chi(\varpi_w)$ is the value of $\chi$ at a uniformiser of $E_w$.
\formallabel{GlobalTheory.AntiLAdic.val_mul_val_conj}
\end{lemma}

\begin{proof}
Let $u=v\cap\OO_F$. The primes of $E$ above $u$ are exactly $v$ and $v'$, since $\Gal(E/F)=\{1,c\}$ acts transitively on
them. Let $\varpi_u$ be a uniformiser of $F_u$. Because $\chi$ is trivial on $F_u^\times$ (embedded diagonally in
$\prod_{w\mid u}E_w^\times$), $\prod_{w\mid u}\chi_w(\varpi_u)=1$. By the hypothesis on valuations, $\varpi_u$ is a
uniformiser of each $E_w$, $w\mid u$, and as $\chi$ is unramified at $w$, $\chi_w(\varpi_u)=\chi(\varpi_w)$. If $v\ne v'$ this
is the assertion. If $v=v'$, the product has one factor, so $\chi(\varpi_v)=1$ and the assertion holds as well.
\end{proof}

\begin{lemma}\label{pf:baselat2:valone}
Let $K$ be a number field and $v$ a prime of $\OO_K$.
\textup{(a)} For $y\in\OO_K\smallsetminus v$ the image of $y$ in $K_v$ has valuation $1$.
\textup{(b)} For $x\in K$, $|x|_v\le1$ (the absolute value of the finite place $v$) if and only if $v(x)\le1$.
\formallabel{Neutral.B1.valued_algebraMap_eq_one}
\formallabel{GlobalTheory.EmbeddingPlace.finitePlace_le_one_iff}
\end{lemma}

\begin{proof}
(a) The valuation of $K_v$ extends that of $K$, and $v(y)=1$ for $y\notin v$. (b) $|x|_v=\Nm(v)^{-\ord_v(x)}$ is a
monotone function of $v(x)$ with $|x|_v=1$ exactly when $v(x)=1$, because $\Nm(v)>1$.
\end{proof}

\begin{lemma}\label{pf:baselat2:locone}
Let $L$ be a number field, $\psi$ a continuous character of the id\`eles of $L$, and $S$ a finite set of primes of
$L$ such that $\prod_{Q\in S}\psi_Q(k)=1$ for all $k\in L^\times$ (where $\psi_Q$ is the restriction of $\psi$ to
$L_Q^\times$). Then $\psi_v=1$ on $L_v^\times$ for every $v\in S$.
\formallabel{GlobalTheory.CuspRep.NormTriv.loc_eq_one_of_prod}
\end{lemma}

\begin{proof}
The diagonal image of $L^\times$ is dense in $\prod_{Q\in S}L_Q^\times$ (weak approximation;
\lref{GlobalTheory.FrobDense.denseRange_units}). The function $y\mapsto\prod_{Q\in S}\psi_Q(y_Q)$ on this product is
continuous and equal to $1$ on a dense subset, hence identically $1$. Evaluating at the element which is $x$ at
$v$ and $1$ elsewhere gives $\psi_v(x)=1$.
\end{proof}

\begin{lemma}\label{pf:baselat2:jointdiag}
Let $f$ be an endomorphism of a vector space $M$ over a field $K$.
\textup{(a)} For a finite set $s\subseteq K$, $\ker\prod_{r\in s}(f-r)=\sum_{r\in s}\ker(f-r)$.
\textup{(b)} If $W\subseteq M$ is $f$-stable and $q\in K[X]$, then $q(f|_W)(y)=q(f)(y)$ for $y\in W$.
\formallabel{GlobalTheory.JointDiag.ker_prod}
\formallabel{GlobalTheory.JointDiag.aeval_restrict_coe}
\end{lemma}

\begin{proof}
(a) Induction on $s$; for $s=s_0\cup\{a\}$ the polynomials $X-a$ and $\prod_{r\in s_0}(X-r)$ are coprime, and the kernel
of $p_1(f)p_2(f)$ for coprime $p_1,p_2$ is $\ker p_1(f)+\ker p_2(f)$. (b) Both sides are $K$-linear in $q$ and
multiplicative, and agree for $q=X$ and constants.
\end{proof}

For a group $G$, an injective homomorphism $\iota:\Gamma\to G$ and a subgroup $J$, a family $(y_i)_{i\in I}$ is a family
of \emph{representatives} of $\Gamma\backslash G/J$ if every $h\in G$ can be written $h=\iota(\gamma)y_ik$ ($\gamma\in\Gamma$,
$k\in J$) for exactly one $i$; and $\mathrm{Stab}_J(y)=\{\gamma\in\Gamma:\ y^{-1}\iota(\gamma)y\in J\}$.

\begin{lemma}\label{pf:baselat2:classsum}
\textup{(a)} If $(y_i)$ are representatives of $\Gamma\backslash G/J$ and $z$ is central in $G$, then $(y_iz)$ are
representatives as well.
\textup{(b)} A surjective homomorphism maps the centre into the centre.
\textup{(c)} Let $J_0\subseteq J$ be subgroups of $G$, $\iota$ injective, $F\subseteq J$ a finite set such that every
$\kappa\in J$ has exactly one $u\in F$ with $u^{-1}\kappa\in J_0$, and let $(y^0_i)_{i\in I_0}$, $(y_i)_{i\in I}$ be finite
families of representatives of $\Gamma\backslash G/J_0$ and $\Gamma\backslash G/J$. Then for every $\Phi:G\to\C$ which is left
$\iota(\Gamma)$-invariant and right $J_0$-invariant,
\[
  \sum_{i\in I_0}\frac{\Phi(y^0_i)}{|\mathrm{Stab}_{J_0}(y^0_i)|}=\sum_{i\in I}\frac{1}{|\mathrm{Stab}_J(y_i)|}\sum_{u\in F}\Phi(y_iu),
\]
where $1/|H|$ is read as $0$ when $H$ is infinite.
\formallabel{GlobalTheory.LiftClassSumPrinc.isReps_mul_central}
\formallabel{GlobalTheory.LiftClassSumPrinc.map_mem_center}
\formallabel{GlobalTheory.T.classSum_regroup_finerLevel_cosets}
\formallabel{GlobalTheory.ClassRefineBridge.sum_isReps_refine}
\end{lemma}

\begin{proof}
(a) Given $h$, write $hz^{-1}=\iota(\gamma)y_ik$; then $h=\iota(\gamma)(y_iz)k$ because $z$ commutes with $k$, and uniqueness
follows in the same way from the uniqueness for $hz^{-1}$. (b) If $f(g)$ is arbitrary, $f(z)f(g)=f(zg)=f(gz)=f(g)f(z)$.
(c) Let $Q=\iota(\Gamma)$. Since $\iota$ is injective, $\iota$ maps $\mathrm{Stab}_J(y)$ isomorphically onto the stabiliser
$\{q\in Q:\ y^{-1}qy\in J\}$, so the two have the same cardinality; the representative families for $\Gamma$ are
representative families for the subgroup $Q$; and $i\mapsto y_i$ is injective (if $y_i=y_{i'}$, then
$y_i=\iota(1)y_{i'}1$, and uniqueness gives $i=i'$), so the sums over $I_0$ and $I$ equal the sums over the images
$\{y^0_i\}$ and $\{y_i\}$. For the subgroup $Q$ and finite sets of representatives, the identity is
\lref{GlobalTheory.ClassRefine.sum_classReps_refine}, which regroups the sum over the finer double cosets
$Q\backslash G/J_0$ according to the coarser ones $Q\backslash G/J$.
\end{proof}

\paragraph{The gap hypothesis.} For a global datum $D$ with ranks $(k,m)$ and $r=k+m-1$ and with Harish-Chandra
parameters $b_j(\tau)$ ($j<m$, $\tau$ a CM type index), the \emph{gap hypothesis} $\mathrm{Gap}(D)$ is the condition
\[
  1801\le r\quad\text{and}\quad b_{j+1}(\tau)-b_j(\tau)\ \ge\ 16\,r^{19}\quad\text{for all }\tau\text{ and }j+1<m .
\]

\begin{lemma}[Gap hypothesis]\label{pf:baselat2:gap}
Eventually in the scale (in the sense of $\mathcal{E}_0$, \cref{sec:pf:baselat2:conv}), $\mathrm{Gap}(D_X)$ holds
for all stages $X$ over the pre-base of $K$.
\formallabel{GlobalTheory.Mem.Top.NodeGap.ev_gapHyp}
\end{lemma}

\begin{proof}
The ranks of a scale context $K$ with scale $(x,\eps,J)$ satisfy $\sqrt{x^{1-\eps/8}}\le m$, $k\le2x^{9/2+10\eps}$ and $m+3\le k$,
and the Harish-Chandra parameters of its words are arithmetic progressions $b_j(\tau)=b_1(\tau)+jh_0$, $e_j=b_j+h_0/2$,
with common difference $h_0=\lceil x^{c_h}\rceil$ (this is part of the definition of a scale context). Hence, for $\eps<1/52$:
if $x\ge M^4$ then $M\le m$ (as $M^2\le x^{1/2}\le x^{1-\eps/8}$); and if $x\ge4$ then $r\le2k\le4x^5\le x^6$. Let $A=16$,
$N=19$, $M=1801$. Given $\eps$ and the other constants, take $c_h\ge6N+1$ and $x\ge\max(4,A,M^4)$. Then $M\le m$ and
\[
  A\,r^N\ \le\ x\cdot x^{6N}=x^{6N+1}\ \le\ x^{c_h}\ \le\ h_0 .
\]
For stages $X$ over the pre-base of $K$ the ranks and parameters are those of $K$. Finally
$r=k+m-1\ge m\ge1801$, and $b_{j+1}-b_j=h_0\ge16r^{19}$; monotonicity of
$\mathcal{E}_0$ concludes.
\end{proof}

\paragraph{Classical results formalized in the region.} The following statements are proved in the
development and are classical; we state them and refer to the literature.

\begin{proposition}[Density of rational points of definite unitary groups]\label{pf:baselat2:zariski}
Let $L$ be a CM field with complex conjugation $c$, $n\in\N$, $K$ a field of characteristic $0$, and $P$ a polynomial
over $K$ in variables $X_{\sigma,a,b}$ ($\sigma:L\to K$, $1\le a,b\le n$). If $P\bigl((\sigma(\gamma_{ab}))\bigr)=0$ for every
$\gamma\in\Uni_n(L)=\{\gamma\in\GL_n(L):\ \gamma\,{}^t\bar\gamma=1\}$, then $P\bigl((B(\sigma)_{ab})\bigr)=0$ for every family
$B=(B(\sigma))_{\sigma}$ in $\GL_n(K)$ with $B(\sigma\circ c)={}^tB(\sigma)^{-1}$ for all $\sigma$.
\formallabel{GlobalTheory.T.zariskiDense_rationalPoints_definiteUnitary_charZero}
\end{proposition}

\begin{proof}
The families $B$ are the $K$-points of the image of $G_K$ under the projection to the factors of
$L\otimes_\Q K$ indexed by embeddings $L\to K$, where $G=\mathrm{Res}_{L^+/\Q}\,\Uni_n$ is the connected reductive group with
$G(\Q)=\Uni_n(L)$. Since $G(\Q)$ is Zariski dense in $G$ \cite[Cor.~18.3]{Borel1991}, it is dense in $G_K$, and the
assertion follows. The formal proof uses the Cayley transform: it chooses a norm-one scalar $u$ such that
$\det(1+u\,B(\sigma))\ne0$ for all $\sigma$,
translates $P$ by the corresponding scalar unitary matrix, and applies the density of Cayley images of rational
skew-Hermitian matrices.
\end{proof}

The irreducibility of tensor products of irreducible algebraic representations of $\GL_n$ attached to distinct
embeddings of a CM field (\lref{GlobalTheory.T.weightTensorRepOfEmbeddings_isIrreducible}; compare
\cite[Lectures~2 and~6]{FultonHarris1991}), which uses \cref{pf:baselat2:zariski}, and the fact that the inertia
groups of two distinct places of $\Omega$ above the same prime of $E$ meet trivially
(\lref{GlobalTheory.T.inertia_two_places_over_same_prime_trivial_meet}; for decomposition groups see
\cite[Ch.~XII, \S1]{NeukirchSchmidtWingberg2008}) are further classical statements of this kind; their formal proofs
are not converted here.

% === END ===
